% =====================================================================
%  utpnotes.tex
%  Shared preamble for the UTP teaching-assistant course notes.
%
%  Every book includes this file, and it is the ONLY file a book needs
%  from .shared/: the cover, the logos and the optional diagram
%  libraries are all reached from here. It owns everything the books
%  have in common, so a change here reaches every course at once.
%
%  ------------------------------------------------------------------
%  HOW A COURSE USES THIS FILE
%  ------------------------------------------------------------------
%      \documentclass[11pt,oneside,a4paper]{report}
%
%      \def\utprunninghead{Machine Learning Notes}   % section 1
%      \def\utplabtotal{8}
%      \def\utpcoursetitle{Machine Learning}
%      \def\utpcourseoutline{One or two lines for the cover.}
%
%      % =====================================================================
%  utpnotes.tex
%  Shared preamble for the UTP teaching-assistant course notes.
%
%  Every book includes this file, and it is the ONLY file a book needs
%  from .shared/: the cover, the logos and the optional diagram
%  libraries are all reached from here. It owns everything the books
%  have in common, so a change here reaches every course at once.
%
%  ------------------------------------------------------------------
%  HOW A COURSE USES THIS FILE
%  ------------------------------------------------------------------
%      \documentclass[11pt,oneside,a4paper]{report}
%
%      \def\utprunninghead{Machine Learning Notes}   % section 1
%      \def\utplabtotal{8}
%      \def\utpcoursetitle{Machine Learning}
%      \def\utpcourseoutline{One or two lines for the cover.}
%
%      % =====================================================================
%  utpnotes.tex
%  Shared preamble for the UTP teaching-assistant course notes.
%
%  Every book includes this file, and it is the ONLY file a book needs
%  from .shared/: the cover, the logos and the optional diagram
%  libraries are all reached from here. It owns everything the books
%  have in common, so a change here reaches every course at once.
%
%  ------------------------------------------------------------------
%  HOW A COURSE USES THIS FILE
%  ------------------------------------------------------------------
%      \documentclass[11pt,oneside,a4paper]{report}
%
%      \def\utprunninghead{Machine Learning Notes}   % section 1
%      \def\utplabtotal{8}
%      \def\utpcoursetitle{Machine Learning}
%      \def\utpcourseoutline{One or two lines for the cover.}
%
%      % =====================================================================
%  utpnotes.tex
%  Shared preamble for the UTP teaching-assistant course notes.
%
%  Every book includes this file, and it is the ONLY file a book needs
%  from .shared/: the cover, the logos and the optional diagram
%  libraries are all reached from here. It owns everything the books
%  have in common, so a change here reaches every course at once.
%
%  ------------------------------------------------------------------
%  HOW A COURSE USES THIS FILE
%  ------------------------------------------------------------------
%      \documentclass[11pt,oneside,a4paper]{report}
%
%      \def\utprunninghead{Machine Learning Notes}   % section 1
%      \def\utplabtotal{8}
%      \def\utpcoursetitle{Machine Learning}
%      \def\utpcourseoutline{One or two lines for the cover.}
%
%      \input{../.shared/utpnotes}
%
%      \utpcontentspage{\small}{\normalsize}{0.65em}{0.07em}{0.16em}
%
%      \begin{document}
%      \utpfrontmatter                     % cover, contents, map
%      \part{Orientation}
%      \input{chapters/01-where-to-begin}
%      \end{document}
%
%  Configuration goes BEFORE the \input; overrides go AFTER it.
%  Paths are relative, so compile from inside the course folder.
%
%  ------------------------------------------------------------------
%  NAMING
%  ------------------------------------------------------------------
%  One name per idea, used by every book. Anything a book can use is
%  either an unprefixed learning box (keyidea, pitfall, labproject,
%  ...) or carries the utp prefix (utpfigure, utptable, utpfig,
%  utpsteps, utppython, UTPNavy, utpcell, ...). Course prefixes such
%  as ml-, ds- or erp- are not used: a chapter moved from one book to
%  another compiles unchanged. .shared/README.md lists every name.
%
%  Adding a new course? Copy .shared/TEMPLATE-main.tex and follow
%  .shared/ADDING-A-COURSE.md. Every step is a numbered checklist.
%
%  ------------------------------------------------------------------
%  CONTENTS
%  ------------------------------------------------------------------
%    1  Course configuration     the knobs a book may set
%    2  Packages                 the shared package set and fonts
%    3  Colour palette           the one UTP palette
%    4  Page layout and tables   margins, lists, table styling
%    5  Table of contents        parts, chapters, labs, figures
%    6  Headings                 chapter badge, part pages, sections
%    7  Page furniture           running heads and footers
%    8  Learning boxes           key idea, pitfall, try it, and friends
%    9  Small commands           \term, \cmd, \chapterguide, utpsteps
%   10  Diagram node styles      the TikZ vocabulary every book shares
%   11  Figure frames            utpfigure, utpfig, source images
%   12  Lab sheets               the lab banner and the lab toolkit
%   13  Stage divider            the part-opening page in longer books
%   14  Front matter             \utpfrontmatter: cover, contents, map
%   15  Listing styles           code, terminal and pseudocode
% =====================================================================

% --------------------------------------------------------------------
%  1.  COURSE CONFIGURATION
% --------------------------------------------------------------------
%  The knobs below, plus \utpstyleparts in section 5 and the box titles
%  in section 8, are the whole configuration surface. A book sets what
%  it needs BEFORE the \input and inherits the rest.
%
%  \providecommand means "define unless the book already did", which is
%  what makes the \def-before-\input pattern work.
%
%  Running text
%    \utprunninghead   the left running head on every page.
%    \utpchapterlabel  the word in the chapter badge. CHAPTER by
%                      default; a book whose chapters are labs sets LAB.
%    \utplabtotal      the N in the "Lab 3 of N" banners.

\providecommand{\utprunninghead}{Course Notes}
\providecommand{\utpchapterlabel}{CHAPTER}
\providecommand{\utplabtotal}{10}

%  Cover, printed by section 14
%    \utpcoursetitle       the title on the title page.
%    \utpdocumentsubtitle  the line under the title; Teaching Assistant Guide by default.
%    \utpcourseoutline     the sentence or two below it; \\ breaks a line.
%    \utpoutlineheading    the small heading above that sentence.
%    \utpauthor            the name in the PREPARED BY block.
%    \utpauthorrole        the line under the name.

\providecommand{\utpcoursetitle}{\utprunninghead}
\providecommand{\utpdocumentsubtitle}{Teaching Assistant Guide}
\providecommand{\utpcourseoutline}{}
\providecommand{\utpoutlineheading}{COURSE OUTLINE}
\providecommand{\utpauthor}{Reeyan Afzal Khan}
\providecommand{\utpauthorrole}{Graduate Research Assistant}

%  Files and folders
%    \utpextragraphicspath  image folders to search before the standard
%                           ones, as brace groups: {{screenshots/}}.
%                           A course's own figures/ folder and the
%                           shared logos are always searched.
%    \utpcoursemap          the overview chapter \utpfrontmatter inputs
%                           after the contents. Empty means none.

\providecommand{\utpextragraphicspath}{}
\providecommand{\utpcoursemap}{chapters/course-map}

%  Contents entries for chapters (section 5 applies them). A book whose
%  chapters are labs says so once here rather than restating the three
%  tocloft settings after the input.
%    \utpchaptertocword      Chapter~ by default, Lab~ in one book.
%    \utpchaptertocindent    left indent of the entry.
%    \utpchaptertocnumwidth  room reserved for the word and number.

\providecommand{\utpchaptertocword}{Chapter~}
\providecommand{\utpchaptertocindent}{1.5em}
\providecommand{\utpchaptertocnumwidth}{6em}

%  Monospace face (section 2)
%    \utpmonofont   sourcecodepro by default. inconsolata is the
%                   alternative, for a book whose pages are mostly
%                   device transcripts.

\providecommand{\utpmonofont}{sourcecodepro}

% --------------------------------------------------------------------
%  2.  PACKAGES
% --------------------------------------------------------------------
%  The union of what the books need. A course that needs something
%  extra loads it in its own file, under the heading "packages only
%  this book needs". Pseudocode is the one optional package the shared
%  file knows about: \utpusealgorithms in section 15 loads and styles it.
%
%  The tikz and pgfplots library lists are the union across all books,
%  so a diagram or chart copied from one course compiles in another.

\usepackage[T1]{fontenc}
\IfFileExists{glyphtounicode.tex}{\input{glyphtounicode}\pdfgentounicode=1}{}
\usepackage[american]{babel}
\IfFileExists{sourcesanspro.sty}
  {\usepackage[default]{sourcesanspro}}
  {\usepackage[scaled]{helvet}\renewcommand{\familydefault}{\sfdefault}}
\DeclareFontShape{T1}{cmss}{b}{n}{<->ssub*cmss/bx/n}{}
\DeclareFontShape{TS1}{cmss}{b}{n}{<->ssub*cmss/bx/n}{}
% Monospace. Without this the tt default stays Computer Modern Typewriter,
% a serif face that clashes with the sans body text in every listing and
% every \cmd. Source Code Pro is the companion to Source Sans Pro; it is
% scaled down because mono digits otherwise sit large next to the sans.
%
% Inconsolata, chosen with \def\utpmonofont{inconsolata}, distinguishes
% 0/O and 1/l/I at a glance, which matters when a reader copies an
% interface name or an address out of a transcript. Only one mono package
% is ever loaded: a second would win the family but inherit the first's
% series names, and every \ttfamily would then resolve by substitution.
\def\utpmonoinconsolata{inconsolata}
\ifx\utpmonofont\utpmonoinconsolata
  \IfFileExists{inconsolata.sty}{\usepackage[varqu,varl,scaled=1.02]{inconsolata}}{}
\else
  \IfFileExists{sourcecodepro.sty}
    {\usepackage[scaled=0.88]{sourcecodepro}}
    {\usepackage[scaled=0.86]{beramono}}
  % Source Sans names its normal weight "regular", but the Computer Modern
  % fallbacks call it "m". Map one onto the other so a \ttfamily reached
  % while the regular series is active resolves instead of warning.
  \DeclareFontShape{T1}{cmtt}{regular}{n}{<->ssub*cmtt/m/n}{}
  \DeclareFontShape{TS1}{cmtt}{regular}{n}{<->ssub*cmtt/m/n}{}
  \DeclareFontShape{T1}{cmtt}{regular}{it}{<->ssub*cmtt/m/it}{}
\fi
\usepackage[margin=21mm,headheight=22pt]{geometry}
\usepackage{amsmath,amssymb,mathtools,bm}
\usepackage{microtype}
\usepackage{enumitem}
\usepackage{booktabs}
\usepackage{tabularx}
\usepackage{array}
\usepackage{ragged2e}
\usepackage{etoolbox}
\usepackage{titlesec}
\usepackage{tocloft}
\usepackage{fancyhdr}
\usepackage[table]{xcolor}
\usepackage{tcolorbox}
\tcbuselibrary{breakable,skins}
\usepackage[hidelinks]{hyperref}
\usepackage{bookmark}
% Document properties, so the PDF identifies itself in a reader's
% title bar, in a library catalogue, and in search results. Set from
% the same configuration macros the cover uses.
\hypersetup{
  pdftitle={\utpcoursetitle},
  pdfauthor={\utpauthor},
  pdfsubject={Universiti Teknologi PETRONAS course notes and laboratory guide},
  pdfkeywords={UTP, course notes, laboratory manual, \utpcoursetitle},
  pdfdisplaydoctitle=true,
  bookmarksnumbered=true,
  bookmarksopenlevel=1
}
\usepackage{listings}
\usepackage{graphicx}
\usepackage{caption}
\usepackage{tikz}
\usetikzlibrary{angles,arrows.meta,backgrounds,calc,chains,decorations.pathmorphing,decorations.pathreplacing,fit,matrix,patterns,positioning,quotes,shadows.blur,shapes.geometric,shapes.misc,shapes.symbols}
\usepackage{pgfplots}
% statistics gives boxplot prepared; groupplots gives multi-panel charts.
\usepgfplotslibrary{statistics}
\usepgfplotslibrary{groupplots}

% Images: any folders the book names first, then its own figures/
% folder, then the shared logos used by the cover.
\expandafter\graphicspath\expandafter{\utpextragraphicspath{figures/}{../.shared/logos/}}

% --------------------------------------------------------------------
%  3.  COLOUR PALETTE
% --------------------------------------------------------------------
%  One palette, defined once under UTP* names, and used by every book,
%  box, diagram and chart. Five accents (blue, teal, green, orange, red)
%  with a pale tint of each for box backgrounds, plus navy for headings,
%  gold for the cover, ink and muted for text, and a sixth accent,
%  purple, for charts and layered diagrams that need one more series.

\definecolor{UTPNavy}{HTML}{18324A}
\definecolor{UTPBlue}{HTML}{246B9E}
\definecolor{UTPTeal}{HTML}{138A80}
\definecolor{UTPGreen}{HTML}{2D7D46}
\definecolor{UTPOrange}{HTML}{C86B24}
\definecolor{UTPRed}{HTML}{B34843}
\definecolor{UTPPurple}{HTML}{6B4FA0}
\definecolor{UTPGold}{HTML}{D4A13E}
\definecolor{UTPInk}{HTML}{263238}
\definecolor{UTPMuted}{HTML}{60717C}
\definecolor{UTPPaleBlue}{HTML}{EEF6FB}
\definecolor{UTPPaleTeal}{HTML}{EEF8F6}
\definecolor{UTPPaleOrange}{HTML}{FFF5EA}
\definecolor{UTPPaleRed}{HTML}{FFF0EF}
\definecolor{UTPPalePurple}{HTML}{F0ECF9}
\definecolor{UTPPaleGray}{HTML}{F4F6F7}
\color{UTPInk}

% --------------------------------------------------------------------
%  4.  PAGE LAYOUT AND TABLES
% --------------------------------------------------------------------
%  Paragraph spacing, list indents, and the table look: alternating
%  row tint, top-aligned paragraph columns, ragged-right cells.
%
%  Two extra column types for tabularx tables:
%      D{width}   a fixed-width, ragged-right, vertically centred column
%      Y          a ragged-right column that fills the remaining width
%
%  Every table in every book has the same shape, the one Algorithms &
%  Data Structures set: a utptable frame around a full-width tabularx,
%  booktabs rules, a tinted header row of \utpth cells (centred, bold,
%  navy), and a first column in bold navy when it names the rows.
%  Cells are vertically centred, so use m{} rather than p{} columns.
%
%      \begin{utptable}
%      \begin{tabularx}{\linewidth}{@{}>{\bfseries\leavevmode\color{UTPNavy}}m{0.22\linewidth}X@{}}
%      \toprule
%      \rowcolor{UTPPaleBlue!92}
%      \utpth{Operation} & \utpth{What it does} \\
%      \midrule
%      Insertion & Add a new element. \\
%      \bottomrule
%      \end{tabularx}
%      \end{utptable}
%
%  A book whose tables carry dense text asks for \utpcompacttables after
%  the input: tabularx tables then set in \footnotesize with tighter
%  padding and vertically centred cells.

\setlength{\parindent}{0pt}
\setlength{\parskip}{0.55em}
\setlist[itemize]{leftmargin=1.45em,itemsep=0.22em,topsep=0.25em}
\setlist[enumerate]{leftmargin=1.75em,itemsep=0.28em,topsep=0.3em}
\setlist[description]{font=\bfseries\color{UTPNavy},itemsep=0.25em,topsep=0.3em}
\renewcommand{\arraystretch}{1.34}
\setlength{\tabcolsep}{6.5pt}
\arrayrulecolor{UTPBlue!45}
% Keep paragraph-style columns ragged and vertically centred. Tables use
% fixed m columns beside flexible X columns; giving X the same centring
% keeps a short cell level with a tall neighbour in a multi-line row.
\renewcommand{\tabularxcolumn}[1]{m{#1}}
\makeatletter
\g@addto@macro\@arrayparboxrestore{\RaggedRight\arraybackslash}
\makeatother
% A header cell: centred, bold and navy, on the tinted header row.
\newcommand{\utpth}[1]{\multicolumn{1}{c}{\bfseries\color{UTPNavy}#1}}
\AtBeginEnvironment{tabular}{\small\setlength{\extrarowheight}{0.6pt}\rowcolors{2}{UTPPaleBlue!72}{white}}
\AtBeginEnvironment{tabularx}{\small\setlength{\extrarowheight}{0.6pt}\rowcolors{2}{UTPPaleBlue!72}{white}}
\AtBeginEnvironment{figure}{\centering}
\newcolumntype{D}[1]{>{\raggedright\arraybackslash}m{#1}}
\newcolumntype{Y}{>{\raggedright\arraybackslash}X}
\newcommand{\utpcompacttables}{%
  \renewcommand{\tabularxcolumn}[1]{m{##1}}%
  \AtBeginEnvironment{tabularx}{%
    \footnotesize
    \setlength{\tabcolsep}{4pt}%
    \renewcommand{\arraystretch}{1.25}%
    \setlength{\extrarowheight}{0.6pt}%
  }%
}
\emergencystretch=1.5em
\hbadness=10000
\hfuzz=8pt

% Page-breaking discipline. Without these a paragraph can leave a single
% line stranded at the foot of one page or the head of the next, which is
% the most visible sign of unattended typesetting.
\widowpenalty=10000          % no first line of a paragraph left alone at a page foot
\clubpenalty=10000           % no last line of a paragraph left alone at a page head
\displaywidowpenalty=10000   % the same around displayed equations
\brokenpenalty=10000         % do not break a page after a hyphenated line
\predisplaypenalty=10000     % keep the line before a display with it
% Pages are built from boxes of many heights, so let the last line float
% up rather than stretching the gaps between boxes to force a flush foot.
\raggedbottom
% Prefer a slightly loose line to an overfull one, but do not hyphenate
% freely -- short sans lines look ragged when over-hyphenated.
\tolerance=1200
\hyphenpenalty=200

% Character protrusion and font expansion. Both are what stop a justified
% sans column from showing visible rivers down the page.
\microtypesetup{protrusion=true,expansion=true,tracking=true}

\pgfplotsset{
  compat=1.18,
  every axis/.append style={
    axis line style={draw=UTPMuted!85,line width=0.55pt},
    tick style={draw=UTPMuted!85},
    tick label style={font=\scriptsize\color{UTPMuted}},
    label style={font=\small\color{UTPNavy}},
    title style={font=\small\bfseries\color{UTPNavy}},
    legend style={font=\scriptsize,draw=UTPBlue!22,fill=white,rounded corners=1pt}
  },
  every axis plot/.append style={line width=1.15pt}
}
\tikzset{every picture/.append style={line cap=round,line join=round}}

% --------------------------------------------------------------------
%  5.  TABLE OF CONTENTS
% --------------------------------------------------------------------
%  Formatting for part, chapter and lab entries, and for a list of
%  figures if a book prints one. The lab level is a custom tocloft list
%  so lab banners appear in the contents.
%
%  Part styling is skipped when \utpstyleparts is 0, which is what a
%  book without \part divisions should set.
%
%  The contents page is meant to fill exactly one page, and how loosely
%  it is set depends on how many entries a book has. One line after the
%  input tunes it:
%
%      \utpcontentspage{<base size>}{<part size>}%
%                      {<space before a part>}{<before a chapter>}{<before a lab>}
%      \utpcontentspage{\normalsize}{\large}{1.15em}{0.38em}{0.30em}   % the default

\addto\captionsamerican{\renewcommand{\contentsname}{Learning Path}}
\newlistentry{lab}{toc}{-1}
\providecommand{\utpstyleparts}{1}
\ifnum\utpstyleparts=1
\renewcommand{\cftpartpresnum}{}
\renewcommand{\cftpartaftersnum}{}
\setlength{\cftpartnumwidth}{5.2em}
\setlength{\cftbeforepartskip}{0.38em}
\renewcommand{\cftpartfont}{\normalsize\bfseries\color{UTPNavy}}
\renewcommand{\cftpartpagefont}{\normalsize\bfseries\color{UTPNavy}}
\renewcommand{\cftpartleader}{\hfill}
\fi
\renewcommand{\cftchappresnum}{\utpchaptertocword}
\renewcommand{\cftchapaftersnum}{\quad}
\setlength{\cftchapindent}{\utpchaptertocindent}
\setlength{\cftchapnumwidth}{\utpchaptertocnumwidth}
\setlength{\cftbeforechapskip}{0pt}
\renewcommand{\cftchapfont}{\normalfont}
\renewcommand{\cftchappagefont}{\normalfont}
\renewcommand{\cftchapleader}{\cftdotfill{\cftdotsep}}
\setlength{\cftlabindent}{3em}
\setlength{\cftlabnumwidth}{0em}
\setlength{\cftbeforelabskip}{0.10em}
\renewcommand{\cftlabfont}{\bfseries\color{UTPOrange!90!black}}
\renewcommand{\cftlabpagefont}{\bfseries\color{UTPOrange!90!black}}
\renewcommand{\cftlableader}{\color{UTPOrange!65!UTPMuted}\cftdotfill{\cftlabdotsep}}

% A list of figures, for a book that prints one (\listoffigures).
\renewcommand{\cftfigfont}{\normalfont\small}
\renewcommand{\cftfigpagefont}{\normalfont\small}
\renewcommand{\cftfigpresnum}{Figure~}
\renewcommand{\cftfigaftersnum}{\quad}
\setlength{\cftfigindent}{1.5em}
\setlength{\cftfignumwidth}{5.6em}
\setlength{\cftbeforefigskip}{0pt}
\renewcommand{\cftfigleader}{\cftdotfill{\cftdotsep}}

% The contents page itself, printed by \utpfrontmatter (section 14).
\newcommand{\utpcontentspage}[5]{%
  \def\utptocsize{#1}\def\utptocpartsize{#2}%
  \def\utptocpartskip{#3}\def\utptocchapterskip{#4}\def\utptoclabskip{#5}}
\utpcontentspage{\normalsize}{\large}{1.15em}{0.38em}{0.30em}
\newcommand{\utptableofcontents}{%
  \begingroup
  \utptocsize
  \setlength{\parskip}{0pt}%
  \setlength{\cftbeforepartskip}{\utptocpartskip}%
  \setlength{\cftbeforechapskip}{\utptocchapterskip}%
  \setlength{\cftbeforelabskip}{\utptoclabskip}%
  \renewcommand{\cftpartfont}{\utptocpartsize\bfseries\color{UTPNavy}}%
  \renewcommand{\cftpartpagefont}{\utptocpartsize\bfseries\color{UTPNavy}}%
  \tableofcontents
  \endgroup
}

% --------------------------------------------------------------------
%  6.  HEADINGS
% --------------------------------------------------------------------
%  The chapter badge (a coloured block with a word and the number),
%  part pages, and section spacing.
%
%  The word in the badge is \utpchapterlabel; Data Communication sets
%  it to LAB because each of its chapters is a lab.

\titleformat{\chapter}[display]
  {\bfseries\color{UTPNavy}}
  {\filleft\colorbox{UTPBlue}{\parbox{2.55cm}{\centering\color{white}\Large
    \utpchapterlabel\\[-0.15em]\Huge\thechapter}}}
  {0.65em}{\Huge\raggedright}
  [\vspace{0.25em}{\color{UTPBlue}\titlerule[1.2pt]}]
\renewcommand{\thepart}{Stage~\arabic{part}}
\titleformat{\part}[display]
  {\centering\color{UTPNavy}\bfseries}
  {\Large\color{UTPTeal}\thepart}{0.45em}{\Huge}
  [\vspace{0.55em}\color{UTPTeal}\rule{0.45\linewidth}{1.5pt}]
\titleformat{\section}{\Large\bfseries\color{UTPNavy}}{\thesection}{0.55em}{}
\titleformat{\subsection}{\large\bfseries\color{UTPBlue}}{\thesubsection}{0.5em}{}
\titlespacing*{\chapter}{0pt}{-1.8em}{1.0em}
\titlespacing*{\part}{0pt}{0pt}{1em}
\titlespacing*{\section}{0pt}{1.2em}{0.35em}
\titlespacing*{\subsection}{0pt}{0.95em}{0.2em}

% --------------------------------------------------------------------
%  7.  PAGE FURNITURE
% --------------------------------------------------------------------
%  Running heads and footers. The left header is \utprunninghead;
%  the right one follows the current chapter.

\pagestyle{fancy}
\fancyhf{}
\lhead{\small\color{UTPMuted}\utprunninghead}
\rhead{\small\color{UTPMuted}\nouppercase{\leftmark}}
\cfoot{\small\color{UTPMuted}\thepage}
\renewcommand{\headrulewidth}{0.5pt}
\renewcommand{\headrule}{\hbox to\headwidth{\color{UTPPaleBlue}\leaders\hrule height \headrulewidth\hfill}}
\renewcommand{\chaptermark}[1]{\markboth{#1}{}}
\fancypagestyle{plain}{%
  \fancyhf{}%
  \cfoot{\small\color{UTPMuted}\thepage}%
  \renewcommand{\headrulewidth}{0pt}%
}

% --------------------------------------------------------------------
%  8.  LEARNING BOXES
% --------------------------------------------------------------------
%  The coloured callouts the chapters use. Every title is a macro, so
%  a book can rename a box without restating its geometry:
%
%      \def\utptitlepitfall{Technical note}   % before the input
%
%  To change how a box LOOKS in one book, use one of the named
%  variations at the end of this section, or \renewtcolorbox after the
%  input. To add a box only one book needs, declare it in that book
%  with \utpnewcallout or \utpnewquietbox.

% Box titles -- override any of these before the input line.
\providecommand{\utptitlekeyidea}{Key idea}
\providecommand{\utptitleworked}{Worked example}
\providecommand{\utptitlepitfall}{Common mistake}
\providecommand{\utptitletryit}{Try it yourself}
\providecommand{\utptitlequickcheck}{Check your understanding}
\providecommand{\utptitlerecap}{Chapter recap}
\providecommand{\utptitledeliverable}{Lab deliverable}
\providecommand{\utptitlelabmode}{How much is scaffolded?}
\providecommand{\utptitleintuition}{Plain idea}
\providecommand{\utptitlefirstread}{Core path}
\providecommand{\utptitleglance}{Chapter at a glance}
\providecommand{\utptitledebug}{Debugging task}
\providecommand{\utptitleoptional}{Beyond the core}
\providecommand{\utptitleverification}{Verification checklist}
\providecommand{\utptitletroubleshoot}{Troubleshooting}
\providecommand{\utptitlesourcebasis}{Source basis}

% Explaining an idea
\newtcolorbox{keyidea}{enhanced,breakable,colback=UTPPaleTeal,colframe=UTPTeal, boxrule=0.8pt,arc=2mm,left=2.5mm,right=2.5mm,top=1.8mm,bottom=1.8mm, fonttitle=\bfseries\color{UTPNavy},title=\utptitlekeyidea}
\newtcolorbox{workedexample}{enhanced,breakable,colback=UTPPaleBlue,colframe=UTPBlue, boxrule=0.8pt,arc=2mm,left=2.5mm,right=2.5mm,top=1.8mm,bottom=1.8mm, fonttitle=\bfseries\color{UTPNavy},title=\utptitleworked}
\newtcolorbox{pitfall}{enhanced,breakable,colback=UTPPaleRed,colframe=UTPRed, boxrule=0.8pt,arc=2mm,left=2.5mm,right=2.5mm,top=1.8mm,bottom=1.8mm, fonttitle=\bfseries\color{UTPNavy},title=\utptitlepitfall}
% Asking the reader to do something
\newtcolorbox{tryit}{enhanced,breakable,colback=white,colframe=UTPTeal, boxrule=0.9pt,arc=2mm,left=2.5mm,right=2.5mm,top=1.8mm,bottom=1.8mm, fonttitle=\bfseries\color{white},colbacktitle=UTPTeal,title=\utptitletryit}
\newtcolorbox{quickcheck}{enhanced,breakable,colback=white,colframe=UTPOrange, boxrule=0.9pt,arc=2mm,left=2.5mm,right=2.5mm,top=1.8mm,bottom=1.8mm, fonttitle=\bfseries\color{white},colbacktitle=UTPOrange,title=\utptitlequickcheck}
\newtcolorbox{chapterrecap}{enhanced,breakable,colback=UTPPaleTeal,colframe=UTPGreen, boxrule=0.9pt,arc=2mm,left=2.5mm,right=2.5mm,top=1.8mm,bottom=1.8mm, fonttitle=\bfseries\color{white},colbacktitle=UTPGreen,title=\utptitlerecap}
% Lab structure
\newtcolorbox{labdeliverable}{enhanced,breakable,colback=white,colframe=UTPGreen, boxrule=0.9pt,arc=2mm,left=2.6mm,right=2.6mm,top=1.9mm,bottom=1.9mm, fonttitle=\bfseries\color{white},colbacktitle=UTPGreen,title=\utptitledeliverable}
\newtcolorbox{labmode}{enhanced,breakable,colback=UTPPaleBlue!52,colframe=UTPBlue, boxrule=0.75pt,arc=2mm,left=2.6mm,right=2.6mm,top=1.7mm,bottom=1.7mm, fonttitle=\bfseries\color{white},colbacktitle=UTPBlue,title=\utptitlelabmode}
% Softer, lighter-weight asides
\newtcolorbox{intuition}{enhanced,breakable,colback=white,colframe=UTPOrange!42, boxrule=0.5pt,arc=2mm,left=2.5mm,right=2.5mm,top=1.5mm,bottom=1.5mm, fonttitle=\bfseries\color{UTPOrange!85!black},title=\utptitleintuition}
\newtcolorbox{firstread}{enhanced,breakable,colback=white,colframe=UTPBlue!42, boxrule=0.5pt,arc=2mm,left=2.5mm,right=2.5mm,top=1.5mm,bottom=1.5mm, fonttitle=\bfseries\color{UTPBlue}, title=\utptitlefirstread}
\newtcolorbox{debugtask}{enhanced,breakable,colback=UTPPaleRed!58,colframe=UTPRed, boxrule=0.85pt,arc=2mm,left=2.6mm,right=2.6mm,top=1.9mm,bottom=1.9mm, fonttitle=\bfseries\color{white},colbacktitle=UTPRed,title=\utptitledebug}
\newtcolorbox{analogy}{enhanced,breakable,colback=white,colframe=UTPOrange!38, boxrule=0.5pt,arc=2mm,left=2.5mm,right=2.5mm,top=1.5mm,bottom=1.5mm, fonttitle=\bfseries\color{UTPOrange!85!black},title=Analogy}
\newtcolorbox{coursechoice}{enhanced,breakable,colback=white,colframe=UTPNavy, boxrule=0.8pt,arc=2mm,left=2.8mm,right=2.8mm,top=2mm,bottom=2mm, fonttitle=\bfseries\color{white},colbacktitle=UTPNavy,title=Course choice}
\newtcolorbox{decisionmemo}{enhanced,breakable,colback=UTPPaleOrange!55,colframe=UTPOrange, boxrule=0.85pt,arc=2mm,left=2.6mm,right=2.6mm,top=1.9mm,bottom=1.9mm, fonttitle=\bfseries\color{white},colbacktitle=UTPOrange,title=Decision memo}
\newtcolorbox{libnote}{enhanced,breakable,colback=white,colframe=UTPMuted!48, boxrule=0.5pt,arc=2mm,left=2.5mm,right=2.5mm,top=1.5mm,bottom=1.5mm, fonttitle=\bfseries\color{UTPNavy},title=Reference note}
\newtcolorbox{miniproject}{enhanced,breakable,colback=UTPPaleTeal!52,colframe=UTPGreen, boxrule=0.9pt,arc=2mm,left=2.6mm,right=2.6mm,top=1.9mm,bottom=1.9mm, fonttitle=\bfseries\color{white},colbacktitle=UTPGreen,title=Mini-project}
% Where the material came from: the lecturer's slides, manual or tutorial
% that a chapter is built on.
\newtcolorbox{sourcebasis}{enhanced,breakable,colback=white,colframe=UTPMuted!55, boxrule=0.5pt,arc=2mm,left=2.5mm,right=2.5mm,top=1.5mm,bottom=1.5mm, fonttitle=\bfseries\color{UTPMuted},title=\utptitlesourcebasis}
% Boxes that take a title argument
\newtcolorbox{modelcard}[1]{enhanced,breakable,colback=white,colframe=UTPBlue, boxrule=0.9pt,arc=2mm,left=2.5mm,right=2.5mm,top=1.8mm,bottom=1.8mm, fonttitle=\bfseries\color{white},colbacktitle=UTPBlue,title=#1}
\newtcolorbox{optionaltopic}{enhanced,breakable,colback=white,colframe=UTPMuted!58, boxrule=0.5pt,arc=2mm,left=2.5mm,right=2.5mm,top=1.5mm,bottom=1.5mm, fonttitle=\bfseries\color{UTPMuted},title=\utptitleoptional}
\newtcolorbox{secondpass}{enhanced,breakable,colback=UTPPaleGray!82,colframe=UTPMuted, boxrule=0.75pt,arc=2mm,left=2.6mm,right=2.6mm,top=1.9mm,bottom=1.9mm, fonttitle=\bfseries\color{white},colbacktitle=UTPMuted,title=Second pass}
\newtcolorbox{labproject}[1]{enhanced,breakable,colback=white,colframe=UTPBlue, boxrule=0.9pt,arc=2mm,left=2.6mm,right=2.6mm,top=1.9mm,bottom=1.9mm, fonttitle=\bfseries\color{white},colbacktitle=UTPBlue,title={#1}}
% Evidence that a lab worked, and what to do when it did not
\newtcolorbox{verification}{enhanced,breakable,colback=UTPPaleTeal!65,colframe=UTPGreen, boxrule=0.8pt,arc=2mm,left=2.6mm,right=2.6mm,top=1.8mm,bottom=1.8mm, fonttitle=\bfseries\color{white},colbacktitle=UTPGreen,title=\utptitleverification}
\newtcolorbox{troubleshoot}{enhanced,breakable,colback=UTPPaleOrange!70,colframe=UTPOrange, boxrule=0.8pt,arc=2mm,left=2.6mm,right=2.6mm,top=1.8mm,bottom=1.8mm, fonttitle=\bfseries\color{white},colbacktitle=UTPOrange,title=\utptitletroubleshoot}

% Named variations a book asks for after the input. Each restyles boxes
% that already exist, so the chapters do not change.
%
%   \utpheavycallouts          a solid colour bar behind the titles of
%                              keyidea, workedexample, intuition and
%                              firstread, with a firmer frame.
%   \utpbreakableframes        let utpfigure and utptable split over a
%                              page break, in a lighter blue (section 11).
%   \utplabprojectcolour{...}  the labproject box in another colour, with
%                              a little more room inside.
\newcommand{\utplabprojectcolour}[1]{%
  \renewtcolorbox{labproject}[1]{enhanced,breakable,colback=white,colframe=#1, boxrule=0.9pt,arc=2mm,left=2.7mm,right=2.7mm,top=2mm,bottom=2mm, fonttitle=\bfseries\color{white},colbacktitle=#1,title={##1}}%
}
\newcommand{\utpheavycallouts}{%
  \renewtcolorbox{keyidea}{enhanced,breakable,colback=UTPPaleTeal,colframe=UTPTeal, boxrule=0.85pt,arc=2mm,left=2.5mm,right=2.5mm,top=1.8mm,bottom=1.8mm, fonttitle=\bfseries\color{white},colbacktitle=UTPTeal,title=\utptitlekeyidea}%
  \renewtcolorbox{workedexample}{enhanced,breakable,colback=UTPPaleBlue,colframe=UTPBlue, boxrule=0.8pt,arc=2mm,left=2.5mm,right=2.5mm,top=1.8mm,bottom=1.8mm, fonttitle=\bfseries\color{white},colbacktitle=UTPBlue,title=\utptitleworked}%
  \renewtcolorbox{intuition}{enhanced,breakable,colback=UTPPaleOrange,colframe=UTPOrange, boxrule=0.8pt,arc=2mm,left=2.5mm,right=2.5mm,top=1.8mm,bottom=1.8mm, fonttitle=\bfseries\color{white},colbacktitle=UTPOrange,title=\utptitleintuition}%
  \renewtcolorbox{firstread}{enhanced,breakable,colback=white,colframe=UTPBlue, boxrule=0.9pt,arc=2mm,left=2.5mm,right=2.5mm,top=1.8mm,bottom=1.8mm, fonttitle=\bfseries\color{white},colbacktitle=UTPBlue,title=\utptitlefirstread}%
}

% Declaring a box a single book needs. Both constructors take the house
% geometry and let a trailing key list override any of it, so a new box
% is one line and still matches the others:
%
%   \utpnewcallout{odoobox}{white}{UTPBlue}{Odoo 19 Community}
%   \utpnewcallout[boxrule=0.9pt]{name}{back}{frame}{Title}
%   \utpnewquietbox{referencebox}{UTPMuted!58}{UTPMuted}{Reference note}
%
% A callout is a tinted box under a solid title bar, for something the
% reader has to act on. A quiet box is a light frame with coloured
% title text, for an aside they can read past. If a second book ever
% wants the same box, move its line into this section.
\newcommand{\utpnewcallout}[5][]{%
  \newtcolorbox{#2}{enhanced,breakable,colback=#3,colframe=#4, boxrule=0.8pt,arc=2mm,left=2.6mm,right=2.6mm,top=1.8mm,bottom=1.8mm, fonttitle=\bfseries\color{white},colbacktitle=#4,title={#5},#1}%
}
\newcommand{\utpnewquietbox}[5][]{%
  \newtcolorbox{#2}{enhanced,breakable,colback=white,colframe=#3, boxrule=0.5pt,arc=2mm,left=2.5mm,right=2.5mm,top=1.5mm,bottom=1.5mm, fonttitle=\bfseries\color{#4},title={#5},#1}%
}

% --------------------------------------------------------------------
%  9.  SMALL COMMANDS
% --------------------------------------------------------------------
%  \term for a defined term, \cmd for inline code, \chapref and
%  \secref for cross-references, and the utpsteps numbered-step list.
%
%      \term{overfitting}          a term being defined
%      \cmd{git status}            inline code, in any language
%      \chapref{ch:trees}          "Chapter 12"
%      \secref{sec:pruning}        "Section 12.3"
%      \chapterguide{prior}{outcomes}   the box that opens a chapter
%      \begin{utpsteps} ... \end{utpsteps}  a numbered Step 1/Step 2 list
%
%  The second group is the shorthand the chapters use in running text:
%
%      \figref{fig:tree}           "Figure 4.2"
%      \chapsref{ch:a}{ch:b}       "Chapters 3--5"
%      \lib{pandas}                a library name
%      \cxg{x}, \cxo{x}, \cxr{x}   maths coloured green/orange/red
%      \given                      the bar in P(a \given b)
%      \coretag, \latertag         read-now / come-back-to markers
%      \kagglesource{url}{name}    the dataset a lab is built on

\newcommand{\term}[1]{\textbf{\color{UTPNavy}#1}}
\newcommand{\cmd}[1]{\texttt{#1}}
\newcommand{\chapref}[1]{Chapter~\ref{#1}}
\newcommand{\secref}[1]{Section~\ref{#1}}
\newcommand{\chapterguide}[2]{%
  \begin{tcolorbox}[enhanced,breakable,colback=white,colframe=UTPNavy,boxrule=0.8pt,
    arc=2mm,left=2.5mm,right=2.5mm,top=2mm,bottom=2mm,
    title=\utptitleglance,fonttitle=\bfseries\color{white},colbacktitle=UTPNavy]
    \textbf{Before you start:} #1\par
    \textbf{By the end, you can:} #2
  \end{tcolorbox}}
\newenvironment{utpsteps}
  {\begin{enumerate}[label=\textcolor{UTPBlue}{\bfseries Step~\arabic*:},leftmargin=4.5em,itemsep=0.38em]}
  {\end{enumerate}}

\newcommand{\figref}[1]{Figure~\ref{#1}}
\newcommand{\chapsref}[2]{Chapters~\ref{#1}--\ref{#2}}
\newcommand{\lib}[1]{\textbf{\color{UTPBlue}\texttt{#1}}}
\newcommand{\cxg}[1]{{\color{UTPGreen}\bfseries $#1$}}
\newcommand{\cxo}[1]{{\color{UTPOrange}\bfseries $#1$}}
\newcommand{\cxr}[1]{{\color{UTPRed}\bfseries $#1$}}
\newcommand{\given}{\mathrel{|}}
\DeclareMathOperator*{\argmin}{arg\,min}
\DeclareMathOperator*{\argmax}{arg\,max}
\newcommand{\coretag}{\par\smallskip\noindent\colorbox{UTPPaleBlue}{\strut\small\bfseries\color{UTPBlue}CORE NOW}\par\smallskip}
\newcommand{\latertag}{\par\smallskip\noindent\colorbox{UTPPaleGray}{\strut\small\bfseries\color{UTPMuted}RETURN LATER}\par\smallskip}
\newcommand{\kagglesource}[2]{%
  \par\smallskip\noindent\textbf{Kaggle source:} \href{#1}{#2}\par
}

% --------------------------------------------------------------------
%  10.  DIAGRAM NODE STYLES
% --------------------------------------------------------------------
%  The vocabulary the hand-drawn TikZ diagrams are built from: a cell
%  in an array, the one cell under the cursor, a circular graph node,
%  the arrow between two of them. Every book has the whole set:
%
%      utpcell, utphot, utpdone, utpdim, utpwall        array cells
%      utpnode, utpnodehot, utpnodedone, utpnodedim     graph nodes
%      utparrow, utpptr, utpdash, utpedge               connectors
%      utplabel, utptag, utpbrace                       annotation
%      utpstep, utpstage, utppanel, utpbox              boxes
%
%  A cell is 1.0cm by 0.8cm. A book whose diagrams were drawn to a
%  different grid resizes the cells once, after the input:
%
%      \utpdiagramsize{1.05cm}{0.85cm}
%
%  Larger vocabularies for one kind of drawing live in .shared/diagrams/
%  and are loaded by name, after the input:
%
%      \utpusediagrams{network}         devices, links, packet headers,
%                                       sequence diagrams, state machines
%      \utpusediagrams{process}         business-process flows, documents,
%                                       swimlanes, status chips
%      \utpusediagrams{data-objects}    vectors, lists and data frames
%                                       drawn as cells; flowcharts
%      \utpusediagrams{network,process} several at once

\newcommand{\utpdiagramsize}[2]{%
  \tikzset{
    utpcell/.style={draw=UTPBlue,fill=UTPPaleBlue,minimum width=#1,minimum height=#2,
                    font=\small\bfseries\color{UTPNavy},inner sep=1pt},
    utphot/.style={draw=UTPOrange,fill=UTPPaleOrange,minimum width=#1,minimum height=#2,
                   font=\small\bfseries\color{UTPNavy},line width=1pt,inner sep=1pt},
    utpdone/.style={draw=UTPGreen,fill=UTPPaleTeal,minimum width=#1,minimum height=#2,
                    font=\small\bfseries\color{UTPNavy},inner sep=1pt},
    utpdim/.style={draw=UTPMuted!45,fill=UTPPaleGray,minimum width=#1,minimum height=#2,
                   font=\small\color{UTPMuted},inner sep=1pt},
    utpwall/.style={draw=UTPNavy,fill=UTPNavy!82,minimum width=#1,minimum height=#2,
                    inner sep=1pt},
    utpnode/.style={circle,draw=UTPBlue,fill=UTPPaleBlue,minimum size=8mm,
                    font=\small\bfseries\color{UTPNavy},inner sep=0.5pt},
    utpnodehot/.style={circle,draw=UTPOrange,fill=UTPPaleOrange,minimum size=8mm,line width=1pt,
                       font=\small\bfseries\color{UTPNavy},inner sep=0.5pt},
    utpnodedone/.style={circle,draw=UTPGreen,fill=UTPPaleTeal,minimum size=8mm,
                        font=\small\bfseries\color{UTPNavy},inner sep=0.5pt},
    utpnodedim/.style={circle,draw=UTPMuted!45,fill=UTPPaleGray,minimum size=8mm,
                       font=\small\color{UTPMuted},inner sep=0.5pt},
    utparrow/.style={-{Stealth[length=2.2mm]},thick,draw=UTPTeal},
    utpptr/.style={-{Stealth[length=2mm]},semithick,draw=UTPBlue},
    utpdash/.style={-{Stealth[length=2mm]},semithick,draw=UTPRed,dashed},
    utpedge/.style={thick,draw=UTPTeal},
    utplabel/.style={font=\scriptsize\color{UTPMuted},inner sep=1pt},
    utptag/.style={font=\scriptsize\bfseries\color{UTPOrange},inner sep=1pt},
    utpbrace/.style={decorate,decoration={brace,amplitude=4pt,raise=2pt},draw=UTPMuted,semithick},
    utpstep/.style={draw=UTPNavy,fill=white,rounded corners=1.5mm,minimum height=0.75cm,
                    font=\scriptsize\color{UTPNavy},align=center,inner sep=2pt},
    utpstage/.style={draw=UTPNavy,fill=white,rounded corners=1.5mm,minimum height=0.75cm,
                     font=\scriptsize\color{UTPNavy},align=center,inner sep=2pt},
    utppanel/.style={draw=UTPBlue!30,fill=white,rounded corners=1.5mm,inner sep=2.5mm},
    utpbox/.style={draw=UTPBlue,fill=UTPPaleBlue,rounded corners=2mm,align=center,
                   font=\small\color{UTPNavy},inner sep=2.2mm}
  }%
}
\utpdiagramsize{1.0cm}{0.8cm}

% \@for, not \foreach: \foreach runs each pass in a group, which would
% throw away the styles a library declares.
\makeatletter
\newcommand{\utpusediagrams}[1]{%
  \@for\utp@library:=#1\do{\input{../.shared/diagrams/\utp@library}}%
}
\makeatother

% --------------------------------------------------------------------
%  11.  FIGURE FRAMES
% --------------------------------------------------------------------
%  Three ways to put a drawing, chart or screenshot on the page:
%
%      \begin{utpfigure} ... \end{utpfigure}
%          a centred frame with no caption, for a diagram the text
%          talks about directly. An optional title puts a bar on top:
%          \begin{utpfigure}[Where the VLAN tag is added] ...
%
%      \begin{utptable} ... \end{utptable}
%          the same frame in teal, around a tabular or tabularx.
%
%      \begin{utpfig}[short entry]{The caption.}{fig:label}
%        \begin{tikzpicture} ... \end{tikzpicture}
%      \end{utpfig}
%          a frame with a numbered caption that \figref can point to.
%          The optional first argument is the short form for a list of
%          figures. A page never breaks between the drawing and its
%          caption.
%
%  A screenshot or slide lifted from the lecturer's material goes in a
%  titled frame, so the reader can tell it from a drawing made for the
%  notes; the optional width defaults to 0.9 and 0.78 of the line:
%
%      \sourceimage[0.8\linewidth]{Static-routing topology}{static-topology.png}
%      \compactsourceimage{The same, kept to a smaller height}{static-topology.png}
%
%  \fignote{...} adds a short muted line under a drawing, inside its frame.

\newtcolorbox{utpfigure}[1][]{enhanced,width=\linewidth,colback=white,colframe=UTPBlue!28, boxrule=0.7pt,arc=2mm, left=2.2mm,right=2.2mm,top=2.2mm,bottom=1.8mm, before skip=0.75em,after skip=0.85em, before upper={\centering}, halign=center, fonttitle=\bfseries\color{white},colbacktitle=UTPBlue,title={#1}}
\newtcolorbox{utptable}{enhanced,width=\linewidth,colback=white,colframe=UTPTeal!32, boxrule=0.7pt,arc=2mm, left=2.2mm,right=2.2mm,top=1.8mm,bottom=1.8mm, before skip=0.75em,after skip=0.85em, halign=center}
\newcommand{\utpbreakableframes}{%
  \renewtcolorbox{utpfigure}[1][]{enhanced,breakable,colback=white,colframe=UTPBlue!35, boxrule=0.5pt,arc=2mm,left=2.5mm,right=2.5mm,top=2.2mm,bottom=2.2mm, halign=center, fonttitle=\bfseries\color{white},colbacktitle=UTPBlue,title={##1}}%
  \renewtcolorbox{utptable}{enhanced,breakable,colback=white,colframe=UTPBlue!28, boxrule=0.45pt,arc=2mm,left=2.2mm,right=2.2mm,top=1.8mm,bottom=1.8mm}%
}

\captionsetup{
  hypcap=false,
  font={small},
  labelfont={bf,color=UTPNavy},
  textfont={color=UTPInk},
  labelsep=period,
  justification=centering,
  skip=4pt
}

\newenvironment{utpfig}[3][\relax]
  {\def\utpfigshort{#1}\def\utpfigcap{#2}\def\utpfiglab{#3}%
   \par\addvspace{0.55em}%
   \begin{tcolorbox}[enhanced,breakable,colback=white,colframe=UTPBlue!32,
     boxrule=0.5pt,arc=2mm,left=2.5mm,right=2.5mm,top=2.6mm,bottom=2.2mm]
   \centering}
  {\par\nopagebreak\vspace{0.45em}%
   % The caption is boxed: \captionof starts with a write and a zero
   % rule, and the glue after them would otherwise be a place where the
   % page could break, leaving the caption alone at the top of the next.
   % Inside the box the caption is told it sits below the drawing, which
   % it can no longer see, so the gap stays above it.
   \nointerlineskip
   \vbox{\setlength{\parskip}{0pt}\captionsetup{position=below}%
   \ifx\utpfigshort\relax
     \captionof{figure}{\utpfigcap}%
   \else
     \captionof{figure}[\utpfigshort]{\utpfigcap}%
   \fi
   \label{\utpfiglab}}%
   \end{tcolorbox}%
   \par\addvspace{0.55em}}

\newcommand{\sourceimage}[3][0.9\linewidth]{%
  \begin{utpfigure}[Source visual --- #2]%
    \includegraphics[width=#1,height=0.56\textheight,keepaspectratio]{#3}%
  \end{utpfigure}}
\newcommand{\compactsourceimage}[3][0.78\linewidth]{%
  \begin{utpfigure}[Source visual --- #2]%
    \includegraphics[width=#1,height=0.42\textheight,keepaspectratio]{#3}%
  \end{utpfigure}}

\newcommand{\fignote}[1]{\par\smallskip{\scriptsize\color{UTPMuted}#1\par}}

% --------------------------------------------------------------------
%  12.  LAB SHEETS
% --------------------------------------------------------------------
%  Every lab sheet opens with the banner, a full-width page break that
%  also puts the lab in the contents:
%
%      \labsection{3}{Scheduling Policies}{sec:lab03}
%
%  The "of N" comes from \utplabtotal, so a book states its lab count
%  once in its configuration rather than in every lab file.
%
%  The body of the sheet is a labproject box (section 8), and inside it
%  the same toolkit serves every course:
%
%      \labflow{chapter link}{goal}{what to explain}
%          the three-line brief under the lab title.
%      \labsourcecard{file}{what it does}{what to inspect}
%          the heading over a listing of one supplied source file.
%      \begin{labquestion}{3}{Short title} ... \end{labquestion}
%      \begin{labanswer} ... \expectedoutput ... \end{labanswer}
%          a question as the student meets it, then the worked code, the
%          output it should print, and what that output means.
%      \begin{checkanswers} ... \end{checkanswers}
%          answers to a quickcheck, set apart so they can be skipped.
%      labmode, verification, troubleshoot, labdeliverable (section 8)
%          how much is scaffolded, the evidence it worked, what to do if
%          it did not, and what to hand in.

\newcommand{\labsection}[3]{%
  \clearpage
  \phantomsection
  \addcontentsline{toc}{lab}{Lab #1 of \utplabtotal\space--- #2}%
  \markboth{Lab #1 of \utplabtotal\space--- #2}{}%
  \begin{tcolorbox}[enhanced,width=\linewidth,colback=UTPNavy,colframe=UTPNavy,
    boxrule=0pt,arc=2mm,left=5mm,right=5mm,top=5mm,bottom=5mm,
    before skip=0pt,after skip=1.2em]
    {\large\bfseries\color{UTPTeal!35!white} PRACTICAL MILESTONE\par}
    \vspace{0.25em}
    {\Huge\bfseries\color{white} Lab #1 of \utplabtotal\par}
    \vspace{0.20em}
    {\Large\bfseries\color{white} #2\par}
  \end{tcolorbox}
  \label{#3}%
}

\newcommand{\labflow}[3]{%
  \par\smallskip\noindent
  \textbf{Chapter connection.} #1\par
  \textbf{Goal.} #2\par
  \textbf{Explain.} #3\par\medskip}
\newcommand{\labsourcecard}[3]{%
  \par\medskip\noindent\textbf{#1 --- #2}\par
  \textbf{Inspect.} #3\par\smallskip}

\newtcolorbox{labquestion}[2]{enhanced,breakable,colback=UTPPaleOrange!55,colframe=UTPOrange,
  boxrule=0.8pt,arc=2mm,left=2.6mm,right=2.6mm,top=1.8mm,bottom=1.8mm,
  fonttitle=\bfseries\color{white},colbacktitle=UTPOrange,title={Question #1 --- #2}}
\newtcolorbox{labanswer}{enhanced,breakable,colback=white,colframe=UTPTeal,
  boxrule=0.8pt,arc=2mm,left=2.6mm,right=2.6mm,top=1.8mm,bottom=1.8mm,
  fonttitle=\bfseries\color{white},colbacktitle=UTPTeal,title=Worked answer}
\newcommand{\expectedoutput}{\par\smallskip\noindent\textbf{\color{UTPNavy}Reference output.}\nopagebreak\par\nopagebreak}
\utpnewcallout[boxrule=0.65pt,left=2.5mm,right=2.5mm,top=1.6mm,bottom=1.6mm,colbacktitle=UTPMuted]{checkanswers}{UTPPaleGray!72}{UTPMuted!62}{Quick-check answers}

% --------------------------------------------------------------------
%  13.  STAGE DIVIDER
% --------------------------------------------------------------------
%  A full-page divider used by the longer books to open a part:
%  \substage{3A}{Supervised Workflow}{one-line summary}.
%  Harmless to leave unused.

\newcommand{\substage}[3]{%
  \cleardoublepage
  \thispagestyle{empty}
  \phantomsection
  \addcontentsline{toc}{chapter}{Stage #1 --- #2}
  \begin{center}
    \vspace*{0.20\textheight}
    {\Large\bfseries\color{UTPTeal} Stage #1\par}
    \vspace{0.45em}
    {\Huge\bfseries\color{UTPNavy} #2\par}
    \vspace{0.65em}
    {\color{UTPTeal}\rule{0.46\linewidth}{1.5pt}\par}
    \vspace{1.1em}
    \begin{minipage}{0.76\linewidth}
      \centering\large\color{UTPMuted} #3
    \end{minipage}
  \end{center}
  \cleardoublepage
}

% --------------------------------------------------------------------
%  14.  FRONT MATTER
% --------------------------------------------------------------------
%  The opening every book shares: the cover, the contents in roman
%  numerals, the course map, then chapter one on an arabic page 1.
%
%      \begin{document}
%      \utpfrontmatter
%
%  What it prints comes from the configuration in section 1 --- the
%  cover text from \utpcoursetitle and \utpcourseoutline, the overview
%  chapter from \utpcoursemap --- and the contents page is set as
%  \utpcontentspage (section 5) describes. A book with no overview
%  chapter sets \def\utpcoursemap{} and the step is skipped.
%
%  Page anchors are off across the cover so hyperref does not warn
%  about the unnumbered pages, and back on before the contents links
%  are used.
%
%  The number of front-matter pages is recorded in the .aux file as
%  \utpfrontpages, which is how .shared/build.py turns the page numbers
%  in the contents into physical pages when it splits the book.

\newcommand{\utpfrontmatter}{%
  \hypersetup{pageanchor=false}%
  \input{../.shared/cover}%
  \cleardoublepage
  \pagenumbering{roman}%
  \setcounter{tocdepth}{0}%
  \utptableofcontents
  \newpage
  \hypersetup{pageanchor=true}%
  \ifx\utpcoursemap\empty\else\input{\utpcoursemap}\fi
  \cleardoublepage
  \ifdefined\ReadonlyShipoutCounter
    \immediate\write\csname @auxout\endcsname
      {\string\gdef\string\utpfrontpages{\the\ReadonlyShipoutCounter}}%
  \fi
  \pagenumbering{arabic}%
}

% --------------------------------------------------------------------
%  15.  LISTING STYLES
% --------------------------------------------------------------------
%  Two base looks, and one named style per language built on them, so
%  the colours and frame are the same in every book:
%
%      utpcode          the base for source code
%      utpterminal      the base for commands typed and output printed
%      utppython        Python
%      utpr             R
%      utpc             C
%      utpcpp           C++
%      utpbash          shell scripts and commands
%      utpcnumbered     C with line numbers, for a complete source file
%      utpbashnumbered  shell with line numbers, likewise
%
%      \begin{lstlisting}[style=utppython] ... \end{lstlisting}
%      \lstinputlisting[style=utpcnumbered]{all_experiments/Lab04/fcfs.c}
%
%  A book that sets most of its listings in one style makes it the
%  default once, after the input:  \lstset{style=utpcpp}

\lstdefinestyle{utpcode}{
  basicstyle=\ttfamily\scriptsize,
  keywordstyle=\bfseries\color{UTPBlue},
  commentstyle=\itshape\color{UTPMuted},
  stringstyle=\color{UTPGreen!70!black},
  backgroundcolor=\color{UTPPaleGray!72},
  frame=single,
  rulecolor=\color{UTPBlue!28},
  framesep=5pt,
  breaklines=true,
  breakatwhitespace=false,
  showstringspaces=false,
  keepspaces=true,
  columns=fullflexible,
  upquote=true,
  tabsize=4,
  aboveskip=0.7em,
  belowskip=0.8em
}
\lstdefinestyle{utpterminal}{
  basicstyle=\ttfamily\scriptsize,
  backgroundcolor=\color{UTPPaleGray!72},
  frame=single,
  rulecolor=\color{UTPMuted!35},
  framesep=5pt,
  breaklines=true,
  showstringspaces=false,
  keepspaces=true,
  columns=fullflexible,
  upquote=true,
  aboveskip=0.6em,
  belowskip=0.7em
}

\lstdefinestyle{utppython}{style=utpcode,language=Python}
\lstdefinestyle{utpr}{style=utpcode,language=R,tabsize=2}
\lstdefinestyle{utpc}{style=utpcode,language=C}
\lstdefinestyle{utpcpp}{style=utpcode,language=C++}
\lstdefinestyle{utpbash}{style=utpcode,language=bash,
  rulecolor=\color{UTPMuted!35},aboveskip=0.6em,belowskip=0.7em}
\lstdefinestyle{utpcnumbered}{style=utpc,numbers=left,numberstyle=\tiny\color{UTPMuted},
  numbersep=7pt,xleftmargin=1.7em,framexleftmargin=1.4em}
\lstdefinestyle{utpbashnumbered}{style=utpbash,numbers=left,numberstyle=\tiny\color{UTPMuted},
  numbersep=7pt,xleftmargin=1.7em,framexleftmargin=1.4em}

% Pseudocode. A book that sets algorithms asks for them once, after the
% input, and gets algorithm2e loaded and coloured to match:
%
%     \utpusealgorithms
%
% \algostyle then sets one algorithm environment a size smaller than the
% text around it.
\newcommand{\utpalgorithmstyle}{%
  \SetAlCapNameFnt{\bfseries\color{UTPNavy}}%
  \SetAlCapFnt{\bfseries\color{UTPNavy}}%
  \SetAlgoCaptionSeparator{.}%
  \SetKwInput{KwIn}{Input}%
  \SetKwInput{KwOut}{Output}%
  \SetNlSty{color}{UTPMuted}{}%
  \SetCommentSty{\itshape\color{UTPMuted}}%
  \SetKwSty{bfseries\color{UTPBlue}}%
  \SetArgSty{}%
  \providecommand{\algostyle}{\small\SetAlgoNlRelativeSize{-1}\DontPrintSemicolon\SetAlgoVlined}%
}
\newcommand{\utpusealgorithms}{%
  \usepackage[linesnumbered,ruled,vlined]{algorithm2e}%
  \utpalgorithmstyle
}

%
%      \utpcontentspage{\small}{\normalsize}{0.65em}{0.07em}{0.16em}
%
%      \begin{document}
%      \utpfrontmatter                     % cover, contents, map
%      \part{Orientation}
%      \input{chapters/01-where-to-begin}
%      \end{document}
%
%  Configuration goes BEFORE the \input; overrides go AFTER it.
%  Paths are relative, so compile from inside the course folder.
%
%  ------------------------------------------------------------------
%  NAMING
%  ------------------------------------------------------------------
%  One name per idea, used by every book. Anything a book can use is
%  either an unprefixed learning box (keyidea, pitfall, labproject,
%  ...) or carries the utp prefix (utpfigure, utptable, utpfig,
%  utpsteps, utppython, UTPNavy, utpcell, ...). Course prefixes such
%  as ml-, ds- or erp- are not used: a chapter moved from one book to
%  another compiles unchanged. .shared/README.md lists every name.
%
%  Adding a new course? Copy .shared/TEMPLATE-main.tex and follow
%  .shared/ADDING-A-COURSE.md. Every step is a numbered checklist.
%
%  ------------------------------------------------------------------
%  CONTENTS
%  ------------------------------------------------------------------
%    1  Course configuration     the knobs a book may set
%    2  Packages                 the shared package set and fonts
%    3  Colour palette           the one UTP palette
%    4  Page layout and tables   margins, lists, table styling
%    5  Table of contents        parts, chapters, labs, figures
%    6  Headings                 chapter badge, part pages, sections
%    7  Page furniture           running heads and footers
%    8  Learning boxes           key idea, pitfall, try it, and friends
%    9  Small commands           \term, \cmd, \chapterguide, utpsteps
%   10  Diagram node styles      the TikZ vocabulary every book shares
%   11  Figure frames            utpfigure, utpfig, source images
%   12  Lab sheets               the lab banner and the lab toolkit
%   13  Stage divider            the part-opening page in longer books
%   14  Front matter             \utpfrontmatter: cover, contents, map
%   15  Listing styles           code, terminal and pseudocode
% =====================================================================

% --------------------------------------------------------------------
%  1.  COURSE CONFIGURATION
% --------------------------------------------------------------------
%  The knobs below, plus \utpstyleparts in section 5 and the box titles
%  in section 8, are the whole configuration surface. A book sets what
%  it needs BEFORE the \input and inherits the rest.
%
%  \providecommand means "define unless the book already did", which is
%  what makes the \def-before-\input pattern work.
%
%  Running text
%    \utprunninghead   the left running head on every page.
%    \utpchapterlabel  the word in the chapter badge. CHAPTER by
%                      default; a book whose chapters are labs sets LAB.
%    \utplabtotal      the N in the "Lab 3 of N" banners.

\providecommand{\utprunninghead}{Course Notes}
\providecommand{\utpchapterlabel}{CHAPTER}
\providecommand{\utplabtotal}{10}

%  Cover, printed by section 14
%    \utpcoursetitle       the title on the title page.
%    \utpdocumentsubtitle  the line under the title; Teaching Assistant Guide by default.
%    \utpcourseoutline     the sentence or two below it; \\ breaks a line.
%    \utpoutlineheading    the small heading above that sentence.
%    \utpauthor            the name in the PREPARED BY block.
%    \utpauthorrole        the line under the name.

\providecommand{\utpcoursetitle}{\utprunninghead}
\providecommand{\utpdocumentsubtitle}{Teaching Assistant Guide}
\providecommand{\utpcourseoutline}{}
\providecommand{\utpoutlineheading}{COURSE OUTLINE}
\providecommand{\utpauthor}{Reeyan Afzal Khan}
\providecommand{\utpauthorrole}{Graduate Research Assistant}

%  Files and folders
%    \utpextragraphicspath  image folders to search before the standard
%                           ones, as brace groups: {{screenshots/}}.
%                           A course's own figures/ folder and the
%                           shared logos are always searched.
%    \utpcoursemap          the overview chapter \utpfrontmatter inputs
%                           after the contents. Empty means none.

\providecommand{\utpextragraphicspath}{}
\providecommand{\utpcoursemap}{chapters/course-map}

%  Contents entries for chapters (section 5 applies them). A book whose
%  chapters are labs says so once here rather than restating the three
%  tocloft settings after the input.
%    \utpchaptertocword      Chapter~ by default, Lab~ in one book.
%    \utpchaptertocindent    left indent of the entry.
%    \utpchaptertocnumwidth  room reserved for the word and number.

\providecommand{\utpchaptertocword}{Chapter~}
\providecommand{\utpchaptertocindent}{1.5em}
\providecommand{\utpchaptertocnumwidth}{6em}

%  Monospace face (section 2)
%    \utpmonofont   sourcecodepro by default. inconsolata is the
%                   alternative, for a book whose pages are mostly
%                   device transcripts.

\providecommand{\utpmonofont}{sourcecodepro}

% --------------------------------------------------------------------
%  2.  PACKAGES
% --------------------------------------------------------------------
%  The union of what the books need. A course that needs something
%  extra loads it in its own file, under the heading "packages only
%  this book needs". Pseudocode is the one optional package the shared
%  file knows about: \utpusealgorithms in section 15 loads and styles it.
%
%  The tikz and pgfplots library lists are the union across all books,
%  so a diagram or chart copied from one course compiles in another.

\usepackage[T1]{fontenc}
\IfFileExists{glyphtounicode.tex}{\input{glyphtounicode}\pdfgentounicode=1}{}
\usepackage[american]{babel}
\IfFileExists{sourcesanspro.sty}
  {\usepackage[default]{sourcesanspro}}
  {\usepackage[scaled]{helvet}\renewcommand{\familydefault}{\sfdefault}}
\DeclareFontShape{T1}{cmss}{b}{n}{<->ssub*cmss/bx/n}{}
\DeclareFontShape{TS1}{cmss}{b}{n}{<->ssub*cmss/bx/n}{}
% Monospace. Without this the tt default stays Computer Modern Typewriter,
% a serif face that clashes with the sans body text in every listing and
% every \cmd. Source Code Pro is the companion to Source Sans Pro; it is
% scaled down because mono digits otherwise sit large next to the sans.
%
% Inconsolata, chosen with \def\utpmonofont{inconsolata}, distinguishes
% 0/O and 1/l/I at a glance, which matters when a reader copies an
% interface name or an address out of a transcript. Only one mono package
% is ever loaded: a second would win the family but inherit the first's
% series names, and every \ttfamily would then resolve by substitution.
\def\utpmonoinconsolata{inconsolata}
\ifx\utpmonofont\utpmonoinconsolata
  \IfFileExists{inconsolata.sty}{\usepackage[varqu,varl,scaled=1.02]{inconsolata}}{}
\else
  \IfFileExists{sourcecodepro.sty}
    {\usepackage[scaled=0.88]{sourcecodepro}}
    {\usepackage[scaled=0.86]{beramono}}
  % Source Sans names its normal weight "regular", but the Computer Modern
  % fallbacks call it "m". Map one onto the other so a \ttfamily reached
  % while the regular series is active resolves instead of warning.
  \DeclareFontShape{T1}{cmtt}{regular}{n}{<->ssub*cmtt/m/n}{}
  \DeclareFontShape{TS1}{cmtt}{regular}{n}{<->ssub*cmtt/m/n}{}
  \DeclareFontShape{T1}{cmtt}{regular}{it}{<->ssub*cmtt/m/it}{}
\fi
\usepackage[margin=21mm,headheight=22pt]{geometry}
\usepackage{amsmath,amssymb,mathtools,bm}
\usepackage{microtype}
\usepackage{enumitem}
\usepackage{booktabs}
\usepackage{tabularx}
\usepackage{array}
\usepackage{ragged2e}
\usepackage{etoolbox}
\usepackage{titlesec}
\usepackage{tocloft}
\usepackage{fancyhdr}
\usepackage[table]{xcolor}
\usepackage{tcolorbox}
\tcbuselibrary{breakable,skins}
\usepackage[hidelinks]{hyperref}
\usepackage{bookmark}
% Document properties, so the PDF identifies itself in a reader's
% title bar, in a library catalogue, and in search results. Set from
% the same configuration macros the cover uses.
\hypersetup{
  pdftitle={\utpcoursetitle},
  pdfauthor={\utpauthor},
  pdfsubject={Universiti Teknologi PETRONAS course notes and laboratory guide},
  pdfkeywords={UTP, course notes, laboratory manual, \utpcoursetitle},
  pdfdisplaydoctitle=true,
  bookmarksnumbered=true,
  bookmarksopenlevel=1
}
\usepackage{listings}
\usepackage{graphicx}
\usepackage{caption}
\usepackage{tikz}
\usetikzlibrary{angles,arrows.meta,backgrounds,calc,chains,decorations.pathmorphing,decorations.pathreplacing,fit,matrix,patterns,positioning,quotes,shadows.blur,shapes.geometric,shapes.misc,shapes.symbols}
\usepackage{pgfplots}
% statistics gives boxplot prepared; groupplots gives multi-panel charts.
\usepgfplotslibrary{statistics}
\usepgfplotslibrary{groupplots}

% Images: any folders the book names first, then its own figures/
% folder, then the shared logos used by the cover.
\expandafter\graphicspath\expandafter{\utpextragraphicspath{figures/}{../.shared/logos/}}

% --------------------------------------------------------------------
%  3.  COLOUR PALETTE
% --------------------------------------------------------------------
%  One palette, defined once under UTP* names, and used by every book,
%  box, diagram and chart. Five accents (blue, teal, green, orange, red)
%  with a pale tint of each for box backgrounds, plus navy for headings,
%  gold for the cover, ink and muted for text, and a sixth accent,
%  purple, for charts and layered diagrams that need one more series.

\definecolor{UTPNavy}{HTML}{18324A}
\definecolor{UTPBlue}{HTML}{246B9E}
\definecolor{UTPTeal}{HTML}{138A80}
\definecolor{UTPGreen}{HTML}{2D7D46}
\definecolor{UTPOrange}{HTML}{C86B24}
\definecolor{UTPRed}{HTML}{B34843}
\definecolor{UTPPurple}{HTML}{6B4FA0}
\definecolor{UTPGold}{HTML}{D4A13E}
\definecolor{UTPInk}{HTML}{263238}
\definecolor{UTPMuted}{HTML}{60717C}
\definecolor{UTPPaleBlue}{HTML}{EEF6FB}
\definecolor{UTPPaleTeal}{HTML}{EEF8F6}
\definecolor{UTPPaleOrange}{HTML}{FFF5EA}
\definecolor{UTPPaleRed}{HTML}{FFF0EF}
\definecolor{UTPPalePurple}{HTML}{F0ECF9}
\definecolor{UTPPaleGray}{HTML}{F4F6F7}
\color{UTPInk}

% --------------------------------------------------------------------
%  4.  PAGE LAYOUT AND TABLES
% --------------------------------------------------------------------
%  Paragraph spacing, list indents, and the table look: alternating
%  row tint, top-aligned paragraph columns, ragged-right cells.
%
%  Two extra column types for tabularx tables:
%      D{width}   a fixed-width, ragged-right, vertically centred column
%      Y          a ragged-right column that fills the remaining width
%
%  Every table in every book has the same shape, the one Algorithms &
%  Data Structures set: a utptable frame around a full-width tabularx,
%  booktabs rules, a tinted header row of \utpth cells (centred, bold,
%  navy), and a first column in bold navy when it names the rows.
%  Cells are vertically centred, so use m{} rather than p{} columns.
%
%      \begin{utptable}
%      \begin{tabularx}{\linewidth}{@{}>{\bfseries\leavevmode\color{UTPNavy}}m{0.22\linewidth}X@{}}
%      \toprule
%      \rowcolor{UTPPaleBlue!92}
%      \utpth{Operation} & \utpth{What it does} \\
%      \midrule
%      Insertion & Add a new element. \\
%      \bottomrule
%      \end{tabularx}
%      \end{utptable}
%
%  A book whose tables carry dense text asks for \utpcompacttables after
%  the input: tabularx tables then set in \footnotesize with tighter
%  padding and vertically centred cells.

\setlength{\parindent}{0pt}
\setlength{\parskip}{0.55em}
\setlist[itemize]{leftmargin=1.45em,itemsep=0.22em,topsep=0.25em}
\setlist[enumerate]{leftmargin=1.75em,itemsep=0.28em,topsep=0.3em}
\setlist[description]{font=\bfseries\color{UTPNavy},itemsep=0.25em,topsep=0.3em}
\renewcommand{\arraystretch}{1.34}
\setlength{\tabcolsep}{6.5pt}
\arrayrulecolor{UTPBlue!45}
% Keep paragraph-style columns ragged and vertically centred. Tables use
% fixed m columns beside flexible X columns; giving X the same centring
% keeps a short cell level with a tall neighbour in a multi-line row.
\renewcommand{\tabularxcolumn}[1]{m{#1}}
\makeatletter
\g@addto@macro\@arrayparboxrestore{\RaggedRight\arraybackslash}
\makeatother
% A header cell: centred, bold and navy, on the tinted header row.
\newcommand{\utpth}[1]{\multicolumn{1}{c}{\bfseries\color{UTPNavy}#1}}
\AtBeginEnvironment{tabular}{\small\setlength{\extrarowheight}{0.6pt}\rowcolors{2}{UTPPaleBlue!72}{white}}
\AtBeginEnvironment{tabularx}{\small\setlength{\extrarowheight}{0.6pt}\rowcolors{2}{UTPPaleBlue!72}{white}}
\AtBeginEnvironment{figure}{\centering}
\newcolumntype{D}[1]{>{\raggedright\arraybackslash}m{#1}}
\newcolumntype{Y}{>{\raggedright\arraybackslash}X}
\newcommand{\utpcompacttables}{%
  \renewcommand{\tabularxcolumn}[1]{m{##1}}%
  \AtBeginEnvironment{tabularx}{%
    \footnotesize
    \setlength{\tabcolsep}{4pt}%
    \renewcommand{\arraystretch}{1.25}%
    \setlength{\extrarowheight}{0.6pt}%
  }%
}
\emergencystretch=1.5em
\hbadness=10000
\hfuzz=8pt

% Page-breaking discipline. Without these a paragraph can leave a single
% line stranded at the foot of one page or the head of the next, which is
% the most visible sign of unattended typesetting.
\widowpenalty=10000          % no first line of a paragraph left alone at a page foot
\clubpenalty=10000           % no last line of a paragraph left alone at a page head
\displaywidowpenalty=10000   % the same around displayed equations
\brokenpenalty=10000         % do not break a page after a hyphenated line
\predisplaypenalty=10000     % keep the line before a display with it
% Pages are built from boxes of many heights, so let the last line float
% up rather than stretching the gaps between boxes to force a flush foot.
\raggedbottom
% Prefer a slightly loose line to an overfull one, but do not hyphenate
% freely -- short sans lines look ragged when over-hyphenated.
\tolerance=1200
\hyphenpenalty=200

% Character protrusion and font expansion. Both are what stop a justified
% sans column from showing visible rivers down the page.
\microtypesetup{protrusion=true,expansion=true,tracking=true}

\pgfplotsset{
  compat=1.18,
  every axis/.append style={
    axis line style={draw=UTPMuted!85,line width=0.55pt},
    tick style={draw=UTPMuted!85},
    tick label style={font=\scriptsize\color{UTPMuted}},
    label style={font=\small\color{UTPNavy}},
    title style={font=\small\bfseries\color{UTPNavy}},
    legend style={font=\scriptsize,draw=UTPBlue!22,fill=white,rounded corners=1pt}
  },
  every axis plot/.append style={line width=1.15pt}
}
\tikzset{every picture/.append style={line cap=round,line join=round}}

% --------------------------------------------------------------------
%  5.  TABLE OF CONTENTS
% --------------------------------------------------------------------
%  Formatting for part, chapter and lab entries, and for a list of
%  figures if a book prints one. The lab level is a custom tocloft list
%  so lab banners appear in the contents.
%
%  Part styling is skipped when \utpstyleparts is 0, which is what a
%  book without \part divisions should set.
%
%  The contents page is meant to fill exactly one page, and how loosely
%  it is set depends on how many entries a book has. One line after the
%  input tunes it:
%
%      \utpcontentspage{<base size>}{<part size>}%
%                      {<space before a part>}{<before a chapter>}{<before a lab>}
%      \utpcontentspage{\normalsize}{\large}{1.15em}{0.38em}{0.30em}   % the default

\addto\captionsamerican{\renewcommand{\contentsname}{Learning Path}}
\newlistentry{lab}{toc}{-1}
\providecommand{\utpstyleparts}{1}
\ifnum\utpstyleparts=1
\renewcommand{\cftpartpresnum}{}
\renewcommand{\cftpartaftersnum}{}
\setlength{\cftpartnumwidth}{5.2em}
\setlength{\cftbeforepartskip}{0.38em}
\renewcommand{\cftpartfont}{\normalsize\bfseries\color{UTPNavy}}
\renewcommand{\cftpartpagefont}{\normalsize\bfseries\color{UTPNavy}}
\renewcommand{\cftpartleader}{\hfill}
\fi
\renewcommand{\cftchappresnum}{\utpchaptertocword}
\renewcommand{\cftchapaftersnum}{\quad}
\setlength{\cftchapindent}{\utpchaptertocindent}
\setlength{\cftchapnumwidth}{\utpchaptertocnumwidth}
\setlength{\cftbeforechapskip}{0pt}
\renewcommand{\cftchapfont}{\normalfont}
\renewcommand{\cftchappagefont}{\normalfont}
\renewcommand{\cftchapleader}{\cftdotfill{\cftdotsep}}
\setlength{\cftlabindent}{3em}
\setlength{\cftlabnumwidth}{0em}
\setlength{\cftbeforelabskip}{0.10em}
\renewcommand{\cftlabfont}{\bfseries\color{UTPOrange!90!black}}
\renewcommand{\cftlabpagefont}{\bfseries\color{UTPOrange!90!black}}
\renewcommand{\cftlableader}{\color{UTPOrange!65!UTPMuted}\cftdotfill{\cftlabdotsep}}

% A list of figures, for a book that prints one (\listoffigures).
\renewcommand{\cftfigfont}{\normalfont\small}
\renewcommand{\cftfigpagefont}{\normalfont\small}
\renewcommand{\cftfigpresnum}{Figure~}
\renewcommand{\cftfigaftersnum}{\quad}
\setlength{\cftfigindent}{1.5em}
\setlength{\cftfignumwidth}{5.6em}
\setlength{\cftbeforefigskip}{0pt}
\renewcommand{\cftfigleader}{\cftdotfill{\cftdotsep}}

% The contents page itself, printed by \utpfrontmatter (section 14).
\newcommand{\utpcontentspage}[5]{%
  \def\utptocsize{#1}\def\utptocpartsize{#2}%
  \def\utptocpartskip{#3}\def\utptocchapterskip{#4}\def\utptoclabskip{#5}}
\utpcontentspage{\normalsize}{\large}{1.15em}{0.38em}{0.30em}
\newcommand{\utptableofcontents}{%
  \begingroup
  \utptocsize
  \setlength{\parskip}{0pt}%
  \setlength{\cftbeforepartskip}{\utptocpartskip}%
  \setlength{\cftbeforechapskip}{\utptocchapterskip}%
  \setlength{\cftbeforelabskip}{\utptoclabskip}%
  \renewcommand{\cftpartfont}{\utptocpartsize\bfseries\color{UTPNavy}}%
  \renewcommand{\cftpartpagefont}{\utptocpartsize\bfseries\color{UTPNavy}}%
  \tableofcontents
  \endgroup
}

% --------------------------------------------------------------------
%  6.  HEADINGS
% --------------------------------------------------------------------
%  The chapter badge (a coloured block with a word and the number),
%  part pages, and section spacing.
%
%  The word in the badge is \utpchapterlabel; Data Communication sets
%  it to LAB because each of its chapters is a lab.

\titleformat{\chapter}[display]
  {\bfseries\color{UTPNavy}}
  {\filleft\colorbox{UTPBlue}{\parbox{2.55cm}{\centering\color{white}\Large
    \utpchapterlabel\\[-0.15em]\Huge\thechapter}}}
  {0.65em}{\Huge\raggedright}
  [\vspace{0.25em}{\color{UTPBlue}\titlerule[1.2pt]}]
\renewcommand{\thepart}{Stage~\arabic{part}}
\titleformat{\part}[display]
  {\centering\color{UTPNavy}\bfseries}
  {\Large\color{UTPTeal}\thepart}{0.45em}{\Huge}
  [\vspace{0.55em}\color{UTPTeal}\rule{0.45\linewidth}{1.5pt}]
\titleformat{\section}{\Large\bfseries\color{UTPNavy}}{\thesection}{0.55em}{}
\titleformat{\subsection}{\large\bfseries\color{UTPBlue}}{\thesubsection}{0.5em}{}
\titlespacing*{\chapter}{0pt}{-1.8em}{1.0em}
\titlespacing*{\part}{0pt}{0pt}{1em}
\titlespacing*{\section}{0pt}{1.2em}{0.35em}
\titlespacing*{\subsection}{0pt}{0.95em}{0.2em}

% --------------------------------------------------------------------
%  7.  PAGE FURNITURE
% --------------------------------------------------------------------
%  Running heads and footers. The left header is \utprunninghead;
%  the right one follows the current chapter.

\pagestyle{fancy}
\fancyhf{}
\lhead{\small\color{UTPMuted}\utprunninghead}
\rhead{\small\color{UTPMuted}\nouppercase{\leftmark}}
\cfoot{\small\color{UTPMuted}\thepage}
\renewcommand{\headrulewidth}{0.5pt}
\renewcommand{\headrule}{\hbox to\headwidth{\color{UTPPaleBlue}\leaders\hrule height \headrulewidth\hfill}}
\renewcommand{\chaptermark}[1]{\markboth{#1}{}}
\fancypagestyle{plain}{%
  \fancyhf{}%
  \cfoot{\small\color{UTPMuted}\thepage}%
  \renewcommand{\headrulewidth}{0pt}%
}

% --------------------------------------------------------------------
%  8.  LEARNING BOXES
% --------------------------------------------------------------------
%  The coloured callouts the chapters use. Every title is a macro, so
%  a book can rename a box without restating its geometry:
%
%      \def\utptitlepitfall{Technical note}   % before the input
%
%  To change how a box LOOKS in one book, use one of the named
%  variations at the end of this section, or \renewtcolorbox after the
%  input. To add a box only one book needs, declare it in that book
%  with \utpnewcallout or \utpnewquietbox.

% Box titles -- override any of these before the input line.
\providecommand{\utptitlekeyidea}{Key idea}
\providecommand{\utptitleworked}{Worked example}
\providecommand{\utptitlepitfall}{Common mistake}
\providecommand{\utptitletryit}{Try it yourself}
\providecommand{\utptitlequickcheck}{Check your understanding}
\providecommand{\utptitlerecap}{Chapter recap}
\providecommand{\utptitledeliverable}{Lab deliverable}
\providecommand{\utptitlelabmode}{How much is scaffolded?}
\providecommand{\utptitleintuition}{Plain idea}
\providecommand{\utptitlefirstread}{Core path}
\providecommand{\utptitleglance}{Chapter at a glance}
\providecommand{\utptitledebug}{Debugging task}
\providecommand{\utptitleoptional}{Beyond the core}
\providecommand{\utptitleverification}{Verification checklist}
\providecommand{\utptitletroubleshoot}{Troubleshooting}
\providecommand{\utptitlesourcebasis}{Source basis}

% Explaining an idea
\newtcolorbox{keyidea}{enhanced,breakable,colback=UTPPaleTeal,colframe=UTPTeal, boxrule=0.8pt,arc=2mm,left=2.5mm,right=2.5mm,top=1.8mm,bottom=1.8mm, fonttitle=\bfseries\color{UTPNavy},title=\utptitlekeyidea}
\newtcolorbox{workedexample}{enhanced,breakable,colback=UTPPaleBlue,colframe=UTPBlue, boxrule=0.8pt,arc=2mm,left=2.5mm,right=2.5mm,top=1.8mm,bottom=1.8mm, fonttitle=\bfseries\color{UTPNavy},title=\utptitleworked}
\newtcolorbox{pitfall}{enhanced,breakable,colback=UTPPaleRed,colframe=UTPRed, boxrule=0.8pt,arc=2mm,left=2.5mm,right=2.5mm,top=1.8mm,bottom=1.8mm, fonttitle=\bfseries\color{UTPNavy},title=\utptitlepitfall}
% Asking the reader to do something
\newtcolorbox{tryit}{enhanced,breakable,colback=white,colframe=UTPTeal, boxrule=0.9pt,arc=2mm,left=2.5mm,right=2.5mm,top=1.8mm,bottom=1.8mm, fonttitle=\bfseries\color{white},colbacktitle=UTPTeal,title=\utptitletryit}
\newtcolorbox{quickcheck}{enhanced,breakable,colback=white,colframe=UTPOrange, boxrule=0.9pt,arc=2mm,left=2.5mm,right=2.5mm,top=1.8mm,bottom=1.8mm, fonttitle=\bfseries\color{white},colbacktitle=UTPOrange,title=\utptitlequickcheck}
\newtcolorbox{chapterrecap}{enhanced,breakable,colback=UTPPaleTeal,colframe=UTPGreen, boxrule=0.9pt,arc=2mm,left=2.5mm,right=2.5mm,top=1.8mm,bottom=1.8mm, fonttitle=\bfseries\color{white},colbacktitle=UTPGreen,title=\utptitlerecap}
% Lab structure
\newtcolorbox{labdeliverable}{enhanced,breakable,colback=white,colframe=UTPGreen, boxrule=0.9pt,arc=2mm,left=2.6mm,right=2.6mm,top=1.9mm,bottom=1.9mm, fonttitle=\bfseries\color{white},colbacktitle=UTPGreen,title=\utptitledeliverable}
\newtcolorbox{labmode}{enhanced,breakable,colback=UTPPaleBlue!52,colframe=UTPBlue, boxrule=0.75pt,arc=2mm,left=2.6mm,right=2.6mm,top=1.7mm,bottom=1.7mm, fonttitle=\bfseries\color{white},colbacktitle=UTPBlue,title=\utptitlelabmode}
% Softer, lighter-weight asides
\newtcolorbox{intuition}{enhanced,breakable,colback=white,colframe=UTPOrange!42, boxrule=0.5pt,arc=2mm,left=2.5mm,right=2.5mm,top=1.5mm,bottom=1.5mm, fonttitle=\bfseries\color{UTPOrange!85!black},title=\utptitleintuition}
\newtcolorbox{firstread}{enhanced,breakable,colback=white,colframe=UTPBlue!42, boxrule=0.5pt,arc=2mm,left=2.5mm,right=2.5mm,top=1.5mm,bottom=1.5mm, fonttitle=\bfseries\color{UTPBlue}, title=\utptitlefirstread}
\newtcolorbox{debugtask}{enhanced,breakable,colback=UTPPaleRed!58,colframe=UTPRed, boxrule=0.85pt,arc=2mm,left=2.6mm,right=2.6mm,top=1.9mm,bottom=1.9mm, fonttitle=\bfseries\color{white},colbacktitle=UTPRed,title=\utptitledebug}
\newtcolorbox{analogy}{enhanced,breakable,colback=white,colframe=UTPOrange!38, boxrule=0.5pt,arc=2mm,left=2.5mm,right=2.5mm,top=1.5mm,bottom=1.5mm, fonttitle=\bfseries\color{UTPOrange!85!black},title=Analogy}
\newtcolorbox{coursechoice}{enhanced,breakable,colback=white,colframe=UTPNavy, boxrule=0.8pt,arc=2mm,left=2.8mm,right=2.8mm,top=2mm,bottom=2mm, fonttitle=\bfseries\color{white},colbacktitle=UTPNavy,title=Course choice}
\newtcolorbox{decisionmemo}{enhanced,breakable,colback=UTPPaleOrange!55,colframe=UTPOrange, boxrule=0.85pt,arc=2mm,left=2.6mm,right=2.6mm,top=1.9mm,bottom=1.9mm, fonttitle=\bfseries\color{white},colbacktitle=UTPOrange,title=Decision memo}
\newtcolorbox{libnote}{enhanced,breakable,colback=white,colframe=UTPMuted!48, boxrule=0.5pt,arc=2mm,left=2.5mm,right=2.5mm,top=1.5mm,bottom=1.5mm, fonttitle=\bfseries\color{UTPNavy},title=Reference note}
\newtcolorbox{miniproject}{enhanced,breakable,colback=UTPPaleTeal!52,colframe=UTPGreen, boxrule=0.9pt,arc=2mm,left=2.6mm,right=2.6mm,top=1.9mm,bottom=1.9mm, fonttitle=\bfseries\color{white},colbacktitle=UTPGreen,title=Mini-project}
% Where the material came from: the lecturer's slides, manual or tutorial
% that a chapter is built on.
\newtcolorbox{sourcebasis}{enhanced,breakable,colback=white,colframe=UTPMuted!55, boxrule=0.5pt,arc=2mm,left=2.5mm,right=2.5mm,top=1.5mm,bottom=1.5mm, fonttitle=\bfseries\color{UTPMuted},title=\utptitlesourcebasis}
% Boxes that take a title argument
\newtcolorbox{modelcard}[1]{enhanced,breakable,colback=white,colframe=UTPBlue, boxrule=0.9pt,arc=2mm,left=2.5mm,right=2.5mm,top=1.8mm,bottom=1.8mm, fonttitle=\bfseries\color{white},colbacktitle=UTPBlue,title=#1}
\newtcolorbox{optionaltopic}{enhanced,breakable,colback=white,colframe=UTPMuted!58, boxrule=0.5pt,arc=2mm,left=2.5mm,right=2.5mm,top=1.5mm,bottom=1.5mm, fonttitle=\bfseries\color{UTPMuted},title=\utptitleoptional}
\newtcolorbox{secondpass}{enhanced,breakable,colback=UTPPaleGray!82,colframe=UTPMuted, boxrule=0.75pt,arc=2mm,left=2.6mm,right=2.6mm,top=1.9mm,bottom=1.9mm, fonttitle=\bfseries\color{white},colbacktitle=UTPMuted,title=Second pass}
\newtcolorbox{labproject}[1]{enhanced,breakable,colback=white,colframe=UTPBlue, boxrule=0.9pt,arc=2mm,left=2.6mm,right=2.6mm,top=1.9mm,bottom=1.9mm, fonttitle=\bfseries\color{white},colbacktitle=UTPBlue,title={#1}}
% Evidence that a lab worked, and what to do when it did not
\newtcolorbox{verification}{enhanced,breakable,colback=UTPPaleTeal!65,colframe=UTPGreen, boxrule=0.8pt,arc=2mm,left=2.6mm,right=2.6mm,top=1.8mm,bottom=1.8mm, fonttitle=\bfseries\color{white},colbacktitle=UTPGreen,title=\utptitleverification}
\newtcolorbox{troubleshoot}{enhanced,breakable,colback=UTPPaleOrange!70,colframe=UTPOrange, boxrule=0.8pt,arc=2mm,left=2.6mm,right=2.6mm,top=1.8mm,bottom=1.8mm, fonttitle=\bfseries\color{white},colbacktitle=UTPOrange,title=\utptitletroubleshoot}

% Named variations a book asks for after the input. Each restyles boxes
% that already exist, so the chapters do not change.
%
%   \utpheavycallouts          a solid colour bar behind the titles of
%                              keyidea, workedexample, intuition and
%                              firstread, with a firmer frame.
%   \utpbreakableframes        let utpfigure and utptable split over a
%                              page break, in a lighter blue (section 11).
%   \utplabprojectcolour{...}  the labproject box in another colour, with
%                              a little more room inside.
\newcommand{\utplabprojectcolour}[1]{%
  \renewtcolorbox{labproject}[1]{enhanced,breakable,colback=white,colframe=#1, boxrule=0.9pt,arc=2mm,left=2.7mm,right=2.7mm,top=2mm,bottom=2mm, fonttitle=\bfseries\color{white},colbacktitle=#1,title={##1}}%
}
\newcommand{\utpheavycallouts}{%
  \renewtcolorbox{keyidea}{enhanced,breakable,colback=UTPPaleTeal,colframe=UTPTeal, boxrule=0.85pt,arc=2mm,left=2.5mm,right=2.5mm,top=1.8mm,bottom=1.8mm, fonttitle=\bfseries\color{white},colbacktitle=UTPTeal,title=\utptitlekeyidea}%
  \renewtcolorbox{workedexample}{enhanced,breakable,colback=UTPPaleBlue,colframe=UTPBlue, boxrule=0.8pt,arc=2mm,left=2.5mm,right=2.5mm,top=1.8mm,bottom=1.8mm, fonttitle=\bfseries\color{white},colbacktitle=UTPBlue,title=\utptitleworked}%
  \renewtcolorbox{intuition}{enhanced,breakable,colback=UTPPaleOrange,colframe=UTPOrange, boxrule=0.8pt,arc=2mm,left=2.5mm,right=2.5mm,top=1.8mm,bottom=1.8mm, fonttitle=\bfseries\color{white},colbacktitle=UTPOrange,title=\utptitleintuition}%
  \renewtcolorbox{firstread}{enhanced,breakable,colback=white,colframe=UTPBlue, boxrule=0.9pt,arc=2mm,left=2.5mm,right=2.5mm,top=1.8mm,bottom=1.8mm, fonttitle=\bfseries\color{white},colbacktitle=UTPBlue,title=\utptitlefirstread}%
}

% Declaring a box a single book needs. Both constructors take the house
% geometry and let a trailing key list override any of it, so a new box
% is one line and still matches the others:
%
%   \utpnewcallout{odoobox}{white}{UTPBlue}{Odoo 19 Community}
%   \utpnewcallout[boxrule=0.9pt]{name}{back}{frame}{Title}
%   \utpnewquietbox{referencebox}{UTPMuted!58}{UTPMuted}{Reference note}
%
% A callout is a tinted box under a solid title bar, for something the
% reader has to act on. A quiet box is a light frame with coloured
% title text, for an aside they can read past. If a second book ever
% wants the same box, move its line into this section.
\newcommand{\utpnewcallout}[5][]{%
  \newtcolorbox{#2}{enhanced,breakable,colback=#3,colframe=#4, boxrule=0.8pt,arc=2mm,left=2.6mm,right=2.6mm,top=1.8mm,bottom=1.8mm, fonttitle=\bfseries\color{white},colbacktitle=#4,title={#5},#1}%
}
\newcommand{\utpnewquietbox}[5][]{%
  \newtcolorbox{#2}{enhanced,breakable,colback=white,colframe=#3, boxrule=0.5pt,arc=2mm,left=2.5mm,right=2.5mm,top=1.5mm,bottom=1.5mm, fonttitle=\bfseries\color{#4},title={#5},#1}%
}

% --------------------------------------------------------------------
%  9.  SMALL COMMANDS
% --------------------------------------------------------------------
%  \term for a defined term, \cmd for inline code, \chapref and
%  \secref for cross-references, and the utpsteps numbered-step list.
%
%      \term{overfitting}          a term being defined
%      \cmd{git status}            inline code, in any language
%      \chapref{ch:trees}          "Chapter 12"
%      \secref{sec:pruning}        "Section 12.3"
%      \chapterguide{prior}{outcomes}   the box that opens a chapter
%      \begin{utpsteps} ... \end{utpsteps}  a numbered Step 1/Step 2 list
%
%  The second group is the shorthand the chapters use in running text:
%
%      \figref{fig:tree}           "Figure 4.2"
%      \chapsref{ch:a}{ch:b}       "Chapters 3--5"
%      \lib{pandas}                a library name
%      \cxg{x}, \cxo{x}, \cxr{x}   maths coloured green/orange/red
%      \given                      the bar in P(a \given b)
%      \coretag, \latertag         read-now / come-back-to markers
%      \kagglesource{url}{name}    the dataset a lab is built on

\newcommand{\term}[1]{\textbf{\color{UTPNavy}#1}}
\newcommand{\cmd}[1]{\texttt{#1}}
\newcommand{\chapref}[1]{Chapter~\ref{#1}}
\newcommand{\secref}[1]{Section~\ref{#1}}
\newcommand{\chapterguide}[2]{%
  \begin{tcolorbox}[enhanced,breakable,colback=white,colframe=UTPNavy,boxrule=0.8pt,
    arc=2mm,left=2.5mm,right=2.5mm,top=2mm,bottom=2mm,
    title=\utptitleglance,fonttitle=\bfseries\color{white},colbacktitle=UTPNavy]
    \textbf{Before you start:} #1\par
    \textbf{By the end, you can:} #2
  \end{tcolorbox}}
\newenvironment{utpsteps}
  {\begin{enumerate}[label=\textcolor{UTPBlue}{\bfseries Step~\arabic*:},leftmargin=4.5em,itemsep=0.38em]}
  {\end{enumerate}}

\newcommand{\figref}[1]{Figure~\ref{#1}}
\newcommand{\chapsref}[2]{Chapters~\ref{#1}--\ref{#2}}
\newcommand{\lib}[1]{\textbf{\color{UTPBlue}\texttt{#1}}}
\newcommand{\cxg}[1]{{\color{UTPGreen}\bfseries $#1$}}
\newcommand{\cxo}[1]{{\color{UTPOrange}\bfseries $#1$}}
\newcommand{\cxr}[1]{{\color{UTPRed}\bfseries $#1$}}
\newcommand{\given}{\mathrel{|}}
\DeclareMathOperator*{\argmin}{arg\,min}
\DeclareMathOperator*{\argmax}{arg\,max}
\newcommand{\coretag}{\par\smallskip\noindent\colorbox{UTPPaleBlue}{\strut\small\bfseries\color{UTPBlue}CORE NOW}\par\smallskip}
\newcommand{\latertag}{\par\smallskip\noindent\colorbox{UTPPaleGray}{\strut\small\bfseries\color{UTPMuted}RETURN LATER}\par\smallskip}
\newcommand{\kagglesource}[2]{%
  \par\smallskip\noindent\textbf{Kaggle source:} \href{#1}{#2}\par
}

% --------------------------------------------------------------------
%  10.  DIAGRAM NODE STYLES
% --------------------------------------------------------------------
%  The vocabulary the hand-drawn TikZ diagrams are built from: a cell
%  in an array, the one cell under the cursor, a circular graph node,
%  the arrow between two of them. Every book has the whole set:
%
%      utpcell, utphot, utpdone, utpdim, utpwall        array cells
%      utpnode, utpnodehot, utpnodedone, utpnodedim     graph nodes
%      utparrow, utpptr, utpdash, utpedge               connectors
%      utplabel, utptag, utpbrace                       annotation
%      utpstep, utpstage, utppanel, utpbox              boxes
%
%  A cell is 1.0cm by 0.8cm. A book whose diagrams were drawn to a
%  different grid resizes the cells once, after the input:
%
%      \utpdiagramsize{1.05cm}{0.85cm}
%
%  Larger vocabularies for one kind of drawing live in .shared/diagrams/
%  and are loaded by name, after the input:
%
%      \utpusediagrams{network}         devices, links, packet headers,
%                                       sequence diagrams, state machines
%      \utpusediagrams{process}         business-process flows, documents,
%                                       swimlanes, status chips
%      \utpusediagrams{data-objects}    vectors, lists and data frames
%                                       drawn as cells; flowcharts
%      \utpusediagrams{network,process} several at once

\newcommand{\utpdiagramsize}[2]{%
  \tikzset{
    utpcell/.style={draw=UTPBlue,fill=UTPPaleBlue,minimum width=#1,minimum height=#2,
                    font=\small\bfseries\color{UTPNavy},inner sep=1pt},
    utphot/.style={draw=UTPOrange,fill=UTPPaleOrange,minimum width=#1,minimum height=#2,
                   font=\small\bfseries\color{UTPNavy},line width=1pt,inner sep=1pt},
    utpdone/.style={draw=UTPGreen,fill=UTPPaleTeal,minimum width=#1,minimum height=#2,
                    font=\small\bfseries\color{UTPNavy},inner sep=1pt},
    utpdim/.style={draw=UTPMuted!45,fill=UTPPaleGray,minimum width=#1,minimum height=#2,
                   font=\small\color{UTPMuted},inner sep=1pt},
    utpwall/.style={draw=UTPNavy,fill=UTPNavy!82,minimum width=#1,minimum height=#2,
                    inner sep=1pt},
    utpnode/.style={circle,draw=UTPBlue,fill=UTPPaleBlue,minimum size=8mm,
                    font=\small\bfseries\color{UTPNavy},inner sep=0.5pt},
    utpnodehot/.style={circle,draw=UTPOrange,fill=UTPPaleOrange,minimum size=8mm,line width=1pt,
                       font=\small\bfseries\color{UTPNavy},inner sep=0.5pt},
    utpnodedone/.style={circle,draw=UTPGreen,fill=UTPPaleTeal,minimum size=8mm,
                        font=\small\bfseries\color{UTPNavy},inner sep=0.5pt},
    utpnodedim/.style={circle,draw=UTPMuted!45,fill=UTPPaleGray,minimum size=8mm,
                       font=\small\color{UTPMuted},inner sep=0.5pt},
    utparrow/.style={-{Stealth[length=2.2mm]},thick,draw=UTPTeal},
    utpptr/.style={-{Stealth[length=2mm]},semithick,draw=UTPBlue},
    utpdash/.style={-{Stealth[length=2mm]},semithick,draw=UTPRed,dashed},
    utpedge/.style={thick,draw=UTPTeal},
    utplabel/.style={font=\scriptsize\color{UTPMuted},inner sep=1pt},
    utptag/.style={font=\scriptsize\bfseries\color{UTPOrange},inner sep=1pt},
    utpbrace/.style={decorate,decoration={brace,amplitude=4pt,raise=2pt},draw=UTPMuted,semithick},
    utpstep/.style={draw=UTPNavy,fill=white,rounded corners=1.5mm,minimum height=0.75cm,
                    font=\scriptsize\color{UTPNavy},align=center,inner sep=2pt},
    utpstage/.style={draw=UTPNavy,fill=white,rounded corners=1.5mm,minimum height=0.75cm,
                     font=\scriptsize\color{UTPNavy},align=center,inner sep=2pt},
    utppanel/.style={draw=UTPBlue!30,fill=white,rounded corners=1.5mm,inner sep=2.5mm},
    utpbox/.style={draw=UTPBlue,fill=UTPPaleBlue,rounded corners=2mm,align=center,
                   font=\small\color{UTPNavy},inner sep=2.2mm}
  }%
}
\utpdiagramsize{1.0cm}{0.8cm}

% \@for, not \foreach: \foreach runs each pass in a group, which would
% throw away the styles a library declares.
\makeatletter
\newcommand{\utpusediagrams}[1]{%
  \@for\utp@library:=#1\do{\input{../.shared/diagrams/\utp@library}}%
}
\makeatother

% --------------------------------------------------------------------
%  11.  FIGURE FRAMES
% --------------------------------------------------------------------
%  Three ways to put a drawing, chart or screenshot on the page:
%
%      \begin{utpfigure} ... \end{utpfigure}
%          a centred frame with no caption, for a diagram the text
%          talks about directly. An optional title puts a bar on top:
%          \begin{utpfigure}[Where the VLAN tag is added] ...
%
%      \begin{utptable} ... \end{utptable}
%          the same frame in teal, around a tabular or tabularx.
%
%      \begin{utpfig}[short entry]{The caption.}{fig:label}
%        \begin{tikzpicture} ... \end{tikzpicture}
%      \end{utpfig}
%          a frame with a numbered caption that \figref can point to.
%          The optional first argument is the short form for a list of
%          figures. A page never breaks between the drawing and its
%          caption.
%
%  A screenshot or slide lifted from the lecturer's material goes in a
%  titled frame, so the reader can tell it from a drawing made for the
%  notes; the optional width defaults to 0.9 and 0.78 of the line:
%
%      \sourceimage[0.8\linewidth]{Static-routing topology}{static-topology.png}
%      \compactsourceimage{The same, kept to a smaller height}{static-topology.png}
%
%  \fignote{...} adds a short muted line under a drawing, inside its frame.

\newtcolorbox{utpfigure}[1][]{enhanced,width=\linewidth,colback=white,colframe=UTPBlue!28, boxrule=0.7pt,arc=2mm, left=2.2mm,right=2.2mm,top=2.2mm,bottom=1.8mm, before skip=0.75em,after skip=0.85em, before upper={\centering}, halign=center, fonttitle=\bfseries\color{white},colbacktitle=UTPBlue,title={#1}}
\newtcolorbox{utptable}{enhanced,width=\linewidth,colback=white,colframe=UTPTeal!32, boxrule=0.7pt,arc=2mm, left=2.2mm,right=2.2mm,top=1.8mm,bottom=1.8mm, before skip=0.75em,after skip=0.85em, halign=center}
\newcommand{\utpbreakableframes}{%
  \renewtcolorbox{utpfigure}[1][]{enhanced,breakable,colback=white,colframe=UTPBlue!35, boxrule=0.5pt,arc=2mm,left=2.5mm,right=2.5mm,top=2.2mm,bottom=2.2mm, halign=center, fonttitle=\bfseries\color{white},colbacktitle=UTPBlue,title={##1}}%
  \renewtcolorbox{utptable}{enhanced,breakable,colback=white,colframe=UTPBlue!28, boxrule=0.45pt,arc=2mm,left=2.2mm,right=2.2mm,top=1.8mm,bottom=1.8mm}%
}

\captionsetup{
  hypcap=false,
  font={small},
  labelfont={bf,color=UTPNavy},
  textfont={color=UTPInk},
  labelsep=period,
  justification=centering,
  skip=4pt
}

\newenvironment{utpfig}[3][\relax]
  {\def\utpfigshort{#1}\def\utpfigcap{#2}\def\utpfiglab{#3}%
   \par\addvspace{0.55em}%
   \begin{tcolorbox}[enhanced,breakable,colback=white,colframe=UTPBlue!32,
     boxrule=0.5pt,arc=2mm,left=2.5mm,right=2.5mm,top=2.6mm,bottom=2.2mm]
   \centering}
  {\par\nopagebreak\vspace{0.45em}%
   % The caption is boxed: \captionof starts with a write and a zero
   % rule, and the glue after them would otherwise be a place where the
   % page could break, leaving the caption alone at the top of the next.
   % Inside the box the caption is told it sits below the drawing, which
   % it can no longer see, so the gap stays above it.
   \nointerlineskip
   \vbox{\setlength{\parskip}{0pt}\captionsetup{position=below}%
   \ifx\utpfigshort\relax
     \captionof{figure}{\utpfigcap}%
   \else
     \captionof{figure}[\utpfigshort]{\utpfigcap}%
   \fi
   \label{\utpfiglab}}%
   \end{tcolorbox}%
   \par\addvspace{0.55em}}

\newcommand{\sourceimage}[3][0.9\linewidth]{%
  \begin{utpfigure}[Source visual --- #2]%
    \includegraphics[width=#1,height=0.56\textheight,keepaspectratio]{#3}%
  \end{utpfigure}}
\newcommand{\compactsourceimage}[3][0.78\linewidth]{%
  \begin{utpfigure}[Source visual --- #2]%
    \includegraphics[width=#1,height=0.42\textheight,keepaspectratio]{#3}%
  \end{utpfigure}}

\newcommand{\fignote}[1]{\par\smallskip{\scriptsize\color{UTPMuted}#1\par}}

% --------------------------------------------------------------------
%  12.  LAB SHEETS
% --------------------------------------------------------------------
%  Every lab sheet opens with the banner, a full-width page break that
%  also puts the lab in the contents:
%
%      \labsection{3}{Scheduling Policies}{sec:lab03}
%
%  The "of N" comes from \utplabtotal, so a book states its lab count
%  once in its configuration rather than in every lab file.
%
%  The body of the sheet is a labproject box (section 8), and inside it
%  the same toolkit serves every course:
%
%      \labflow{chapter link}{goal}{what to explain}
%          the three-line brief under the lab title.
%      \labsourcecard{file}{what it does}{what to inspect}
%          the heading over a listing of one supplied source file.
%      \begin{labquestion}{3}{Short title} ... \end{labquestion}
%      \begin{labanswer} ... \expectedoutput ... \end{labanswer}
%          a question as the student meets it, then the worked code, the
%          output it should print, and what that output means.
%      \begin{checkanswers} ... \end{checkanswers}
%          answers to a quickcheck, set apart so they can be skipped.
%      labmode, verification, troubleshoot, labdeliverable (section 8)
%          how much is scaffolded, the evidence it worked, what to do if
%          it did not, and what to hand in.

\newcommand{\labsection}[3]{%
  \clearpage
  \phantomsection
  \addcontentsline{toc}{lab}{Lab #1 of \utplabtotal\space--- #2}%
  \markboth{Lab #1 of \utplabtotal\space--- #2}{}%
  \begin{tcolorbox}[enhanced,width=\linewidth,colback=UTPNavy,colframe=UTPNavy,
    boxrule=0pt,arc=2mm,left=5mm,right=5mm,top=5mm,bottom=5mm,
    before skip=0pt,after skip=1.2em]
    {\large\bfseries\color{UTPTeal!35!white} PRACTICAL MILESTONE\par}
    \vspace{0.25em}
    {\Huge\bfseries\color{white} Lab #1 of \utplabtotal\par}
    \vspace{0.20em}
    {\Large\bfseries\color{white} #2\par}
  \end{tcolorbox}
  \label{#3}%
}

\newcommand{\labflow}[3]{%
  \par\smallskip\noindent
  \textbf{Chapter connection.} #1\par
  \textbf{Goal.} #2\par
  \textbf{Explain.} #3\par\medskip}
\newcommand{\labsourcecard}[3]{%
  \par\medskip\noindent\textbf{#1 --- #2}\par
  \textbf{Inspect.} #3\par\smallskip}

\newtcolorbox{labquestion}[2]{enhanced,breakable,colback=UTPPaleOrange!55,colframe=UTPOrange,
  boxrule=0.8pt,arc=2mm,left=2.6mm,right=2.6mm,top=1.8mm,bottom=1.8mm,
  fonttitle=\bfseries\color{white},colbacktitle=UTPOrange,title={Question #1 --- #2}}
\newtcolorbox{labanswer}{enhanced,breakable,colback=white,colframe=UTPTeal,
  boxrule=0.8pt,arc=2mm,left=2.6mm,right=2.6mm,top=1.8mm,bottom=1.8mm,
  fonttitle=\bfseries\color{white},colbacktitle=UTPTeal,title=Worked answer}
\newcommand{\expectedoutput}{\par\smallskip\noindent\textbf{\color{UTPNavy}Reference output.}\nopagebreak\par\nopagebreak}
\utpnewcallout[boxrule=0.65pt,left=2.5mm,right=2.5mm,top=1.6mm,bottom=1.6mm,colbacktitle=UTPMuted]{checkanswers}{UTPPaleGray!72}{UTPMuted!62}{Quick-check answers}

% --------------------------------------------------------------------
%  13.  STAGE DIVIDER
% --------------------------------------------------------------------
%  A full-page divider used by the longer books to open a part:
%  \substage{3A}{Supervised Workflow}{one-line summary}.
%  Harmless to leave unused.

\newcommand{\substage}[3]{%
  \cleardoublepage
  \thispagestyle{empty}
  \phantomsection
  \addcontentsline{toc}{chapter}{Stage #1 --- #2}
  \begin{center}
    \vspace*{0.20\textheight}
    {\Large\bfseries\color{UTPTeal} Stage #1\par}
    \vspace{0.45em}
    {\Huge\bfseries\color{UTPNavy} #2\par}
    \vspace{0.65em}
    {\color{UTPTeal}\rule{0.46\linewidth}{1.5pt}\par}
    \vspace{1.1em}
    \begin{minipage}{0.76\linewidth}
      \centering\large\color{UTPMuted} #3
    \end{minipage}
  \end{center}
  \cleardoublepage
}

% --------------------------------------------------------------------
%  14.  FRONT MATTER
% --------------------------------------------------------------------
%  The opening every book shares: the cover, the contents in roman
%  numerals, the course map, then chapter one on an arabic page 1.
%
%      \begin{document}
%      \utpfrontmatter
%
%  What it prints comes from the configuration in section 1 --- the
%  cover text from \utpcoursetitle and \utpcourseoutline, the overview
%  chapter from \utpcoursemap --- and the contents page is set as
%  \utpcontentspage (section 5) describes. A book with no overview
%  chapter sets \def\utpcoursemap{} and the step is skipped.
%
%  Page anchors are off across the cover so hyperref does not warn
%  about the unnumbered pages, and back on before the contents links
%  are used.
%
%  The number of front-matter pages is recorded in the .aux file as
%  \utpfrontpages, which is how .shared/build.py turns the page numbers
%  in the contents into physical pages when it splits the book.

\newcommand{\utpfrontmatter}{%
  \hypersetup{pageanchor=false}%
  % ---------- Cover ----------
%  The title page, printed by \utpfrontmatter. Everything that changes
%  between books is a configuration macro from section 1 of utpnotes:
%  \utpcoursetitle, \utpcourseoutline, \utpoutlineheading, \utpauthor
%  and \utpauthorrole. The colours are the shared palette.

\begin{titlepage}
\thispagestyle{empty}
\begin{tikzpicture}[remember picture,overlay]
  \fill[UTPGold]
    (current page.north west) rectangle
    ([yshift=-0.12cm]current page.north east);
  \fill[UTPGold]
    (current page.south west) rectangle
    ([yshift=0.12cm]current page.south east);
\end{tikzpicture}

\centering
\includegraphics[
  height=2.85cm,
  clip
]{logo-UTP.png}\par

\vspace{0.50cm}
{\fontsize{29}{33}\selectfont\bfseries\color{UTPNavy}\utpcoursetitle\par}
\vspace{0.24cm}
{\Large\color{UTPMuted} \utpdocumentsubtitle\par}
\vspace{0.42cm}
{\color{UTPGold}\rule{0.26\linewidth}{1.3pt}\par}
\vspace{0.72cm}
\begin{tikzpicture}[
  core/.style={circle,draw=UTPGold,fill=UTPNavy,minimum size=1.95cm,
    align=center,font=\small\bfseries\color{white},line width=1pt},
  stage/.style={rounded corners=1.5mm,draw=UTPBlue!65,fill=UTPPaleBlue!42,
    text width=3.15cm,minimum height=1.42cm,align=center,
    font=\small\bfseries\color{UTPNavy},inner xsep=3mm,
    inner ysep=3.5mm,line width=0.65pt},
  cycle/.style={-{Stealth[length=2.1mm]},semithick,draw=UTPMuted,
    rounded corners=4mm,shorten <=7pt,shorten >=7pt}
]
\node[core] (guide) at (0,-2.55) {TA\\WORKFLOW};
\node[stage] (prepare) at (0,0) {SESSION\\PLANNING};
\node[stage] (explain) at (4.75,-2.55) {INSTRUCTIONAL\\SUPPORT};
\node[stage] (facilitate) at (0,-5.10) {STUDENT\\ENGAGEMENT};
\node[stage] (assess) at (-4.75,-2.55) {LEARNING\\ASSESSMENT};

\draw[cycle] (prepare.east) -- (explain.north);
\draw[cycle] (explain.south) -- (facilitate.east);
\draw[cycle] (facilitate.west) -- (assess.south);
\draw[cycle] (assess.north) -- (prepare.west);
\end{tikzpicture}

\vspace{0.72cm}
\begin{tcolorbox}[
  enhanced,width=0.82\linewidth,colback=white,colframe=UTPBlue!18,
  arc=1.5mm,boxrule=0.45pt,left=6mm,right=6mm,top=3.2mm,bottom=3.2mm,
  borderline west={1.6pt}{0pt}{UTPGold}]
\centering\normalsize\color{UTPInk}
{\large\bfseries\color{UTPBlue}\utpoutlineheading\par}
\vspace{0.3em}
\utpcourseoutline
\end{tcolorbox}

\vfill
\begin{minipage}{0.82\textwidth}
\begin{tabular}{@{}>{\RaggedRight\arraybackslash}m{0.42\linewidth}@{\hspace{0.10\linewidth}}>{\RaggedLeft\arraybackslash}m{0.48\linewidth}@{}}
{\normalsize\bfseries\color{UTPBlue}PREPARED BY\par\vspace{0.25em}
 \Large\bfseries\color{UTPNavy}\utpauthor\par\vspace{0.12em}
 \normalfont\small\color{UTPMuted}\utpauthorrole}
&
\includegraphics[height=2.10cm,keepaspectratio]{logo-DC.png}
\end{tabular}
\end{minipage}
\vspace{0.5cm}
\end{titlepage}%
  \cleardoublepage
  \pagenumbering{roman}%
  \setcounter{tocdepth}{0}%
  \utptableofcontents
  \newpage
  \hypersetup{pageanchor=true}%
  \ifx\utpcoursemap\empty\else\input{\utpcoursemap}\fi
  \cleardoublepage
  \ifdefined\ReadonlyShipoutCounter
    \immediate\write\csname @auxout\endcsname
      {\string\gdef\string\utpfrontpages{\the\ReadonlyShipoutCounter}}%
  \fi
  \pagenumbering{arabic}%
}

% --------------------------------------------------------------------
%  15.  LISTING STYLES
% --------------------------------------------------------------------
%  Two base looks, and one named style per language built on them, so
%  the colours and frame are the same in every book:
%
%      utpcode          the base for source code
%      utpterminal      the base for commands typed and output printed
%      utppython        Python
%      utpr             R
%      utpc             C
%      utpcpp           C++
%      utpbash          shell scripts and commands
%      utpcnumbered     C with line numbers, for a complete source file
%      utpbashnumbered  shell with line numbers, likewise
%
%      \begin{lstlisting}[style=utppython] ... \end{lstlisting}
%      \lstinputlisting[style=utpcnumbered]{all_experiments/Lab04/fcfs.c}
%
%  A book that sets most of its listings in one style makes it the
%  default once, after the input:  \lstset{style=utpcpp}

\lstdefinestyle{utpcode}{
  basicstyle=\ttfamily\scriptsize,
  keywordstyle=\bfseries\color{UTPBlue},
  commentstyle=\itshape\color{UTPMuted},
  stringstyle=\color{UTPGreen!70!black},
  backgroundcolor=\color{UTPPaleGray!72},
  frame=single,
  rulecolor=\color{UTPBlue!28},
  framesep=5pt,
  breaklines=true,
  breakatwhitespace=false,
  showstringspaces=false,
  keepspaces=true,
  columns=fullflexible,
  upquote=true,
  tabsize=4,
  aboveskip=0.7em,
  belowskip=0.8em
}
\lstdefinestyle{utpterminal}{
  basicstyle=\ttfamily\scriptsize,
  backgroundcolor=\color{UTPPaleGray!72},
  frame=single,
  rulecolor=\color{UTPMuted!35},
  framesep=5pt,
  breaklines=true,
  showstringspaces=false,
  keepspaces=true,
  columns=fullflexible,
  upquote=true,
  aboveskip=0.6em,
  belowskip=0.7em
}

\lstdefinestyle{utppython}{style=utpcode,language=Python}
\lstdefinestyle{utpr}{style=utpcode,language=R,tabsize=2}
\lstdefinestyle{utpc}{style=utpcode,language=C}
\lstdefinestyle{utpcpp}{style=utpcode,language=C++}
\lstdefinestyle{utpbash}{style=utpcode,language=bash,
  rulecolor=\color{UTPMuted!35},aboveskip=0.6em,belowskip=0.7em}
\lstdefinestyle{utpcnumbered}{style=utpc,numbers=left,numberstyle=\tiny\color{UTPMuted},
  numbersep=7pt,xleftmargin=1.7em,framexleftmargin=1.4em}
\lstdefinestyle{utpbashnumbered}{style=utpbash,numbers=left,numberstyle=\tiny\color{UTPMuted},
  numbersep=7pt,xleftmargin=1.7em,framexleftmargin=1.4em}

% Pseudocode. A book that sets algorithms asks for them once, after the
% input, and gets algorithm2e loaded and coloured to match:
%
%     \utpusealgorithms
%
% \algostyle then sets one algorithm environment a size smaller than the
% text around it.
\newcommand{\utpalgorithmstyle}{%
  \SetAlCapNameFnt{\bfseries\color{UTPNavy}}%
  \SetAlCapFnt{\bfseries\color{UTPNavy}}%
  \SetAlgoCaptionSeparator{.}%
  \SetKwInput{KwIn}{Input}%
  \SetKwInput{KwOut}{Output}%
  \SetNlSty{color}{UTPMuted}{}%
  \SetCommentSty{\itshape\color{UTPMuted}}%
  \SetKwSty{bfseries\color{UTPBlue}}%
  \SetArgSty{}%
  \providecommand{\algostyle}{\small\SetAlgoNlRelativeSize{-1}\DontPrintSemicolon\SetAlgoVlined}%
}
\newcommand{\utpusealgorithms}{%
  \usepackage[linesnumbered,ruled,vlined]{algorithm2e}%
  \utpalgorithmstyle
}

%
%      \utpcontentspage{\small}{\normalsize}{0.65em}{0.07em}{0.16em}
%
%      \begin{document}
%      \utpfrontmatter                     % cover, contents, map
%      \part{Orientation}
%      \input{chapters/01-where-to-begin}
%      \end{document}
%
%  Configuration goes BEFORE the \input; overrides go AFTER it.
%  Paths are relative, so compile from inside the course folder.
%
%  ------------------------------------------------------------------
%  NAMING
%  ------------------------------------------------------------------
%  One name per idea, used by every book. Anything a book can use is
%  either an unprefixed learning box (keyidea, pitfall, labproject,
%  ...) or carries the utp prefix (utpfigure, utptable, utpfig,
%  utpsteps, utppython, UTPNavy, utpcell, ...). Course prefixes such
%  as ml-, ds- or erp- are not used: a chapter moved from one book to
%  another compiles unchanged. .shared/README.md lists every name.
%
%  Adding a new course? Copy .shared/TEMPLATE-main.tex and follow
%  .shared/ADDING-A-COURSE.md. Every step is a numbered checklist.
%
%  ------------------------------------------------------------------
%  CONTENTS
%  ------------------------------------------------------------------
%    1  Course configuration     the knobs a book may set
%    2  Packages                 the shared package set and fonts
%    3  Colour palette           the one UTP palette
%    4  Page layout and tables   margins, lists, table styling
%    5  Table of contents        parts, chapters, labs, figures
%    6  Headings                 chapter badge, part pages, sections
%    7  Page furniture           running heads and footers
%    8  Learning boxes           key idea, pitfall, try it, and friends
%    9  Small commands           \term, \cmd, \chapterguide, utpsteps
%   10  Diagram node styles      the TikZ vocabulary every book shares
%   11  Figure frames            utpfigure, utpfig, source images
%   12  Lab sheets               the lab banner and the lab toolkit
%   13  Stage divider            the part-opening page in longer books
%   14  Front matter             \utpfrontmatter: cover, contents, map
%   15  Listing styles           code, terminal and pseudocode
% =====================================================================

% --------------------------------------------------------------------
%  1.  COURSE CONFIGURATION
% --------------------------------------------------------------------
%  The knobs below, plus \utpstyleparts in section 5 and the box titles
%  in section 8, are the whole configuration surface. A book sets what
%  it needs BEFORE the \input and inherits the rest.
%
%  \providecommand means "define unless the book already did", which is
%  what makes the \def-before-\input pattern work.
%
%  Running text
%    \utprunninghead   the left running head on every page.
%    \utpchapterlabel  the word in the chapter badge. CHAPTER by
%                      default; a book whose chapters are labs sets LAB.
%    \utplabtotal      the N in the "Lab 3 of N" banners.

\providecommand{\utprunninghead}{Course Notes}
\providecommand{\utpchapterlabel}{CHAPTER}
\providecommand{\utplabtotal}{10}

%  Cover, printed by section 14
%    \utpcoursetitle       the title on the title page.
%    \utpdocumentsubtitle  the line under the title; Teaching Assistant Guide by default.
%    \utpcourseoutline     the sentence or two below it; \\ breaks a line.
%    \utpoutlineheading    the small heading above that sentence.
%    \utpauthor            the name in the PREPARED BY block.
%    \utpauthorrole        the line under the name.

\providecommand{\utpcoursetitle}{\utprunninghead}
\providecommand{\utpdocumentsubtitle}{Teaching Assistant Guide}
\providecommand{\utpcourseoutline}{}
\providecommand{\utpoutlineheading}{COURSE OUTLINE}
\providecommand{\utpauthor}{Reeyan Afzal Khan}
\providecommand{\utpauthorrole}{Graduate Research Assistant}

%  Files and folders
%    \utpextragraphicspath  image folders to search before the standard
%                           ones, as brace groups: {{screenshots/}}.
%                           A course's own figures/ folder and the
%                           shared logos are always searched.
%    \utpcoursemap          the overview chapter \utpfrontmatter inputs
%                           after the contents. Empty means none.

\providecommand{\utpextragraphicspath}{}
\providecommand{\utpcoursemap}{chapters/course-map}

%  Contents entries for chapters (section 5 applies them). A book whose
%  chapters are labs says so once here rather than restating the three
%  tocloft settings after the input.
%    \utpchaptertocword      Chapter~ by default, Lab~ in one book.
%    \utpchaptertocindent    left indent of the entry.
%    \utpchaptertocnumwidth  room reserved for the word and number.

\providecommand{\utpchaptertocword}{Chapter~}
\providecommand{\utpchaptertocindent}{1.5em}
\providecommand{\utpchaptertocnumwidth}{6em}

%  Monospace face (section 2)
%    \utpmonofont   sourcecodepro by default. inconsolata is the
%                   alternative, for a book whose pages are mostly
%                   device transcripts.

\providecommand{\utpmonofont}{sourcecodepro}

% --------------------------------------------------------------------
%  2.  PACKAGES
% --------------------------------------------------------------------
%  The union of what the books need. A course that needs something
%  extra loads it in its own file, under the heading "packages only
%  this book needs". Pseudocode is the one optional package the shared
%  file knows about: \utpusealgorithms in section 15 loads and styles it.
%
%  The tikz and pgfplots library lists are the union across all books,
%  so a diagram or chart copied from one course compiles in another.

\usepackage[T1]{fontenc}
\IfFileExists{glyphtounicode.tex}{\input{glyphtounicode}\pdfgentounicode=1}{}
\usepackage[american]{babel}
\IfFileExists{sourcesanspro.sty}
  {\usepackage[default]{sourcesanspro}}
  {\usepackage[scaled]{helvet}\renewcommand{\familydefault}{\sfdefault}}
\DeclareFontShape{T1}{cmss}{b}{n}{<->ssub*cmss/bx/n}{}
\DeclareFontShape{TS1}{cmss}{b}{n}{<->ssub*cmss/bx/n}{}
% Monospace. Without this the tt default stays Computer Modern Typewriter,
% a serif face that clashes with the sans body text in every listing and
% every \cmd. Source Code Pro is the companion to Source Sans Pro; it is
% scaled down because mono digits otherwise sit large next to the sans.
%
% Inconsolata, chosen with \def\utpmonofont{inconsolata}, distinguishes
% 0/O and 1/l/I at a glance, which matters when a reader copies an
% interface name or an address out of a transcript. Only one mono package
% is ever loaded: a second would win the family but inherit the first's
% series names, and every \ttfamily would then resolve by substitution.
\def\utpmonoinconsolata{inconsolata}
\ifx\utpmonofont\utpmonoinconsolata
  \IfFileExists{inconsolata.sty}{\usepackage[varqu,varl,scaled=1.02]{inconsolata}}{}
\else
  \IfFileExists{sourcecodepro.sty}
    {\usepackage[scaled=0.88]{sourcecodepro}}
    {\usepackage[scaled=0.86]{beramono}}
  % Source Sans names its normal weight "regular", but the Computer Modern
  % fallbacks call it "m". Map one onto the other so a \ttfamily reached
  % while the regular series is active resolves instead of warning.
  \DeclareFontShape{T1}{cmtt}{regular}{n}{<->ssub*cmtt/m/n}{}
  \DeclareFontShape{TS1}{cmtt}{regular}{n}{<->ssub*cmtt/m/n}{}
  \DeclareFontShape{T1}{cmtt}{regular}{it}{<->ssub*cmtt/m/it}{}
\fi
\usepackage[margin=21mm,headheight=22pt]{geometry}
\usepackage{amsmath,amssymb,mathtools,bm}
\usepackage{microtype}
\usepackage{enumitem}
\usepackage{booktabs}
\usepackage{tabularx}
\usepackage{array}
\usepackage{ragged2e}
\usepackage{etoolbox}
\usepackage{titlesec}
\usepackage{tocloft}
\usepackage{fancyhdr}
\usepackage[table]{xcolor}
\usepackage{tcolorbox}
\tcbuselibrary{breakable,skins}
\usepackage[hidelinks]{hyperref}
\usepackage{bookmark}
% Document properties, so the PDF identifies itself in a reader's
% title bar, in a library catalogue, and in search results. Set from
% the same configuration macros the cover uses.
\hypersetup{
  pdftitle={\utpcoursetitle},
  pdfauthor={\utpauthor},
  pdfsubject={Universiti Teknologi PETRONAS course notes and laboratory guide},
  pdfkeywords={UTP, course notes, laboratory manual, \utpcoursetitle},
  pdfdisplaydoctitle=true,
  bookmarksnumbered=true,
  bookmarksopenlevel=1
}
\usepackage{listings}
\usepackage{graphicx}
\usepackage{caption}
\usepackage{tikz}
\usetikzlibrary{angles,arrows.meta,backgrounds,calc,chains,decorations.pathmorphing,decorations.pathreplacing,fit,matrix,patterns,positioning,quotes,shadows.blur,shapes.geometric,shapes.misc,shapes.symbols}
\usepackage{pgfplots}
% statistics gives boxplot prepared; groupplots gives multi-panel charts.
\usepgfplotslibrary{statistics}
\usepgfplotslibrary{groupplots}

% Images: any folders the book names first, then its own figures/
% folder, then the shared logos used by the cover.
\expandafter\graphicspath\expandafter{\utpextragraphicspath{figures/}{../.shared/logos/}}

% --------------------------------------------------------------------
%  3.  COLOUR PALETTE
% --------------------------------------------------------------------
%  One palette, defined once under UTP* names, and used by every book,
%  box, diagram and chart. Five accents (blue, teal, green, orange, red)
%  with a pale tint of each for box backgrounds, plus navy for headings,
%  gold for the cover, ink and muted for text, and a sixth accent,
%  purple, for charts and layered diagrams that need one more series.

\definecolor{UTPNavy}{HTML}{18324A}
\definecolor{UTPBlue}{HTML}{246B9E}
\definecolor{UTPTeal}{HTML}{138A80}
\definecolor{UTPGreen}{HTML}{2D7D46}
\definecolor{UTPOrange}{HTML}{C86B24}
\definecolor{UTPRed}{HTML}{B34843}
\definecolor{UTPPurple}{HTML}{6B4FA0}
\definecolor{UTPGold}{HTML}{D4A13E}
\definecolor{UTPInk}{HTML}{263238}
\definecolor{UTPMuted}{HTML}{60717C}
\definecolor{UTPPaleBlue}{HTML}{EEF6FB}
\definecolor{UTPPaleTeal}{HTML}{EEF8F6}
\definecolor{UTPPaleOrange}{HTML}{FFF5EA}
\definecolor{UTPPaleRed}{HTML}{FFF0EF}
\definecolor{UTPPalePurple}{HTML}{F0ECF9}
\definecolor{UTPPaleGray}{HTML}{F4F6F7}
\color{UTPInk}

% --------------------------------------------------------------------
%  4.  PAGE LAYOUT AND TABLES
% --------------------------------------------------------------------
%  Paragraph spacing, list indents, and the table look: alternating
%  row tint, top-aligned paragraph columns, ragged-right cells.
%
%  Two extra column types for tabularx tables:
%      D{width}   a fixed-width, ragged-right, vertically centred column
%      Y          a ragged-right column that fills the remaining width
%
%  Every table in every book has the same shape, the one Algorithms &
%  Data Structures set: a utptable frame around a full-width tabularx,
%  booktabs rules, a tinted header row of \utpth cells (centred, bold,
%  navy), and a first column in bold navy when it names the rows.
%  Cells are vertically centred, so use m{} rather than p{} columns.
%
%      \begin{utptable}
%      \begin{tabularx}{\linewidth}{@{}>{\bfseries\leavevmode\color{UTPNavy}}m{0.22\linewidth}X@{}}
%      \toprule
%      \rowcolor{UTPPaleBlue!92}
%      \utpth{Operation} & \utpth{What it does} \\
%      \midrule
%      Insertion & Add a new element. \\
%      \bottomrule
%      \end{tabularx}
%      \end{utptable}
%
%  A book whose tables carry dense text asks for \utpcompacttables after
%  the input: tabularx tables then set in \footnotesize with tighter
%  padding and vertically centred cells.

\setlength{\parindent}{0pt}
\setlength{\parskip}{0.55em}
\setlist[itemize]{leftmargin=1.45em,itemsep=0.22em,topsep=0.25em}
\setlist[enumerate]{leftmargin=1.75em,itemsep=0.28em,topsep=0.3em}
\setlist[description]{font=\bfseries\color{UTPNavy},itemsep=0.25em,topsep=0.3em}
\renewcommand{\arraystretch}{1.34}
\setlength{\tabcolsep}{6.5pt}
\arrayrulecolor{UTPBlue!45}
% Keep paragraph-style columns ragged and vertically centred. Tables use
% fixed m columns beside flexible X columns; giving X the same centring
% keeps a short cell level with a tall neighbour in a multi-line row.
\renewcommand{\tabularxcolumn}[1]{m{#1}}
\makeatletter
\g@addto@macro\@arrayparboxrestore{\RaggedRight\arraybackslash}
\makeatother
% A header cell: centred, bold and navy, on the tinted header row.
\newcommand{\utpth}[1]{\multicolumn{1}{c}{\bfseries\color{UTPNavy}#1}}
\AtBeginEnvironment{tabular}{\small\setlength{\extrarowheight}{0.6pt}\rowcolors{2}{UTPPaleBlue!72}{white}}
\AtBeginEnvironment{tabularx}{\small\setlength{\extrarowheight}{0.6pt}\rowcolors{2}{UTPPaleBlue!72}{white}}
\AtBeginEnvironment{figure}{\centering}
\newcolumntype{D}[1]{>{\raggedright\arraybackslash}m{#1}}
\newcolumntype{Y}{>{\raggedright\arraybackslash}X}
\newcommand{\utpcompacttables}{%
  \renewcommand{\tabularxcolumn}[1]{m{##1}}%
  \AtBeginEnvironment{tabularx}{%
    \footnotesize
    \setlength{\tabcolsep}{4pt}%
    \renewcommand{\arraystretch}{1.25}%
    \setlength{\extrarowheight}{0.6pt}%
  }%
}
\emergencystretch=1.5em
\hbadness=10000
\hfuzz=8pt

% Page-breaking discipline. Without these a paragraph can leave a single
% line stranded at the foot of one page or the head of the next, which is
% the most visible sign of unattended typesetting.
\widowpenalty=10000          % no first line of a paragraph left alone at a page foot
\clubpenalty=10000           % no last line of a paragraph left alone at a page head
\displaywidowpenalty=10000   % the same around displayed equations
\brokenpenalty=10000         % do not break a page after a hyphenated line
\predisplaypenalty=10000     % keep the line before a display with it
% Pages are built from boxes of many heights, so let the last line float
% up rather than stretching the gaps between boxes to force a flush foot.
\raggedbottom
% Prefer a slightly loose line to an overfull one, but do not hyphenate
% freely -- short sans lines look ragged when over-hyphenated.
\tolerance=1200
\hyphenpenalty=200

% Character protrusion and font expansion. Both are what stop a justified
% sans column from showing visible rivers down the page.
\microtypesetup{protrusion=true,expansion=true,tracking=true}

\pgfplotsset{
  compat=1.18,
  every axis/.append style={
    axis line style={draw=UTPMuted!85,line width=0.55pt},
    tick style={draw=UTPMuted!85},
    tick label style={font=\scriptsize\color{UTPMuted}},
    label style={font=\small\color{UTPNavy}},
    title style={font=\small\bfseries\color{UTPNavy}},
    legend style={font=\scriptsize,draw=UTPBlue!22,fill=white,rounded corners=1pt}
  },
  every axis plot/.append style={line width=1.15pt}
}
\tikzset{every picture/.append style={line cap=round,line join=round}}

% --------------------------------------------------------------------
%  5.  TABLE OF CONTENTS
% --------------------------------------------------------------------
%  Formatting for part, chapter and lab entries, and for a list of
%  figures if a book prints one. The lab level is a custom tocloft list
%  so lab banners appear in the contents.
%
%  Part styling is skipped when \utpstyleparts is 0, which is what a
%  book without \part divisions should set.
%
%  The contents page is meant to fill exactly one page, and how loosely
%  it is set depends on how many entries a book has. One line after the
%  input tunes it:
%
%      \utpcontentspage{<base size>}{<part size>}%
%                      {<space before a part>}{<before a chapter>}{<before a lab>}
%      \utpcontentspage{\normalsize}{\large}{1.15em}{0.38em}{0.30em}   % the default

\addto\captionsamerican{\renewcommand{\contentsname}{Learning Path}}
\newlistentry{lab}{toc}{-1}
\providecommand{\utpstyleparts}{1}
\ifnum\utpstyleparts=1
\renewcommand{\cftpartpresnum}{}
\renewcommand{\cftpartaftersnum}{}
\setlength{\cftpartnumwidth}{5.2em}
\setlength{\cftbeforepartskip}{0.38em}
\renewcommand{\cftpartfont}{\normalsize\bfseries\color{UTPNavy}}
\renewcommand{\cftpartpagefont}{\normalsize\bfseries\color{UTPNavy}}
\renewcommand{\cftpartleader}{\hfill}
\fi
\renewcommand{\cftchappresnum}{\utpchaptertocword}
\renewcommand{\cftchapaftersnum}{\quad}
\setlength{\cftchapindent}{\utpchaptertocindent}
\setlength{\cftchapnumwidth}{\utpchaptertocnumwidth}
\setlength{\cftbeforechapskip}{0pt}
\renewcommand{\cftchapfont}{\normalfont}
\renewcommand{\cftchappagefont}{\normalfont}
\renewcommand{\cftchapleader}{\cftdotfill{\cftdotsep}}
\setlength{\cftlabindent}{3em}
\setlength{\cftlabnumwidth}{0em}
\setlength{\cftbeforelabskip}{0.10em}
\renewcommand{\cftlabfont}{\bfseries\color{UTPOrange!90!black}}
\renewcommand{\cftlabpagefont}{\bfseries\color{UTPOrange!90!black}}
\renewcommand{\cftlableader}{\color{UTPOrange!65!UTPMuted}\cftdotfill{\cftlabdotsep}}

% A list of figures, for a book that prints one (\listoffigures).
\renewcommand{\cftfigfont}{\normalfont\small}
\renewcommand{\cftfigpagefont}{\normalfont\small}
\renewcommand{\cftfigpresnum}{Figure~}
\renewcommand{\cftfigaftersnum}{\quad}
\setlength{\cftfigindent}{1.5em}
\setlength{\cftfignumwidth}{5.6em}
\setlength{\cftbeforefigskip}{0pt}
\renewcommand{\cftfigleader}{\cftdotfill{\cftdotsep}}

% The contents page itself, printed by \utpfrontmatter (section 14).
\newcommand{\utpcontentspage}[5]{%
  \def\utptocsize{#1}\def\utptocpartsize{#2}%
  \def\utptocpartskip{#3}\def\utptocchapterskip{#4}\def\utptoclabskip{#5}}
\utpcontentspage{\normalsize}{\large}{1.15em}{0.38em}{0.30em}
\newcommand{\utptableofcontents}{%
  \begingroup
  \utptocsize
  \setlength{\parskip}{0pt}%
  \setlength{\cftbeforepartskip}{\utptocpartskip}%
  \setlength{\cftbeforechapskip}{\utptocchapterskip}%
  \setlength{\cftbeforelabskip}{\utptoclabskip}%
  \renewcommand{\cftpartfont}{\utptocpartsize\bfseries\color{UTPNavy}}%
  \renewcommand{\cftpartpagefont}{\utptocpartsize\bfseries\color{UTPNavy}}%
  \tableofcontents
  \endgroup
}

% --------------------------------------------------------------------
%  6.  HEADINGS
% --------------------------------------------------------------------
%  The chapter badge (a coloured block with a word and the number),
%  part pages, and section spacing.
%
%  The word in the badge is \utpchapterlabel; Data Communication sets
%  it to LAB because each of its chapters is a lab.

\titleformat{\chapter}[display]
  {\bfseries\color{UTPNavy}}
  {\filleft\colorbox{UTPBlue}{\parbox{2.55cm}{\centering\color{white}\Large
    \utpchapterlabel\\[-0.15em]\Huge\thechapter}}}
  {0.65em}{\Huge\raggedright}
  [\vspace{0.25em}{\color{UTPBlue}\titlerule[1.2pt]}]
\renewcommand{\thepart}{Stage~\arabic{part}}
\titleformat{\part}[display]
  {\centering\color{UTPNavy}\bfseries}
  {\Large\color{UTPTeal}\thepart}{0.45em}{\Huge}
  [\vspace{0.55em}\color{UTPTeal}\rule{0.45\linewidth}{1.5pt}]
\titleformat{\section}{\Large\bfseries\color{UTPNavy}}{\thesection}{0.55em}{}
\titleformat{\subsection}{\large\bfseries\color{UTPBlue}}{\thesubsection}{0.5em}{}
\titlespacing*{\chapter}{0pt}{-1.8em}{1.0em}
\titlespacing*{\part}{0pt}{0pt}{1em}
\titlespacing*{\section}{0pt}{1.2em}{0.35em}
\titlespacing*{\subsection}{0pt}{0.95em}{0.2em}

% --------------------------------------------------------------------
%  7.  PAGE FURNITURE
% --------------------------------------------------------------------
%  Running heads and footers. The left header is \utprunninghead;
%  the right one follows the current chapter.

\pagestyle{fancy}
\fancyhf{}
\lhead{\small\color{UTPMuted}\utprunninghead}
\rhead{\small\color{UTPMuted}\nouppercase{\leftmark}}
\cfoot{\small\color{UTPMuted}\thepage}
\renewcommand{\headrulewidth}{0.5pt}
\renewcommand{\headrule}{\hbox to\headwidth{\color{UTPPaleBlue}\leaders\hrule height \headrulewidth\hfill}}
\renewcommand{\chaptermark}[1]{\markboth{#1}{}}
\fancypagestyle{plain}{%
  \fancyhf{}%
  \cfoot{\small\color{UTPMuted}\thepage}%
  \renewcommand{\headrulewidth}{0pt}%
}

% --------------------------------------------------------------------
%  8.  LEARNING BOXES
% --------------------------------------------------------------------
%  The coloured callouts the chapters use. Every title is a macro, so
%  a book can rename a box without restating its geometry:
%
%      \def\utptitlepitfall{Technical note}   % before the input
%
%  To change how a box LOOKS in one book, use one of the named
%  variations at the end of this section, or \renewtcolorbox after the
%  input. To add a box only one book needs, declare it in that book
%  with \utpnewcallout or \utpnewquietbox.

% Box titles -- override any of these before the input line.
\providecommand{\utptitlekeyidea}{Key idea}
\providecommand{\utptitleworked}{Worked example}
\providecommand{\utptitlepitfall}{Common mistake}
\providecommand{\utptitletryit}{Try it yourself}
\providecommand{\utptitlequickcheck}{Check your understanding}
\providecommand{\utptitlerecap}{Chapter recap}
\providecommand{\utptitledeliverable}{Lab deliverable}
\providecommand{\utptitlelabmode}{How much is scaffolded?}
\providecommand{\utptitleintuition}{Plain idea}
\providecommand{\utptitlefirstread}{Core path}
\providecommand{\utptitleglance}{Chapter at a glance}
\providecommand{\utptitledebug}{Debugging task}
\providecommand{\utptitleoptional}{Beyond the core}
\providecommand{\utptitleverification}{Verification checklist}
\providecommand{\utptitletroubleshoot}{Troubleshooting}
\providecommand{\utptitlesourcebasis}{Source basis}

% Explaining an idea
\newtcolorbox{keyidea}{enhanced,breakable,colback=UTPPaleTeal,colframe=UTPTeal, boxrule=0.8pt,arc=2mm,left=2.5mm,right=2.5mm,top=1.8mm,bottom=1.8mm, fonttitle=\bfseries\color{UTPNavy},title=\utptitlekeyidea}
\newtcolorbox{workedexample}{enhanced,breakable,colback=UTPPaleBlue,colframe=UTPBlue, boxrule=0.8pt,arc=2mm,left=2.5mm,right=2.5mm,top=1.8mm,bottom=1.8mm, fonttitle=\bfseries\color{UTPNavy},title=\utptitleworked}
\newtcolorbox{pitfall}{enhanced,breakable,colback=UTPPaleRed,colframe=UTPRed, boxrule=0.8pt,arc=2mm,left=2.5mm,right=2.5mm,top=1.8mm,bottom=1.8mm, fonttitle=\bfseries\color{UTPNavy},title=\utptitlepitfall}
% Asking the reader to do something
\newtcolorbox{tryit}{enhanced,breakable,colback=white,colframe=UTPTeal, boxrule=0.9pt,arc=2mm,left=2.5mm,right=2.5mm,top=1.8mm,bottom=1.8mm, fonttitle=\bfseries\color{white},colbacktitle=UTPTeal,title=\utptitletryit}
\newtcolorbox{quickcheck}{enhanced,breakable,colback=white,colframe=UTPOrange, boxrule=0.9pt,arc=2mm,left=2.5mm,right=2.5mm,top=1.8mm,bottom=1.8mm, fonttitle=\bfseries\color{white},colbacktitle=UTPOrange,title=\utptitlequickcheck}
\newtcolorbox{chapterrecap}{enhanced,breakable,colback=UTPPaleTeal,colframe=UTPGreen, boxrule=0.9pt,arc=2mm,left=2.5mm,right=2.5mm,top=1.8mm,bottom=1.8mm, fonttitle=\bfseries\color{white},colbacktitle=UTPGreen,title=\utptitlerecap}
% Lab structure
\newtcolorbox{labdeliverable}{enhanced,breakable,colback=white,colframe=UTPGreen, boxrule=0.9pt,arc=2mm,left=2.6mm,right=2.6mm,top=1.9mm,bottom=1.9mm, fonttitle=\bfseries\color{white},colbacktitle=UTPGreen,title=\utptitledeliverable}
\newtcolorbox{labmode}{enhanced,breakable,colback=UTPPaleBlue!52,colframe=UTPBlue, boxrule=0.75pt,arc=2mm,left=2.6mm,right=2.6mm,top=1.7mm,bottom=1.7mm, fonttitle=\bfseries\color{white},colbacktitle=UTPBlue,title=\utptitlelabmode}
% Softer, lighter-weight asides
\newtcolorbox{intuition}{enhanced,breakable,colback=white,colframe=UTPOrange!42, boxrule=0.5pt,arc=2mm,left=2.5mm,right=2.5mm,top=1.5mm,bottom=1.5mm, fonttitle=\bfseries\color{UTPOrange!85!black},title=\utptitleintuition}
\newtcolorbox{firstread}{enhanced,breakable,colback=white,colframe=UTPBlue!42, boxrule=0.5pt,arc=2mm,left=2.5mm,right=2.5mm,top=1.5mm,bottom=1.5mm, fonttitle=\bfseries\color{UTPBlue}, title=\utptitlefirstread}
\newtcolorbox{debugtask}{enhanced,breakable,colback=UTPPaleRed!58,colframe=UTPRed, boxrule=0.85pt,arc=2mm,left=2.6mm,right=2.6mm,top=1.9mm,bottom=1.9mm, fonttitle=\bfseries\color{white},colbacktitle=UTPRed,title=\utptitledebug}
\newtcolorbox{analogy}{enhanced,breakable,colback=white,colframe=UTPOrange!38, boxrule=0.5pt,arc=2mm,left=2.5mm,right=2.5mm,top=1.5mm,bottom=1.5mm, fonttitle=\bfseries\color{UTPOrange!85!black},title=Analogy}
\newtcolorbox{coursechoice}{enhanced,breakable,colback=white,colframe=UTPNavy, boxrule=0.8pt,arc=2mm,left=2.8mm,right=2.8mm,top=2mm,bottom=2mm, fonttitle=\bfseries\color{white},colbacktitle=UTPNavy,title=Course choice}
\newtcolorbox{decisionmemo}{enhanced,breakable,colback=UTPPaleOrange!55,colframe=UTPOrange, boxrule=0.85pt,arc=2mm,left=2.6mm,right=2.6mm,top=1.9mm,bottom=1.9mm, fonttitle=\bfseries\color{white},colbacktitle=UTPOrange,title=Decision memo}
\newtcolorbox{libnote}{enhanced,breakable,colback=white,colframe=UTPMuted!48, boxrule=0.5pt,arc=2mm,left=2.5mm,right=2.5mm,top=1.5mm,bottom=1.5mm, fonttitle=\bfseries\color{UTPNavy},title=Reference note}
\newtcolorbox{miniproject}{enhanced,breakable,colback=UTPPaleTeal!52,colframe=UTPGreen, boxrule=0.9pt,arc=2mm,left=2.6mm,right=2.6mm,top=1.9mm,bottom=1.9mm, fonttitle=\bfseries\color{white},colbacktitle=UTPGreen,title=Mini-project}
% Where the material came from: the lecturer's slides, manual or tutorial
% that a chapter is built on.
\newtcolorbox{sourcebasis}{enhanced,breakable,colback=white,colframe=UTPMuted!55, boxrule=0.5pt,arc=2mm,left=2.5mm,right=2.5mm,top=1.5mm,bottom=1.5mm, fonttitle=\bfseries\color{UTPMuted},title=\utptitlesourcebasis}
% Boxes that take a title argument
\newtcolorbox{modelcard}[1]{enhanced,breakable,colback=white,colframe=UTPBlue, boxrule=0.9pt,arc=2mm,left=2.5mm,right=2.5mm,top=1.8mm,bottom=1.8mm, fonttitle=\bfseries\color{white},colbacktitle=UTPBlue,title=#1}
\newtcolorbox{optionaltopic}{enhanced,breakable,colback=white,colframe=UTPMuted!58, boxrule=0.5pt,arc=2mm,left=2.5mm,right=2.5mm,top=1.5mm,bottom=1.5mm, fonttitle=\bfseries\color{UTPMuted},title=\utptitleoptional}
\newtcolorbox{secondpass}{enhanced,breakable,colback=UTPPaleGray!82,colframe=UTPMuted, boxrule=0.75pt,arc=2mm,left=2.6mm,right=2.6mm,top=1.9mm,bottom=1.9mm, fonttitle=\bfseries\color{white},colbacktitle=UTPMuted,title=Second pass}
\newtcolorbox{labproject}[1]{enhanced,breakable,colback=white,colframe=UTPBlue, boxrule=0.9pt,arc=2mm,left=2.6mm,right=2.6mm,top=1.9mm,bottom=1.9mm, fonttitle=\bfseries\color{white},colbacktitle=UTPBlue,title={#1}}
% Evidence that a lab worked, and what to do when it did not
\newtcolorbox{verification}{enhanced,breakable,colback=UTPPaleTeal!65,colframe=UTPGreen, boxrule=0.8pt,arc=2mm,left=2.6mm,right=2.6mm,top=1.8mm,bottom=1.8mm, fonttitle=\bfseries\color{white},colbacktitle=UTPGreen,title=\utptitleverification}
\newtcolorbox{troubleshoot}{enhanced,breakable,colback=UTPPaleOrange!70,colframe=UTPOrange, boxrule=0.8pt,arc=2mm,left=2.6mm,right=2.6mm,top=1.8mm,bottom=1.8mm, fonttitle=\bfseries\color{white},colbacktitle=UTPOrange,title=\utptitletroubleshoot}

% Named variations a book asks for after the input. Each restyles boxes
% that already exist, so the chapters do not change.
%
%   \utpheavycallouts          a solid colour bar behind the titles of
%                              keyidea, workedexample, intuition and
%                              firstread, with a firmer frame.
%   \utpbreakableframes        let utpfigure and utptable split over a
%                              page break, in a lighter blue (section 11).
%   \utplabprojectcolour{...}  the labproject box in another colour, with
%                              a little more room inside.
\newcommand{\utplabprojectcolour}[1]{%
  \renewtcolorbox{labproject}[1]{enhanced,breakable,colback=white,colframe=#1, boxrule=0.9pt,arc=2mm,left=2.7mm,right=2.7mm,top=2mm,bottom=2mm, fonttitle=\bfseries\color{white},colbacktitle=#1,title={##1}}%
}
\newcommand{\utpheavycallouts}{%
  \renewtcolorbox{keyidea}{enhanced,breakable,colback=UTPPaleTeal,colframe=UTPTeal, boxrule=0.85pt,arc=2mm,left=2.5mm,right=2.5mm,top=1.8mm,bottom=1.8mm, fonttitle=\bfseries\color{white},colbacktitle=UTPTeal,title=\utptitlekeyidea}%
  \renewtcolorbox{workedexample}{enhanced,breakable,colback=UTPPaleBlue,colframe=UTPBlue, boxrule=0.8pt,arc=2mm,left=2.5mm,right=2.5mm,top=1.8mm,bottom=1.8mm, fonttitle=\bfseries\color{white},colbacktitle=UTPBlue,title=\utptitleworked}%
  \renewtcolorbox{intuition}{enhanced,breakable,colback=UTPPaleOrange,colframe=UTPOrange, boxrule=0.8pt,arc=2mm,left=2.5mm,right=2.5mm,top=1.8mm,bottom=1.8mm, fonttitle=\bfseries\color{white},colbacktitle=UTPOrange,title=\utptitleintuition}%
  \renewtcolorbox{firstread}{enhanced,breakable,colback=white,colframe=UTPBlue, boxrule=0.9pt,arc=2mm,left=2.5mm,right=2.5mm,top=1.8mm,bottom=1.8mm, fonttitle=\bfseries\color{white},colbacktitle=UTPBlue,title=\utptitlefirstread}%
}

% Declaring a box a single book needs. Both constructors take the house
% geometry and let a trailing key list override any of it, so a new box
% is one line and still matches the others:
%
%   \utpnewcallout{odoobox}{white}{UTPBlue}{Odoo 19 Community}
%   \utpnewcallout[boxrule=0.9pt]{name}{back}{frame}{Title}
%   \utpnewquietbox{referencebox}{UTPMuted!58}{UTPMuted}{Reference note}
%
% A callout is a tinted box under a solid title bar, for something the
% reader has to act on. A quiet box is a light frame with coloured
% title text, for an aside they can read past. If a second book ever
% wants the same box, move its line into this section.
\newcommand{\utpnewcallout}[5][]{%
  \newtcolorbox{#2}{enhanced,breakable,colback=#3,colframe=#4, boxrule=0.8pt,arc=2mm,left=2.6mm,right=2.6mm,top=1.8mm,bottom=1.8mm, fonttitle=\bfseries\color{white},colbacktitle=#4,title={#5},#1}%
}
\newcommand{\utpnewquietbox}[5][]{%
  \newtcolorbox{#2}{enhanced,breakable,colback=white,colframe=#3, boxrule=0.5pt,arc=2mm,left=2.5mm,right=2.5mm,top=1.5mm,bottom=1.5mm, fonttitle=\bfseries\color{#4},title={#5},#1}%
}

% --------------------------------------------------------------------
%  9.  SMALL COMMANDS
% --------------------------------------------------------------------
%  \term for a defined term, \cmd for inline code, \chapref and
%  \secref for cross-references, and the utpsteps numbered-step list.
%
%      \term{overfitting}          a term being defined
%      \cmd{git status}            inline code, in any language
%      \chapref{ch:trees}          "Chapter 12"
%      \secref{sec:pruning}        "Section 12.3"
%      \chapterguide{prior}{outcomes}   the box that opens a chapter
%      \begin{utpsteps} ... \end{utpsteps}  a numbered Step 1/Step 2 list
%
%  The second group is the shorthand the chapters use in running text:
%
%      \figref{fig:tree}           "Figure 4.2"
%      \chapsref{ch:a}{ch:b}       "Chapters 3--5"
%      \lib{pandas}                a library name
%      \cxg{x}, \cxo{x}, \cxr{x}   maths coloured green/orange/red
%      \given                      the bar in P(a \given b)
%      \coretag, \latertag         read-now / come-back-to markers
%      \kagglesource{url}{name}    the dataset a lab is built on

\newcommand{\term}[1]{\textbf{\color{UTPNavy}#1}}
\newcommand{\cmd}[1]{\texttt{#1}}
\newcommand{\chapref}[1]{Chapter~\ref{#1}}
\newcommand{\secref}[1]{Section~\ref{#1}}
\newcommand{\chapterguide}[2]{%
  \begin{tcolorbox}[enhanced,breakable,colback=white,colframe=UTPNavy,boxrule=0.8pt,
    arc=2mm,left=2.5mm,right=2.5mm,top=2mm,bottom=2mm,
    title=\utptitleglance,fonttitle=\bfseries\color{white},colbacktitle=UTPNavy]
    \textbf{Before you start:} #1\par
    \textbf{By the end, you can:} #2
  \end{tcolorbox}}
\newenvironment{utpsteps}
  {\begin{enumerate}[label=\textcolor{UTPBlue}{\bfseries Step~\arabic*:},leftmargin=4.5em,itemsep=0.38em]}
  {\end{enumerate}}

\newcommand{\figref}[1]{Figure~\ref{#1}}
\newcommand{\chapsref}[2]{Chapters~\ref{#1}--\ref{#2}}
\newcommand{\lib}[1]{\textbf{\color{UTPBlue}\texttt{#1}}}
\newcommand{\cxg}[1]{{\color{UTPGreen}\bfseries $#1$}}
\newcommand{\cxo}[1]{{\color{UTPOrange}\bfseries $#1$}}
\newcommand{\cxr}[1]{{\color{UTPRed}\bfseries $#1$}}
\newcommand{\given}{\mathrel{|}}
\DeclareMathOperator*{\argmin}{arg\,min}
\DeclareMathOperator*{\argmax}{arg\,max}
\newcommand{\coretag}{\par\smallskip\noindent\colorbox{UTPPaleBlue}{\strut\small\bfseries\color{UTPBlue}CORE NOW}\par\smallskip}
\newcommand{\latertag}{\par\smallskip\noindent\colorbox{UTPPaleGray}{\strut\small\bfseries\color{UTPMuted}RETURN LATER}\par\smallskip}
\newcommand{\kagglesource}[2]{%
  \par\smallskip\noindent\textbf{Kaggle source:} \href{#1}{#2}\par
}

% --------------------------------------------------------------------
%  10.  DIAGRAM NODE STYLES
% --------------------------------------------------------------------
%  The vocabulary the hand-drawn TikZ diagrams are built from: a cell
%  in an array, the one cell under the cursor, a circular graph node,
%  the arrow between two of them. Every book has the whole set:
%
%      utpcell, utphot, utpdone, utpdim, utpwall        array cells
%      utpnode, utpnodehot, utpnodedone, utpnodedim     graph nodes
%      utparrow, utpptr, utpdash, utpedge               connectors
%      utplabel, utptag, utpbrace                       annotation
%      utpstep, utpstage, utppanel, utpbox              boxes
%
%  A cell is 1.0cm by 0.8cm. A book whose diagrams were drawn to a
%  different grid resizes the cells once, after the input:
%
%      \utpdiagramsize{1.05cm}{0.85cm}
%
%  Larger vocabularies for one kind of drawing live in .shared/diagrams/
%  and are loaded by name, after the input:
%
%      \utpusediagrams{network}         devices, links, packet headers,
%                                       sequence diagrams, state machines
%      \utpusediagrams{process}         business-process flows, documents,
%                                       swimlanes, status chips
%      \utpusediagrams{data-objects}    vectors, lists and data frames
%                                       drawn as cells; flowcharts
%      \utpusediagrams{network,process} several at once

\newcommand{\utpdiagramsize}[2]{%
  \tikzset{
    utpcell/.style={draw=UTPBlue,fill=UTPPaleBlue,minimum width=#1,minimum height=#2,
                    font=\small\bfseries\color{UTPNavy},inner sep=1pt},
    utphot/.style={draw=UTPOrange,fill=UTPPaleOrange,minimum width=#1,minimum height=#2,
                   font=\small\bfseries\color{UTPNavy},line width=1pt,inner sep=1pt},
    utpdone/.style={draw=UTPGreen,fill=UTPPaleTeal,minimum width=#1,minimum height=#2,
                    font=\small\bfseries\color{UTPNavy},inner sep=1pt},
    utpdim/.style={draw=UTPMuted!45,fill=UTPPaleGray,minimum width=#1,minimum height=#2,
                   font=\small\color{UTPMuted},inner sep=1pt},
    utpwall/.style={draw=UTPNavy,fill=UTPNavy!82,minimum width=#1,minimum height=#2,
                    inner sep=1pt},
    utpnode/.style={circle,draw=UTPBlue,fill=UTPPaleBlue,minimum size=8mm,
                    font=\small\bfseries\color{UTPNavy},inner sep=0.5pt},
    utpnodehot/.style={circle,draw=UTPOrange,fill=UTPPaleOrange,minimum size=8mm,line width=1pt,
                       font=\small\bfseries\color{UTPNavy},inner sep=0.5pt},
    utpnodedone/.style={circle,draw=UTPGreen,fill=UTPPaleTeal,minimum size=8mm,
                        font=\small\bfseries\color{UTPNavy},inner sep=0.5pt},
    utpnodedim/.style={circle,draw=UTPMuted!45,fill=UTPPaleGray,minimum size=8mm,
                       font=\small\color{UTPMuted},inner sep=0.5pt},
    utparrow/.style={-{Stealth[length=2.2mm]},thick,draw=UTPTeal},
    utpptr/.style={-{Stealth[length=2mm]},semithick,draw=UTPBlue},
    utpdash/.style={-{Stealth[length=2mm]},semithick,draw=UTPRed,dashed},
    utpedge/.style={thick,draw=UTPTeal},
    utplabel/.style={font=\scriptsize\color{UTPMuted},inner sep=1pt},
    utptag/.style={font=\scriptsize\bfseries\color{UTPOrange},inner sep=1pt},
    utpbrace/.style={decorate,decoration={brace,amplitude=4pt,raise=2pt},draw=UTPMuted,semithick},
    utpstep/.style={draw=UTPNavy,fill=white,rounded corners=1.5mm,minimum height=0.75cm,
                    font=\scriptsize\color{UTPNavy},align=center,inner sep=2pt},
    utpstage/.style={draw=UTPNavy,fill=white,rounded corners=1.5mm,minimum height=0.75cm,
                     font=\scriptsize\color{UTPNavy},align=center,inner sep=2pt},
    utppanel/.style={draw=UTPBlue!30,fill=white,rounded corners=1.5mm,inner sep=2.5mm},
    utpbox/.style={draw=UTPBlue,fill=UTPPaleBlue,rounded corners=2mm,align=center,
                   font=\small\color{UTPNavy},inner sep=2.2mm}
  }%
}
\utpdiagramsize{1.0cm}{0.8cm}

% \@for, not \foreach: \foreach runs each pass in a group, which would
% throw away the styles a library declares.
\makeatletter
\newcommand{\utpusediagrams}[1]{%
  \@for\utp@library:=#1\do{\input{../.shared/diagrams/\utp@library}}%
}
\makeatother

% --------------------------------------------------------------------
%  11.  FIGURE FRAMES
% --------------------------------------------------------------------
%  Three ways to put a drawing, chart or screenshot on the page:
%
%      \begin{utpfigure} ... \end{utpfigure}
%          a centred frame with no caption, for a diagram the text
%          talks about directly. An optional title puts a bar on top:
%          \begin{utpfigure}[Where the VLAN tag is added] ...
%
%      \begin{utptable} ... \end{utptable}
%          the same frame in teal, around a tabular or tabularx.
%
%      \begin{utpfig}[short entry]{The caption.}{fig:label}
%        \begin{tikzpicture} ... \end{tikzpicture}
%      \end{utpfig}
%          a frame with a numbered caption that \figref can point to.
%          The optional first argument is the short form for a list of
%          figures. A page never breaks between the drawing and its
%          caption.
%
%  A screenshot or slide lifted from the lecturer's material goes in a
%  titled frame, so the reader can tell it from a drawing made for the
%  notes; the optional width defaults to 0.9 and 0.78 of the line:
%
%      \sourceimage[0.8\linewidth]{Static-routing topology}{static-topology.png}
%      \compactsourceimage{The same, kept to a smaller height}{static-topology.png}
%
%  \fignote{...} adds a short muted line under a drawing, inside its frame.

\newtcolorbox{utpfigure}[1][]{enhanced,width=\linewidth,colback=white,colframe=UTPBlue!28, boxrule=0.7pt,arc=2mm, left=2.2mm,right=2.2mm,top=2.2mm,bottom=1.8mm, before skip=0.75em,after skip=0.85em, before upper={\centering}, halign=center, fonttitle=\bfseries\color{white},colbacktitle=UTPBlue,title={#1}}
\newtcolorbox{utptable}{enhanced,width=\linewidth,colback=white,colframe=UTPTeal!32, boxrule=0.7pt,arc=2mm, left=2.2mm,right=2.2mm,top=1.8mm,bottom=1.8mm, before skip=0.75em,after skip=0.85em, halign=center}
\newcommand{\utpbreakableframes}{%
  \renewtcolorbox{utpfigure}[1][]{enhanced,breakable,colback=white,colframe=UTPBlue!35, boxrule=0.5pt,arc=2mm,left=2.5mm,right=2.5mm,top=2.2mm,bottom=2.2mm, halign=center, fonttitle=\bfseries\color{white},colbacktitle=UTPBlue,title={##1}}%
  \renewtcolorbox{utptable}{enhanced,breakable,colback=white,colframe=UTPBlue!28, boxrule=0.45pt,arc=2mm,left=2.2mm,right=2.2mm,top=1.8mm,bottom=1.8mm}%
}

\captionsetup{
  hypcap=false,
  font={small},
  labelfont={bf,color=UTPNavy},
  textfont={color=UTPInk},
  labelsep=period,
  justification=centering,
  skip=4pt
}

\newenvironment{utpfig}[3][\relax]
  {\def\utpfigshort{#1}\def\utpfigcap{#2}\def\utpfiglab{#3}%
   \par\addvspace{0.55em}%
   \begin{tcolorbox}[enhanced,breakable,colback=white,colframe=UTPBlue!32,
     boxrule=0.5pt,arc=2mm,left=2.5mm,right=2.5mm,top=2.6mm,bottom=2.2mm]
   \centering}
  {\par\nopagebreak\vspace{0.45em}%
   % The caption is boxed: \captionof starts with a write and a zero
   % rule, and the glue after them would otherwise be a place where the
   % page could break, leaving the caption alone at the top of the next.
   % Inside the box the caption is told it sits below the drawing, which
   % it can no longer see, so the gap stays above it.
   \nointerlineskip
   \vbox{\setlength{\parskip}{0pt}\captionsetup{position=below}%
   \ifx\utpfigshort\relax
     \captionof{figure}{\utpfigcap}%
   \else
     \captionof{figure}[\utpfigshort]{\utpfigcap}%
   \fi
   \label{\utpfiglab}}%
   \end{tcolorbox}%
   \par\addvspace{0.55em}}

\newcommand{\sourceimage}[3][0.9\linewidth]{%
  \begin{utpfigure}[Source visual --- #2]%
    \includegraphics[width=#1,height=0.56\textheight,keepaspectratio]{#3}%
  \end{utpfigure}}
\newcommand{\compactsourceimage}[3][0.78\linewidth]{%
  \begin{utpfigure}[Source visual --- #2]%
    \includegraphics[width=#1,height=0.42\textheight,keepaspectratio]{#3}%
  \end{utpfigure}}

\newcommand{\fignote}[1]{\par\smallskip{\scriptsize\color{UTPMuted}#1\par}}

% --------------------------------------------------------------------
%  12.  LAB SHEETS
% --------------------------------------------------------------------
%  Every lab sheet opens with the banner, a full-width page break that
%  also puts the lab in the contents:
%
%      \labsection{3}{Scheduling Policies}{sec:lab03}
%
%  The "of N" comes from \utplabtotal, so a book states its lab count
%  once in its configuration rather than in every lab file.
%
%  The body of the sheet is a labproject box (section 8), and inside it
%  the same toolkit serves every course:
%
%      \labflow{chapter link}{goal}{what to explain}
%          the three-line brief under the lab title.
%      \labsourcecard{file}{what it does}{what to inspect}
%          the heading over a listing of one supplied source file.
%      \begin{labquestion}{3}{Short title} ... \end{labquestion}
%      \begin{labanswer} ... \expectedoutput ... \end{labanswer}
%          a question as the student meets it, then the worked code, the
%          output it should print, and what that output means.
%      \begin{checkanswers} ... \end{checkanswers}
%          answers to a quickcheck, set apart so they can be skipped.
%      labmode, verification, troubleshoot, labdeliverable (section 8)
%          how much is scaffolded, the evidence it worked, what to do if
%          it did not, and what to hand in.

\newcommand{\labsection}[3]{%
  \clearpage
  \phantomsection
  \addcontentsline{toc}{lab}{Lab #1 of \utplabtotal\space--- #2}%
  \markboth{Lab #1 of \utplabtotal\space--- #2}{}%
  \begin{tcolorbox}[enhanced,width=\linewidth,colback=UTPNavy,colframe=UTPNavy,
    boxrule=0pt,arc=2mm,left=5mm,right=5mm,top=5mm,bottom=5mm,
    before skip=0pt,after skip=1.2em]
    {\large\bfseries\color{UTPTeal!35!white} PRACTICAL MILESTONE\par}
    \vspace{0.25em}
    {\Huge\bfseries\color{white} Lab #1 of \utplabtotal\par}
    \vspace{0.20em}
    {\Large\bfseries\color{white} #2\par}
  \end{tcolorbox}
  \label{#3}%
}

\newcommand{\labflow}[3]{%
  \par\smallskip\noindent
  \textbf{Chapter connection.} #1\par
  \textbf{Goal.} #2\par
  \textbf{Explain.} #3\par\medskip}
\newcommand{\labsourcecard}[3]{%
  \par\medskip\noindent\textbf{#1 --- #2}\par
  \textbf{Inspect.} #3\par\smallskip}

\newtcolorbox{labquestion}[2]{enhanced,breakable,colback=UTPPaleOrange!55,colframe=UTPOrange,
  boxrule=0.8pt,arc=2mm,left=2.6mm,right=2.6mm,top=1.8mm,bottom=1.8mm,
  fonttitle=\bfseries\color{white},colbacktitle=UTPOrange,title={Question #1 --- #2}}
\newtcolorbox{labanswer}{enhanced,breakable,colback=white,colframe=UTPTeal,
  boxrule=0.8pt,arc=2mm,left=2.6mm,right=2.6mm,top=1.8mm,bottom=1.8mm,
  fonttitle=\bfseries\color{white},colbacktitle=UTPTeal,title=Worked answer}
\newcommand{\expectedoutput}{\par\smallskip\noindent\textbf{\color{UTPNavy}Reference output.}\nopagebreak\par\nopagebreak}
\utpnewcallout[boxrule=0.65pt,left=2.5mm,right=2.5mm,top=1.6mm,bottom=1.6mm,colbacktitle=UTPMuted]{checkanswers}{UTPPaleGray!72}{UTPMuted!62}{Quick-check answers}

% --------------------------------------------------------------------
%  13.  STAGE DIVIDER
% --------------------------------------------------------------------
%  A full-page divider used by the longer books to open a part:
%  \substage{3A}{Supervised Workflow}{one-line summary}.
%  Harmless to leave unused.

\newcommand{\substage}[3]{%
  \cleardoublepage
  \thispagestyle{empty}
  \phantomsection
  \addcontentsline{toc}{chapter}{Stage #1 --- #2}
  \begin{center}
    \vspace*{0.20\textheight}
    {\Large\bfseries\color{UTPTeal} Stage #1\par}
    \vspace{0.45em}
    {\Huge\bfseries\color{UTPNavy} #2\par}
    \vspace{0.65em}
    {\color{UTPTeal}\rule{0.46\linewidth}{1.5pt}\par}
    \vspace{1.1em}
    \begin{minipage}{0.76\linewidth}
      \centering\large\color{UTPMuted} #3
    \end{minipage}
  \end{center}
  \cleardoublepage
}

% --------------------------------------------------------------------
%  14.  FRONT MATTER
% --------------------------------------------------------------------
%  The opening every book shares: the cover, the contents in roman
%  numerals, the course map, then chapter one on an arabic page 1.
%
%      \begin{document}
%      \utpfrontmatter
%
%  What it prints comes from the configuration in section 1 --- the
%  cover text from \utpcoursetitle and \utpcourseoutline, the overview
%  chapter from \utpcoursemap --- and the contents page is set as
%  \utpcontentspage (section 5) describes. A book with no overview
%  chapter sets \def\utpcoursemap{} and the step is skipped.
%
%  Page anchors are off across the cover so hyperref does not warn
%  about the unnumbered pages, and back on before the contents links
%  are used.
%
%  The number of front-matter pages is recorded in the .aux file as
%  \utpfrontpages, which is how .shared/build.py turns the page numbers
%  in the contents into physical pages when it splits the book.

\newcommand{\utpfrontmatter}{%
  \hypersetup{pageanchor=false}%
  % ---------- Cover ----------
%  The title page, printed by \utpfrontmatter. Everything that changes
%  between books is a configuration macro from section 1 of utpnotes:
%  \utpcoursetitle, \utpcourseoutline, \utpoutlineheading, \utpauthor
%  and \utpauthorrole. The colours are the shared palette.

\begin{titlepage}
\thispagestyle{empty}
\begin{tikzpicture}[remember picture,overlay]
  \fill[UTPGold]
    (current page.north west) rectangle
    ([yshift=-0.12cm]current page.north east);
  \fill[UTPGold]
    (current page.south west) rectangle
    ([yshift=0.12cm]current page.south east);
\end{tikzpicture}

\centering
\includegraphics[
  height=2.85cm,
  clip
]{logo-UTP.png}\par

\vspace{0.50cm}
{\fontsize{29}{33}\selectfont\bfseries\color{UTPNavy}\utpcoursetitle\par}
\vspace{0.24cm}
{\Large\color{UTPMuted} \utpdocumentsubtitle\par}
\vspace{0.42cm}
{\color{UTPGold}\rule{0.26\linewidth}{1.3pt}\par}
\vspace{0.72cm}
\begin{tikzpicture}[
  core/.style={circle,draw=UTPGold,fill=UTPNavy,minimum size=1.95cm,
    align=center,font=\small\bfseries\color{white},line width=1pt},
  stage/.style={rounded corners=1.5mm,draw=UTPBlue!65,fill=UTPPaleBlue!42,
    text width=3.15cm,minimum height=1.42cm,align=center,
    font=\small\bfseries\color{UTPNavy},inner xsep=3mm,
    inner ysep=3.5mm,line width=0.65pt},
  cycle/.style={-{Stealth[length=2.1mm]},semithick,draw=UTPMuted,
    rounded corners=4mm,shorten <=7pt,shorten >=7pt}
]
\node[core] (guide) at (0,-2.55) {TA\\WORKFLOW};
\node[stage] (prepare) at (0,0) {SESSION\\PLANNING};
\node[stage] (explain) at (4.75,-2.55) {INSTRUCTIONAL\\SUPPORT};
\node[stage] (facilitate) at (0,-5.10) {STUDENT\\ENGAGEMENT};
\node[stage] (assess) at (-4.75,-2.55) {LEARNING\\ASSESSMENT};

\draw[cycle] (prepare.east) -- (explain.north);
\draw[cycle] (explain.south) -- (facilitate.east);
\draw[cycle] (facilitate.west) -- (assess.south);
\draw[cycle] (assess.north) -- (prepare.west);
\end{tikzpicture}

\vspace{0.72cm}
\begin{tcolorbox}[
  enhanced,width=0.82\linewidth,colback=white,colframe=UTPBlue!18,
  arc=1.5mm,boxrule=0.45pt,left=6mm,right=6mm,top=3.2mm,bottom=3.2mm,
  borderline west={1.6pt}{0pt}{UTPGold}]
\centering\normalsize\color{UTPInk}
{\large\bfseries\color{UTPBlue}\utpoutlineheading\par}
\vspace{0.3em}
\utpcourseoutline
\end{tcolorbox}

\vfill
\begin{minipage}{0.82\textwidth}
\begin{tabular}{@{}>{\RaggedRight\arraybackslash}m{0.42\linewidth}@{\hspace{0.10\linewidth}}>{\RaggedLeft\arraybackslash}m{0.48\linewidth}@{}}
{\normalsize\bfseries\color{UTPBlue}PREPARED BY\par\vspace{0.25em}
 \Large\bfseries\color{UTPNavy}\utpauthor\par\vspace{0.12em}
 \normalfont\small\color{UTPMuted}\utpauthorrole}
&
\includegraphics[height=2.10cm,keepaspectratio]{logo-DC.png}
\end{tabular}
\end{minipage}
\vspace{0.5cm}
\end{titlepage}%
  \cleardoublepage
  \pagenumbering{roman}%
  \setcounter{tocdepth}{0}%
  \utptableofcontents
  \newpage
  \hypersetup{pageanchor=true}%
  \ifx\utpcoursemap\empty\else\input{\utpcoursemap}\fi
  \cleardoublepage
  \ifdefined\ReadonlyShipoutCounter
    \immediate\write\csname @auxout\endcsname
      {\string\gdef\string\utpfrontpages{\the\ReadonlyShipoutCounter}}%
  \fi
  \pagenumbering{arabic}%
}

% --------------------------------------------------------------------
%  15.  LISTING STYLES
% --------------------------------------------------------------------
%  Two base looks, and one named style per language built on them, so
%  the colours and frame are the same in every book:
%
%      utpcode          the base for source code
%      utpterminal      the base for commands typed and output printed
%      utppython        Python
%      utpr             R
%      utpc             C
%      utpcpp           C++
%      utpbash          shell scripts and commands
%      utpcnumbered     C with line numbers, for a complete source file
%      utpbashnumbered  shell with line numbers, likewise
%
%      \begin{lstlisting}[style=utppython] ... \end{lstlisting}
%      \lstinputlisting[style=utpcnumbered]{all_experiments/Lab04/fcfs.c}
%
%  A book that sets most of its listings in one style makes it the
%  default once, after the input:  \lstset{style=utpcpp}

\lstdefinestyle{utpcode}{
  basicstyle=\ttfamily\scriptsize,
  keywordstyle=\bfseries\color{UTPBlue},
  commentstyle=\itshape\color{UTPMuted},
  stringstyle=\color{UTPGreen!70!black},
  backgroundcolor=\color{UTPPaleGray!72},
  frame=single,
  rulecolor=\color{UTPBlue!28},
  framesep=5pt,
  breaklines=true,
  breakatwhitespace=false,
  showstringspaces=false,
  keepspaces=true,
  columns=fullflexible,
  upquote=true,
  tabsize=4,
  aboveskip=0.7em,
  belowskip=0.8em
}
\lstdefinestyle{utpterminal}{
  basicstyle=\ttfamily\scriptsize,
  backgroundcolor=\color{UTPPaleGray!72},
  frame=single,
  rulecolor=\color{UTPMuted!35},
  framesep=5pt,
  breaklines=true,
  showstringspaces=false,
  keepspaces=true,
  columns=fullflexible,
  upquote=true,
  aboveskip=0.6em,
  belowskip=0.7em
}

\lstdefinestyle{utppython}{style=utpcode,language=Python}
\lstdefinestyle{utpr}{style=utpcode,language=R,tabsize=2}
\lstdefinestyle{utpc}{style=utpcode,language=C}
\lstdefinestyle{utpcpp}{style=utpcode,language=C++}
\lstdefinestyle{utpbash}{style=utpcode,language=bash,
  rulecolor=\color{UTPMuted!35},aboveskip=0.6em,belowskip=0.7em}
\lstdefinestyle{utpcnumbered}{style=utpc,numbers=left,numberstyle=\tiny\color{UTPMuted},
  numbersep=7pt,xleftmargin=1.7em,framexleftmargin=1.4em}
\lstdefinestyle{utpbashnumbered}{style=utpbash,numbers=left,numberstyle=\tiny\color{UTPMuted},
  numbersep=7pt,xleftmargin=1.7em,framexleftmargin=1.4em}

% Pseudocode. A book that sets algorithms asks for them once, after the
% input, and gets algorithm2e loaded and coloured to match:
%
%     \utpusealgorithms
%
% \algostyle then sets one algorithm environment a size smaller than the
% text around it.
\newcommand{\utpalgorithmstyle}{%
  \SetAlCapNameFnt{\bfseries\color{UTPNavy}}%
  \SetAlCapFnt{\bfseries\color{UTPNavy}}%
  \SetAlgoCaptionSeparator{.}%
  \SetKwInput{KwIn}{Input}%
  \SetKwInput{KwOut}{Output}%
  \SetNlSty{color}{UTPMuted}{}%
  \SetCommentSty{\itshape\color{UTPMuted}}%
  \SetKwSty{bfseries\color{UTPBlue}}%
  \SetArgSty{}%
  \providecommand{\algostyle}{\small\SetAlgoNlRelativeSize{-1}\DontPrintSemicolon\SetAlgoVlined}%
}
\newcommand{\utpusealgorithms}{%
  \usepackage[linesnumbered,ruled,vlined]{algorithm2e}%
  \utpalgorithmstyle
}

%
%      \utpcontentspage{\small}{\normalsize}{0.65em}{0.07em}{0.16em}
%
%      \begin{document}
%      \utpfrontmatter                     % cover, contents, map
%      \part{Orientation}
%      \input{chapters/01-where-to-begin}
%      \end{document}
%
%  Configuration goes BEFORE the \input; overrides go AFTER it.
%  Paths are relative, so compile from inside the course folder.
%
%  ------------------------------------------------------------------
%  NAMING
%  ------------------------------------------------------------------
%  One name per idea, used by every book. Anything a book can use is
%  either an unprefixed learning box (keyidea, pitfall, labproject,
%  ...) or carries the utp prefix (utpfigure, utptable, utpfig,
%  utpsteps, utppython, UTPNavy, utpcell, ...). Course prefixes such
%  as ml-, ds- or erp- are not used: a chapter moved from one book to
%  another compiles unchanged. .shared/README.md lists every name.
%
%  Adding a new course? Copy .shared/TEMPLATE-main.tex and follow
%  .shared/ADDING-A-COURSE.md. Every step is a numbered checklist.
%
%  ------------------------------------------------------------------
%  CONTENTS
%  ------------------------------------------------------------------
%    1  Course configuration     the knobs a book may set
%    2  Packages                 the shared package set and fonts
%    3  Colour palette           the one UTP palette
%    4  Page layout and tables   margins, lists, table styling
%    5  Table of contents        parts, chapters, labs, figures
%    6  Headings                 chapter badge, part pages, sections
%    7  Page furniture           running heads and footers
%    8  Learning boxes           key idea, pitfall, try it, and friends
%    9  Small commands           \term, \cmd, \chapterguide, utpsteps
%   10  Diagram node styles      the TikZ vocabulary every book shares
%   11  Figure frames            utpfigure, utpfig, source images
%   12  Lab sheets               the lab banner and the lab toolkit
%   13  Stage divider            the part-opening page in longer books
%   14  Front matter             \utpfrontmatter: cover, contents, map
%   15  Listing styles           code, terminal and pseudocode
% =====================================================================

% --------------------------------------------------------------------
%  1.  COURSE CONFIGURATION
% --------------------------------------------------------------------
%  The knobs below, plus \utpstyleparts in section 5 and the box titles
%  in section 8, are the whole configuration surface. A book sets what
%  it needs BEFORE the \input and inherits the rest.
%
%  \providecommand means "define unless the book already did", which is
%  what makes the \def-before-\input pattern work.
%
%  Running text
%    \utprunninghead   the left running head on every page.
%    \utpchapterlabel  the word in the chapter badge. CHAPTER by
%                      default; a book whose chapters are labs sets LAB.
%    \utplabtotal      the N in the "Lab 3 of N" banners.

\providecommand{\utprunninghead}{Course Notes}
\providecommand{\utpchapterlabel}{CHAPTER}
\providecommand{\utplabtotal}{10}

%  Cover, printed by section 14
%    \utpcoursetitle       the title on the title page.
%    \utpdocumentsubtitle  the line under the title; Teaching Assistant Guide by default.
%    \utpcourseoutline     the sentence or two below it; \\ breaks a line.
%    \utpoutlineheading    the small heading above that sentence.
%    \utpauthor            the name in the PREPARED BY block.
%    \utpauthorrole        the line under the name.

\providecommand{\utpcoursetitle}{\utprunninghead}
\providecommand{\utpdocumentsubtitle}{Teaching Assistant Guide}
\providecommand{\utpcourseoutline}{}
\providecommand{\utpoutlineheading}{COURSE OUTLINE}
\providecommand{\utpauthor}{Reeyan Afzal Khan}
\providecommand{\utpauthorrole}{Graduate Research Assistant}

%  Files and folders
%    \utpextragraphicspath  image folders to search before the standard
%                           ones, as brace groups: {{screenshots/}}.
%                           A course's own figures/ folder and the
%                           shared logos are always searched.
%    \utpcoursemap          the overview chapter \utpfrontmatter inputs
%                           after the contents. Empty means none.

\providecommand{\utpextragraphicspath}{}
\providecommand{\utpcoursemap}{chapters/course-map}

%  Contents entries for chapters (section 5 applies them). A book whose
%  chapters are labs says so once here rather than restating the three
%  tocloft settings after the input.
%    \utpchaptertocword      Chapter~ by default, Lab~ in one book.
%    \utpchaptertocindent    left indent of the entry.
%    \utpchaptertocnumwidth  room reserved for the word and number.

\providecommand{\utpchaptertocword}{Chapter~}
\providecommand{\utpchaptertocindent}{1.5em}
\providecommand{\utpchaptertocnumwidth}{6em}

%  Monospace face (section 2)
%    \utpmonofont   sourcecodepro by default. inconsolata is the
%                   alternative, for a book whose pages are mostly
%                   device transcripts.

\providecommand{\utpmonofont}{sourcecodepro}

% --------------------------------------------------------------------
%  2.  PACKAGES
% --------------------------------------------------------------------
%  The union of what the books need. A course that needs something
%  extra loads it in its own file, under the heading "packages only
%  this book needs". Pseudocode is the one optional package the shared
%  file knows about: \utpusealgorithms in section 15 loads and styles it.
%
%  The tikz and pgfplots library lists are the union across all books,
%  so a diagram or chart copied from one course compiles in another.

\usepackage[T1]{fontenc}
\IfFileExists{glyphtounicode.tex}{\input{glyphtounicode}\pdfgentounicode=1}{}
\usepackage[american]{babel}
\IfFileExists{sourcesanspro.sty}
  {\usepackage[default]{sourcesanspro}}
  {\usepackage[scaled]{helvet}\renewcommand{\familydefault}{\sfdefault}}
\DeclareFontShape{T1}{cmss}{b}{n}{<->ssub*cmss/bx/n}{}
\DeclareFontShape{TS1}{cmss}{b}{n}{<->ssub*cmss/bx/n}{}
% Monospace. Without this the tt default stays Computer Modern Typewriter,
% a serif face that clashes with the sans body text in every listing and
% every \cmd. Source Code Pro is the companion to Source Sans Pro; it is
% scaled down because mono digits otherwise sit large next to the sans.
%
% Inconsolata, chosen with \def\utpmonofont{inconsolata}, distinguishes
% 0/O and 1/l/I at a glance, which matters when a reader copies an
% interface name or an address out of a transcript. Only one mono package
% is ever loaded: a second would win the family but inherit the first's
% series names, and every \ttfamily would then resolve by substitution.
\def\utpmonoinconsolata{inconsolata}
\ifx\utpmonofont\utpmonoinconsolata
  \IfFileExists{inconsolata.sty}{\usepackage[varqu,varl,scaled=1.02]{inconsolata}}{}
\else
  \IfFileExists{sourcecodepro.sty}
    {\usepackage[scaled=0.88]{sourcecodepro}}
    {\usepackage[scaled=0.86]{beramono}}
  % Source Sans names its normal weight "regular", but the Computer Modern
  % fallbacks call it "m". Map one onto the other so a \ttfamily reached
  % while the regular series is active resolves instead of warning.
  \DeclareFontShape{T1}{cmtt}{regular}{n}{<->ssub*cmtt/m/n}{}
  \DeclareFontShape{TS1}{cmtt}{regular}{n}{<->ssub*cmtt/m/n}{}
  \DeclareFontShape{T1}{cmtt}{regular}{it}{<->ssub*cmtt/m/it}{}
\fi
\usepackage[margin=21mm,headheight=22pt]{geometry}
\usepackage{amsmath,amssymb,mathtools,bm}
\usepackage{microtype}
\usepackage{enumitem}
\usepackage{booktabs}
\usepackage{tabularx}
\usepackage{array}
\usepackage{ragged2e}
\usepackage{etoolbox}
\usepackage{titlesec}
\usepackage{tocloft}
\usepackage{fancyhdr}
\usepackage[table]{xcolor}
\usepackage{tcolorbox}
\tcbuselibrary{breakable,skins}
\usepackage[hidelinks]{hyperref}
\usepackage{bookmark}
% Document properties, so the PDF identifies itself in a reader's
% title bar, in a library catalogue, and in search results. Set from
% the same configuration macros the cover uses.
\hypersetup{
  pdftitle={\utpcoursetitle},
  pdfauthor={\utpauthor},
  pdfsubject={Universiti Teknologi PETRONAS course notes and laboratory guide},
  pdfkeywords={UTP, course notes, laboratory manual, \utpcoursetitle},
  pdfdisplaydoctitle=true,
  bookmarksnumbered=true,
  bookmarksopenlevel=1
}
\usepackage{listings}
\usepackage{graphicx}
\usepackage{caption}
\usepackage{tikz}
\usetikzlibrary{angles,arrows.meta,backgrounds,calc,chains,decorations.pathmorphing,decorations.pathreplacing,fit,matrix,patterns,positioning,quotes,shadows.blur,shapes.geometric,shapes.misc,shapes.symbols}
\usepackage{pgfplots}
% statistics gives boxplot prepared; groupplots gives multi-panel charts.
\usepgfplotslibrary{statistics}
\usepgfplotslibrary{groupplots}

% Images: any folders the book names first, then its own figures/
% folder, then the shared logos used by the cover.
\expandafter\graphicspath\expandafter{\utpextragraphicspath{figures/}{../.shared/logos/}}

% --------------------------------------------------------------------
%  3.  COLOUR PALETTE
% --------------------------------------------------------------------
%  One palette, defined once under UTP* names, and used by every book,
%  box, diagram and chart. Five accents (blue, teal, green, orange, red)
%  with a pale tint of each for box backgrounds, plus navy for headings,
%  gold for the cover, ink and muted for text, and a sixth accent,
%  purple, for charts and layered diagrams that need one more series.

\definecolor{UTPNavy}{HTML}{18324A}
\definecolor{UTPBlue}{HTML}{246B9E}
\definecolor{UTPTeal}{HTML}{138A80}
\definecolor{UTPGreen}{HTML}{2D7D46}
\definecolor{UTPOrange}{HTML}{C86B24}
\definecolor{UTPRed}{HTML}{B34843}
\definecolor{UTPPurple}{HTML}{6B4FA0}
\definecolor{UTPGold}{HTML}{D4A13E}
\definecolor{UTPInk}{HTML}{263238}
\definecolor{UTPMuted}{HTML}{60717C}
\definecolor{UTPPaleBlue}{HTML}{EEF6FB}
\definecolor{UTPPaleTeal}{HTML}{EEF8F6}
\definecolor{UTPPaleOrange}{HTML}{FFF5EA}
\definecolor{UTPPaleRed}{HTML}{FFF0EF}
\definecolor{UTPPalePurple}{HTML}{F0ECF9}
\definecolor{UTPPaleGray}{HTML}{F4F6F7}
\color{UTPInk}

% --------------------------------------------------------------------
%  4.  PAGE LAYOUT AND TABLES
% --------------------------------------------------------------------
%  Paragraph spacing, list indents, and the table look: alternating
%  row tint, top-aligned paragraph columns, ragged-right cells.
%
%  Two extra column types for tabularx tables:
%      D{width}   a fixed-width, ragged-right, vertically centred column
%      Y          a ragged-right column that fills the remaining width
%
%  Every table in every book has the same shape, the one Algorithms &
%  Data Structures set: a utptable frame around a full-width tabularx,
%  booktabs rules, a tinted header row of \utpth cells (centred, bold,
%  navy), and a first column in bold navy when it names the rows.
%  Cells are vertically centred, so use m{} rather than p{} columns.
%
%      \begin{utptable}
%      \begin{tabularx}{\linewidth}{@{}>{\bfseries\leavevmode\color{UTPNavy}}m{0.22\linewidth}X@{}}
%      \toprule
%      \rowcolor{UTPPaleBlue!92}
%      \utpth{Operation} & \utpth{What it does} \\
%      \midrule
%      Insertion & Add a new element. \\
%      \bottomrule
%      \end{tabularx}
%      \end{utptable}
%
%  A book whose tables carry dense text asks for \utpcompacttables after
%  the input: tabularx tables then set in \footnotesize with tighter
%  padding and vertically centred cells.

\setlength{\parindent}{0pt}
\setlength{\parskip}{0.55em}
\setlist[itemize]{leftmargin=1.45em,itemsep=0.22em,topsep=0.25em}
\setlist[enumerate]{leftmargin=1.75em,itemsep=0.28em,topsep=0.3em}
\setlist[description]{font=\bfseries\color{UTPNavy},itemsep=0.25em,topsep=0.3em}
\renewcommand{\arraystretch}{1.34}
\setlength{\tabcolsep}{6.5pt}
\arrayrulecolor{UTPBlue!45}
% Keep paragraph-style columns ragged and vertically centred. Tables use
% fixed m columns beside flexible X columns; giving X the same centring
% keeps a short cell level with a tall neighbour in a multi-line row.
\renewcommand{\tabularxcolumn}[1]{m{#1}}
\makeatletter
\g@addto@macro\@arrayparboxrestore{\RaggedRight\arraybackslash}
\makeatother
% A header cell: centred, bold and navy, on the tinted header row.
\newcommand{\utpth}[1]{\multicolumn{1}{c}{\bfseries\color{UTPNavy}#1}}
\AtBeginEnvironment{tabular}{\small\setlength{\extrarowheight}{0.6pt}\rowcolors{2}{UTPPaleBlue!72}{white}}
\AtBeginEnvironment{tabularx}{\small\setlength{\extrarowheight}{0.6pt}\rowcolors{2}{UTPPaleBlue!72}{white}}
\AtBeginEnvironment{figure}{\centering}
\newcolumntype{D}[1]{>{\raggedright\arraybackslash}m{#1}}
\newcolumntype{Y}{>{\raggedright\arraybackslash}X}
\newcommand{\utpcompacttables}{%
  \renewcommand{\tabularxcolumn}[1]{m{##1}}%
  \AtBeginEnvironment{tabularx}{%
    \footnotesize
    \setlength{\tabcolsep}{4pt}%
    \renewcommand{\arraystretch}{1.25}%
    \setlength{\extrarowheight}{0.6pt}%
  }%
}
\emergencystretch=1.5em
\hbadness=10000
\hfuzz=8pt

% Page-breaking discipline. Without these a paragraph can leave a single
% line stranded at the foot of one page or the head of the next, which is
% the most visible sign of unattended typesetting.
\widowpenalty=10000          % no first line of a paragraph left alone at a page foot
\clubpenalty=10000           % no last line of a paragraph left alone at a page head
\displaywidowpenalty=10000   % the same around displayed equations
\brokenpenalty=10000         % do not break a page after a hyphenated line
\predisplaypenalty=10000     % keep the line before a display with it
% Pages are built from boxes of many heights, so let the last line float
% up rather than stretching the gaps between boxes to force a flush foot.
\raggedbottom
% Prefer a slightly loose line to an overfull one, but do not hyphenate
% freely -- short sans lines look ragged when over-hyphenated.
\tolerance=1200
\hyphenpenalty=200

% Character protrusion and font expansion. Both are what stop a justified
% sans column from showing visible rivers down the page.
\microtypesetup{protrusion=true,expansion=true,tracking=true}

\pgfplotsset{
  compat=1.18,
  every axis/.append style={
    axis line style={draw=UTPMuted!85,line width=0.55pt},
    tick style={draw=UTPMuted!85},
    tick label style={font=\scriptsize\color{UTPMuted}},
    label style={font=\small\color{UTPNavy}},
    title style={font=\small\bfseries\color{UTPNavy}},
    legend style={font=\scriptsize,draw=UTPBlue!22,fill=white,rounded corners=1pt}
  },
  every axis plot/.append style={line width=1.15pt}
}
\tikzset{every picture/.append style={line cap=round,line join=round}}

% --------------------------------------------------------------------
%  5.  TABLE OF CONTENTS
% --------------------------------------------------------------------
%  Formatting for part, chapter and lab entries, and for a list of
%  figures if a book prints one. The lab level is a custom tocloft list
%  so lab banners appear in the contents.
%
%  Part styling is skipped when \utpstyleparts is 0, which is what a
%  book without \part divisions should set.
%
%  The contents page is meant to fill exactly one page, and how loosely
%  it is set depends on how many entries a book has. One line after the
%  input tunes it:
%
%      \utpcontentspage{<base size>}{<part size>}%
%                      {<space before a part>}{<before a chapter>}{<before a lab>}
%      \utpcontentspage{\normalsize}{\large}{1.15em}{0.38em}{0.30em}   % the default

\addto\captionsamerican{\renewcommand{\contentsname}{Learning Path}}
\newlistentry{lab}{toc}{-1}
\providecommand{\utpstyleparts}{1}
\ifnum\utpstyleparts=1
\renewcommand{\cftpartpresnum}{}
\renewcommand{\cftpartaftersnum}{}
\setlength{\cftpartnumwidth}{5.2em}
\setlength{\cftbeforepartskip}{0.38em}
\renewcommand{\cftpartfont}{\normalsize\bfseries\color{UTPNavy}}
\renewcommand{\cftpartpagefont}{\normalsize\bfseries\color{UTPNavy}}
\renewcommand{\cftpartleader}{\hfill}
\fi
\renewcommand{\cftchappresnum}{\utpchaptertocword}
\renewcommand{\cftchapaftersnum}{\quad}
\setlength{\cftchapindent}{\utpchaptertocindent}
\setlength{\cftchapnumwidth}{\utpchaptertocnumwidth}
\setlength{\cftbeforechapskip}{0pt}
\renewcommand{\cftchapfont}{\normalfont}
\renewcommand{\cftchappagefont}{\normalfont}
\renewcommand{\cftchapleader}{\cftdotfill{\cftdotsep}}
\setlength{\cftlabindent}{3em}
\setlength{\cftlabnumwidth}{0em}
\setlength{\cftbeforelabskip}{0.10em}
\renewcommand{\cftlabfont}{\bfseries\color{UTPOrange!90!black}}
\renewcommand{\cftlabpagefont}{\bfseries\color{UTPOrange!90!black}}
\renewcommand{\cftlableader}{\color{UTPOrange!65!UTPMuted}\cftdotfill{\cftlabdotsep}}

% A list of figures, for a book that prints one (\listoffigures).
\renewcommand{\cftfigfont}{\normalfont\small}
\renewcommand{\cftfigpagefont}{\normalfont\small}
\renewcommand{\cftfigpresnum}{Figure~}
\renewcommand{\cftfigaftersnum}{\quad}
\setlength{\cftfigindent}{1.5em}
\setlength{\cftfignumwidth}{5.6em}
\setlength{\cftbeforefigskip}{0pt}
\renewcommand{\cftfigleader}{\cftdotfill{\cftdotsep}}

% The contents page itself, printed by \utpfrontmatter (section 14).
\newcommand{\utpcontentspage}[5]{%
  \def\utptocsize{#1}\def\utptocpartsize{#2}%
  \def\utptocpartskip{#3}\def\utptocchapterskip{#4}\def\utptoclabskip{#5}}
\utpcontentspage{\normalsize}{\large}{1.15em}{0.38em}{0.30em}
\newcommand{\utptableofcontents}{%
  \begingroup
  \utptocsize
  \setlength{\parskip}{0pt}%
  \setlength{\cftbeforepartskip}{\utptocpartskip}%
  \setlength{\cftbeforechapskip}{\utptocchapterskip}%
  \setlength{\cftbeforelabskip}{\utptoclabskip}%
  \renewcommand{\cftpartfont}{\utptocpartsize\bfseries\color{UTPNavy}}%
  \renewcommand{\cftpartpagefont}{\utptocpartsize\bfseries\color{UTPNavy}}%
  \tableofcontents
  \endgroup
}

% --------------------------------------------------------------------
%  6.  HEADINGS
% --------------------------------------------------------------------
%  The chapter badge (a coloured block with a word and the number),
%  part pages, and section spacing.
%
%  The word in the badge is \utpchapterlabel; Data Communication sets
%  it to LAB because each of its chapters is a lab.

\titleformat{\chapter}[display]
  {\bfseries\color{UTPNavy}}
  {\filleft\colorbox{UTPBlue}{\parbox{2.55cm}{\centering\color{white}\Large
    \utpchapterlabel\\[-0.15em]\Huge\thechapter}}}
  {0.65em}{\Huge\raggedright}
  [\vspace{0.25em}{\color{UTPBlue}\titlerule[1.2pt]}]
\renewcommand{\thepart}{Stage~\arabic{part}}
\titleformat{\part}[display]
  {\centering\color{UTPNavy}\bfseries}
  {\Large\color{UTPTeal}\thepart}{0.45em}{\Huge}
  [\vspace{0.55em}\color{UTPTeal}\rule{0.45\linewidth}{1.5pt}]
\titleformat{\section}{\Large\bfseries\color{UTPNavy}}{\thesection}{0.55em}{}
\titleformat{\subsection}{\large\bfseries\color{UTPBlue}}{\thesubsection}{0.5em}{}
\titlespacing*{\chapter}{0pt}{-1.8em}{1.0em}
\titlespacing*{\part}{0pt}{0pt}{1em}
\titlespacing*{\section}{0pt}{1.2em}{0.35em}
\titlespacing*{\subsection}{0pt}{0.95em}{0.2em}

% --------------------------------------------------------------------
%  7.  PAGE FURNITURE
% --------------------------------------------------------------------
%  Running heads and footers. The left header is \utprunninghead;
%  the right one follows the current chapter.

\pagestyle{fancy}
\fancyhf{}
\lhead{\small\color{UTPMuted}\utprunninghead}
\rhead{\small\color{UTPMuted}\nouppercase{\leftmark}}
\cfoot{\small\color{UTPMuted}\thepage}
\renewcommand{\headrulewidth}{0.5pt}
\renewcommand{\headrule}{\hbox to\headwidth{\color{UTPPaleBlue}\leaders\hrule height \headrulewidth\hfill}}
\renewcommand{\chaptermark}[1]{\markboth{#1}{}}
\fancypagestyle{plain}{%
  \fancyhf{}%
  \cfoot{\small\color{UTPMuted}\thepage}%
  \renewcommand{\headrulewidth}{0pt}%
}

% --------------------------------------------------------------------
%  8.  LEARNING BOXES
% --------------------------------------------------------------------
%  The coloured callouts the chapters use. Every title is a macro, so
%  a book can rename a box without restating its geometry:
%
%      \def\utptitlepitfall{Technical note}   % before the input
%
%  To change how a box LOOKS in one book, use one of the named
%  variations at the end of this section, or \renewtcolorbox after the
%  input. To add a box only one book needs, declare it in that book
%  with \utpnewcallout or \utpnewquietbox.

% Box titles -- override any of these before the input line.
\providecommand{\utptitlekeyidea}{Key idea}
\providecommand{\utptitleworked}{Worked example}
\providecommand{\utptitlepitfall}{Common mistake}
\providecommand{\utptitletryit}{Try it yourself}
\providecommand{\utptitlequickcheck}{Check your understanding}
\providecommand{\utptitlerecap}{Chapter recap}
\providecommand{\utptitledeliverable}{Lab deliverable}
\providecommand{\utptitlelabmode}{How much is scaffolded?}
\providecommand{\utptitleintuition}{Plain idea}
\providecommand{\utptitlefirstread}{Core path}
\providecommand{\utptitleglance}{Chapter at a glance}
\providecommand{\utptitledebug}{Debugging task}
\providecommand{\utptitleoptional}{Beyond the core}
\providecommand{\utptitleverification}{Verification checklist}
\providecommand{\utptitletroubleshoot}{Troubleshooting}
\providecommand{\utptitlesourcebasis}{Source basis}

% Explaining an idea
\newtcolorbox{keyidea}{enhanced,breakable,colback=UTPPaleTeal,colframe=UTPTeal, boxrule=0.8pt,arc=2mm,left=2.5mm,right=2.5mm,top=1.8mm,bottom=1.8mm, fonttitle=\bfseries\color{UTPNavy},title=\utptitlekeyidea}
\newtcolorbox{workedexample}{enhanced,breakable,colback=UTPPaleBlue,colframe=UTPBlue, boxrule=0.8pt,arc=2mm,left=2.5mm,right=2.5mm,top=1.8mm,bottom=1.8mm, fonttitle=\bfseries\color{UTPNavy},title=\utptitleworked}
\newtcolorbox{pitfall}{enhanced,breakable,colback=UTPPaleRed,colframe=UTPRed, boxrule=0.8pt,arc=2mm,left=2.5mm,right=2.5mm,top=1.8mm,bottom=1.8mm, fonttitle=\bfseries\color{UTPNavy},title=\utptitlepitfall}
% Asking the reader to do something
\newtcolorbox{tryit}{enhanced,breakable,colback=white,colframe=UTPTeal, boxrule=0.9pt,arc=2mm,left=2.5mm,right=2.5mm,top=1.8mm,bottom=1.8mm, fonttitle=\bfseries\color{white},colbacktitle=UTPTeal,title=\utptitletryit}
\newtcolorbox{quickcheck}{enhanced,breakable,colback=white,colframe=UTPOrange, boxrule=0.9pt,arc=2mm,left=2.5mm,right=2.5mm,top=1.8mm,bottom=1.8mm, fonttitle=\bfseries\color{white},colbacktitle=UTPOrange,title=\utptitlequickcheck}
\newtcolorbox{chapterrecap}{enhanced,breakable,colback=UTPPaleTeal,colframe=UTPGreen, boxrule=0.9pt,arc=2mm,left=2.5mm,right=2.5mm,top=1.8mm,bottom=1.8mm, fonttitle=\bfseries\color{white},colbacktitle=UTPGreen,title=\utptitlerecap}
% Lab structure
\newtcolorbox{labdeliverable}{enhanced,breakable,colback=white,colframe=UTPGreen, boxrule=0.9pt,arc=2mm,left=2.6mm,right=2.6mm,top=1.9mm,bottom=1.9mm, fonttitle=\bfseries\color{white},colbacktitle=UTPGreen,title=\utptitledeliverable}
\newtcolorbox{labmode}{enhanced,breakable,colback=UTPPaleBlue!52,colframe=UTPBlue, boxrule=0.75pt,arc=2mm,left=2.6mm,right=2.6mm,top=1.7mm,bottom=1.7mm, fonttitle=\bfseries\color{white},colbacktitle=UTPBlue,title=\utptitlelabmode}
% Softer, lighter-weight asides
\newtcolorbox{intuition}{enhanced,breakable,colback=white,colframe=UTPOrange!42, boxrule=0.5pt,arc=2mm,left=2.5mm,right=2.5mm,top=1.5mm,bottom=1.5mm, fonttitle=\bfseries\color{UTPOrange!85!black},title=\utptitleintuition}
\newtcolorbox{firstread}{enhanced,breakable,colback=white,colframe=UTPBlue!42, boxrule=0.5pt,arc=2mm,left=2.5mm,right=2.5mm,top=1.5mm,bottom=1.5mm, fonttitle=\bfseries\color{UTPBlue}, title=\utptitlefirstread}
\newtcolorbox{debugtask}{enhanced,breakable,colback=UTPPaleRed!58,colframe=UTPRed, boxrule=0.85pt,arc=2mm,left=2.6mm,right=2.6mm,top=1.9mm,bottom=1.9mm, fonttitle=\bfseries\color{white},colbacktitle=UTPRed,title=\utptitledebug}
\newtcolorbox{analogy}{enhanced,breakable,colback=white,colframe=UTPOrange!38, boxrule=0.5pt,arc=2mm,left=2.5mm,right=2.5mm,top=1.5mm,bottom=1.5mm, fonttitle=\bfseries\color{UTPOrange!85!black},title=Analogy}
\newtcolorbox{coursechoice}{enhanced,breakable,colback=white,colframe=UTPNavy, boxrule=0.8pt,arc=2mm,left=2.8mm,right=2.8mm,top=2mm,bottom=2mm, fonttitle=\bfseries\color{white},colbacktitle=UTPNavy,title=Course choice}
\newtcolorbox{decisionmemo}{enhanced,breakable,colback=UTPPaleOrange!55,colframe=UTPOrange, boxrule=0.85pt,arc=2mm,left=2.6mm,right=2.6mm,top=1.9mm,bottom=1.9mm, fonttitle=\bfseries\color{white},colbacktitle=UTPOrange,title=Decision memo}
\newtcolorbox{libnote}{enhanced,breakable,colback=white,colframe=UTPMuted!48, boxrule=0.5pt,arc=2mm,left=2.5mm,right=2.5mm,top=1.5mm,bottom=1.5mm, fonttitle=\bfseries\color{UTPNavy},title=Reference note}
\newtcolorbox{miniproject}{enhanced,breakable,colback=UTPPaleTeal!52,colframe=UTPGreen, boxrule=0.9pt,arc=2mm,left=2.6mm,right=2.6mm,top=1.9mm,bottom=1.9mm, fonttitle=\bfseries\color{white},colbacktitle=UTPGreen,title=Mini-project}
% Where the material came from: the lecturer's slides, manual or tutorial
% that a chapter is built on.
\newtcolorbox{sourcebasis}{enhanced,breakable,colback=white,colframe=UTPMuted!55, boxrule=0.5pt,arc=2mm,left=2.5mm,right=2.5mm,top=1.5mm,bottom=1.5mm, fonttitle=\bfseries\color{UTPMuted},title=\utptitlesourcebasis}
% Boxes that take a title argument
\newtcolorbox{modelcard}[1]{enhanced,breakable,colback=white,colframe=UTPBlue, boxrule=0.9pt,arc=2mm,left=2.5mm,right=2.5mm,top=1.8mm,bottom=1.8mm, fonttitle=\bfseries\color{white},colbacktitle=UTPBlue,title=#1}
\newtcolorbox{optionaltopic}{enhanced,breakable,colback=white,colframe=UTPMuted!58, boxrule=0.5pt,arc=2mm,left=2.5mm,right=2.5mm,top=1.5mm,bottom=1.5mm, fonttitle=\bfseries\color{UTPMuted},title=\utptitleoptional}
\newtcolorbox{secondpass}{enhanced,breakable,colback=UTPPaleGray!82,colframe=UTPMuted, boxrule=0.75pt,arc=2mm,left=2.6mm,right=2.6mm,top=1.9mm,bottom=1.9mm, fonttitle=\bfseries\color{white},colbacktitle=UTPMuted,title=Second pass}
\newtcolorbox{labproject}[1]{enhanced,breakable,colback=white,colframe=UTPBlue, boxrule=0.9pt,arc=2mm,left=2.6mm,right=2.6mm,top=1.9mm,bottom=1.9mm, fonttitle=\bfseries\color{white},colbacktitle=UTPBlue,title={#1}}
% Evidence that a lab worked, and what to do when it did not
\newtcolorbox{verification}{enhanced,breakable,colback=UTPPaleTeal!65,colframe=UTPGreen, boxrule=0.8pt,arc=2mm,left=2.6mm,right=2.6mm,top=1.8mm,bottom=1.8mm, fonttitle=\bfseries\color{white},colbacktitle=UTPGreen,title=\utptitleverification}
\newtcolorbox{troubleshoot}{enhanced,breakable,colback=UTPPaleOrange!70,colframe=UTPOrange, boxrule=0.8pt,arc=2mm,left=2.6mm,right=2.6mm,top=1.8mm,bottom=1.8mm, fonttitle=\bfseries\color{white},colbacktitle=UTPOrange,title=\utptitletroubleshoot}

% Named variations a book asks for after the input. Each restyles boxes
% that already exist, so the chapters do not change.
%
%   \utpheavycallouts          a solid colour bar behind the titles of
%                              keyidea, workedexample, intuition and
%                              firstread, with a firmer frame.
%   \utpbreakableframes        let utpfigure and utptable split over a
%                              page break, in a lighter blue (section 11).
%   \utplabprojectcolour{...}  the labproject box in another colour, with
%                              a little more room inside.
\newcommand{\utplabprojectcolour}[1]{%
  \renewtcolorbox{labproject}[1]{enhanced,breakable,colback=white,colframe=#1, boxrule=0.9pt,arc=2mm,left=2.7mm,right=2.7mm,top=2mm,bottom=2mm, fonttitle=\bfseries\color{white},colbacktitle=#1,title={##1}}%
}
\newcommand{\utpheavycallouts}{%
  \renewtcolorbox{keyidea}{enhanced,breakable,colback=UTPPaleTeal,colframe=UTPTeal, boxrule=0.85pt,arc=2mm,left=2.5mm,right=2.5mm,top=1.8mm,bottom=1.8mm, fonttitle=\bfseries\color{white},colbacktitle=UTPTeal,title=\utptitlekeyidea}%
  \renewtcolorbox{workedexample}{enhanced,breakable,colback=UTPPaleBlue,colframe=UTPBlue, boxrule=0.8pt,arc=2mm,left=2.5mm,right=2.5mm,top=1.8mm,bottom=1.8mm, fonttitle=\bfseries\color{white},colbacktitle=UTPBlue,title=\utptitleworked}%
  \renewtcolorbox{intuition}{enhanced,breakable,colback=UTPPaleOrange,colframe=UTPOrange, boxrule=0.8pt,arc=2mm,left=2.5mm,right=2.5mm,top=1.8mm,bottom=1.8mm, fonttitle=\bfseries\color{white},colbacktitle=UTPOrange,title=\utptitleintuition}%
  \renewtcolorbox{firstread}{enhanced,breakable,colback=white,colframe=UTPBlue, boxrule=0.9pt,arc=2mm,left=2.5mm,right=2.5mm,top=1.8mm,bottom=1.8mm, fonttitle=\bfseries\color{white},colbacktitle=UTPBlue,title=\utptitlefirstread}%
}

% Declaring a box a single book needs. Both constructors take the house
% geometry and let a trailing key list override any of it, so a new box
% is one line and still matches the others:
%
%   \utpnewcallout{odoobox}{white}{UTPBlue}{Odoo 19 Community}
%   \utpnewcallout[boxrule=0.9pt]{name}{back}{frame}{Title}
%   \utpnewquietbox{referencebox}{UTPMuted!58}{UTPMuted}{Reference note}
%
% A callout is a tinted box under a solid title bar, for something the
% reader has to act on. A quiet box is a light frame with coloured
% title text, for an aside they can read past. If a second book ever
% wants the same box, move its line into this section.
\newcommand{\utpnewcallout}[5][]{%
  \newtcolorbox{#2}{enhanced,breakable,colback=#3,colframe=#4, boxrule=0.8pt,arc=2mm,left=2.6mm,right=2.6mm,top=1.8mm,bottom=1.8mm, fonttitle=\bfseries\color{white},colbacktitle=#4,title={#5},#1}%
}
\newcommand{\utpnewquietbox}[5][]{%
  \newtcolorbox{#2}{enhanced,breakable,colback=white,colframe=#3, boxrule=0.5pt,arc=2mm,left=2.5mm,right=2.5mm,top=1.5mm,bottom=1.5mm, fonttitle=\bfseries\color{#4},title={#5},#1}%
}

% --------------------------------------------------------------------
%  9.  SMALL COMMANDS
% --------------------------------------------------------------------
%  \term for a defined term, \cmd for inline code, \chapref and
%  \secref for cross-references, and the utpsteps numbered-step list.
%
%      \term{overfitting}          a term being defined
%      \cmd{git status}            inline code, in any language
%      \chapref{ch:trees}          "Chapter 12"
%      \secref{sec:pruning}        "Section 12.3"
%      \chapterguide{prior}{outcomes}   the box that opens a chapter
%      \begin{utpsteps} ... \end{utpsteps}  a numbered Step 1/Step 2 list
%
%  The second group is the shorthand the chapters use in running text:
%
%      \figref{fig:tree}           "Figure 4.2"
%      \chapsref{ch:a}{ch:b}       "Chapters 3--5"
%      \lib{pandas}                a library name
%      \cxg{x}, \cxo{x}, \cxr{x}   maths coloured green/orange/red
%      \given                      the bar in P(a \given b)
%      \coretag, \latertag         read-now / come-back-to markers
%      \kagglesource{url}{name}    the dataset a lab is built on

\newcommand{\term}[1]{\textbf{\color{UTPNavy}#1}}
\newcommand{\cmd}[1]{\texttt{#1}}
\newcommand{\chapref}[1]{Chapter~\ref{#1}}
\newcommand{\secref}[1]{Section~\ref{#1}}
\newcommand{\chapterguide}[2]{%
  \begin{tcolorbox}[enhanced,breakable,colback=white,colframe=UTPNavy,boxrule=0.8pt,
    arc=2mm,left=2.5mm,right=2.5mm,top=2mm,bottom=2mm,
    title=\utptitleglance,fonttitle=\bfseries\color{white},colbacktitle=UTPNavy]
    \textbf{Before you start:} #1\par
    \textbf{By the end, you can:} #2
  \end{tcolorbox}}
\newenvironment{utpsteps}
  {\begin{enumerate}[label=\textcolor{UTPBlue}{\bfseries Step~\arabic*:},leftmargin=4.5em,itemsep=0.38em]}
  {\end{enumerate}}

\newcommand{\figref}[1]{Figure~\ref{#1}}
\newcommand{\chapsref}[2]{Chapters~\ref{#1}--\ref{#2}}
\newcommand{\lib}[1]{\textbf{\color{UTPBlue}\texttt{#1}}}
\newcommand{\cxg}[1]{{\color{UTPGreen}\bfseries $#1$}}
\newcommand{\cxo}[1]{{\color{UTPOrange}\bfseries $#1$}}
\newcommand{\cxr}[1]{{\color{UTPRed}\bfseries $#1$}}
\newcommand{\given}{\mathrel{|}}
\DeclareMathOperator*{\argmin}{arg\,min}
\DeclareMathOperator*{\argmax}{arg\,max}
\newcommand{\coretag}{\par\smallskip\noindent\colorbox{UTPPaleBlue}{\strut\small\bfseries\color{UTPBlue}CORE NOW}\par\smallskip}
\newcommand{\latertag}{\par\smallskip\noindent\colorbox{UTPPaleGray}{\strut\small\bfseries\color{UTPMuted}RETURN LATER}\par\smallskip}
\newcommand{\kagglesource}[2]{%
  \par\smallskip\noindent\textbf{Kaggle source:} \href{#1}{#2}\par
}

% --------------------------------------------------------------------
%  10.  DIAGRAM NODE STYLES
% --------------------------------------------------------------------
%  The vocabulary the hand-drawn TikZ diagrams are built from: a cell
%  in an array, the one cell under the cursor, a circular graph node,
%  the arrow between two of them. Every book has the whole set:
%
%      utpcell, utphot, utpdone, utpdim, utpwall        array cells
%      utpnode, utpnodehot, utpnodedone, utpnodedim     graph nodes
%      utparrow, utpptr, utpdash, utpedge               connectors
%      utplabel, utptag, utpbrace                       annotation
%      utpstep, utpstage, utppanel, utpbox              boxes
%
%  A cell is 1.0cm by 0.8cm. A book whose diagrams were drawn to a
%  different grid resizes the cells once, after the input:
%
%      \utpdiagramsize{1.05cm}{0.85cm}
%
%  Larger vocabularies for one kind of drawing live in .shared/diagrams/
%  and are loaded by name, after the input:
%
%      \utpusediagrams{network}         devices, links, packet headers,
%                                       sequence diagrams, state machines
%      \utpusediagrams{process}         business-process flows, documents,
%                                       swimlanes, status chips
%      \utpusediagrams{data-objects}    vectors, lists and data frames
%                                       drawn as cells; flowcharts
%      \utpusediagrams{network,process} several at once

\newcommand{\utpdiagramsize}[2]{%
  \tikzset{
    utpcell/.style={draw=UTPBlue,fill=UTPPaleBlue,minimum width=#1,minimum height=#2,
                    font=\small\bfseries\color{UTPNavy},inner sep=1pt},
    utphot/.style={draw=UTPOrange,fill=UTPPaleOrange,minimum width=#1,minimum height=#2,
                   font=\small\bfseries\color{UTPNavy},line width=1pt,inner sep=1pt},
    utpdone/.style={draw=UTPGreen,fill=UTPPaleTeal,minimum width=#1,minimum height=#2,
                    font=\small\bfseries\color{UTPNavy},inner sep=1pt},
    utpdim/.style={draw=UTPMuted!45,fill=UTPPaleGray,minimum width=#1,minimum height=#2,
                   font=\small\color{UTPMuted},inner sep=1pt},
    utpwall/.style={draw=UTPNavy,fill=UTPNavy!82,minimum width=#1,minimum height=#2,
                    inner sep=1pt},
    utpnode/.style={circle,draw=UTPBlue,fill=UTPPaleBlue,minimum size=8mm,
                    font=\small\bfseries\color{UTPNavy},inner sep=0.5pt},
    utpnodehot/.style={circle,draw=UTPOrange,fill=UTPPaleOrange,minimum size=8mm,line width=1pt,
                       font=\small\bfseries\color{UTPNavy},inner sep=0.5pt},
    utpnodedone/.style={circle,draw=UTPGreen,fill=UTPPaleTeal,minimum size=8mm,
                        font=\small\bfseries\color{UTPNavy},inner sep=0.5pt},
    utpnodedim/.style={circle,draw=UTPMuted!45,fill=UTPPaleGray,minimum size=8mm,
                       font=\small\color{UTPMuted},inner sep=0.5pt},
    utparrow/.style={-{Stealth[length=2.2mm]},thick,draw=UTPTeal},
    utpptr/.style={-{Stealth[length=2mm]},semithick,draw=UTPBlue},
    utpdash/.style={-{Stealth[length=2mm]},semithick,draw=UTPRed,dashed},
    utpedge/.style={thick,draw=UTPTeal},
    utplabel/.style={font=\scriptsize\color{UTPMuted},inner sep=1pt},
    utptag/.style={font=\scriptsize\bfseries\color{UTPOrange},inner sep=1pt},
    utpbrace/.style={decorate,decoration={brace,amplitude=4pt,raise=2pt},draw=UTPMuted,semithick},
    utpstep/.style={draw=UTPNavy,fill=white,rounded corners=1.5mm,minimum height=0.75cm,
                    font=\scriptsize\color{UTPNavy},align=center,inner sep=2pt},
    utpstage/.style={draw=UTPNavy,fill=white,rounded corners=1.5mm,minimum height=0.75cm,
                     font=\scriptsize\color{UTPNavy},align=center,inner sep=2pt},
    utppanel/.style={draw=UTPBlue!30,fill=white,rounded corners=1.5mm,inner sep=2.5mm},
    utpbox/.style={draw=UTPBlue,fill=UTPPaleBlue,rounded corners=2mm,align=center,
                   font=\small\color{UTPNavy},inner sep=2.2mm}
  }%
}
\utpdiagramsize{1.0cm}{0.8cm}

% \@for, not \foreach: \foreach runs each pass in a group, which would
% throw away the styles a library declares.
\makeatletter
\newcommand{\utpusediagrams}[1]{%
  \@for\utp@library:=#1\do{\input{../.shared/diagrams/\utp@library}}%
}
\makeatother

% --------------------------------------------------------------------
%  11.  FIGURE FRAMES
% --------------------------------------------------------------------
%  Three ways to put a drawing, chart or screenshot on the page:
%
%      \begin{utpfigure} ... \end{utpfigure}
%          a centred frame with no caption, for a diagram the text
%          talks about directly. An optional title puts a bar on top:
%          \begin{utpfigure}[Where the VLAN tag is added] ...
%
%      \begin{utptable} ... \end{utptable}
%          the same frame in teal, around a tabular or tabularx.
%
%      \begin{utpfig}[short entry]{The caption.}{fig:label}
%        \begin{tikzpicture} ... \end{tikzpicture}
%      \end{utpfig}
%          a frame with a numbered caption that \figref can point to.
%          The optional first argument is the short form for a list of
%          figures. A page never breaks between the drawing and its
%          caption.
%
%  A screenshot or slide lifted from the lecturer's material goes in a
%  titled frame, so the reader can tell it from a drawing made for the
%  notes; the optional width defaults to 0.9 and 0.78 of the line:
%
%      \sourceimage[0.8\linewidth]{Static-routing topology}{static-topology.png}
%      \compactsourceimage{The same, kept to a smaller height}{static-topology.png}
%
%  \fignote{...} adds a short muted line under a drawing, inside its frame.

\newtcolorbox{utpfigure}[1][]{enhanced,width=\linewidth,colback=white,colframe=UTPBlue!28, boxrule=0.7pt,arc=2mm, left=2.2mm,right=2.2mm,top=2.2mm,bottom=1.8mm, before skip=0.75em,after skip=0.85em, before upper={\centering}, halign=center, fonttitle=\bfseries\color{white},colbacktitle=UTPBlue,title={#1}}
\newtcolorbox{utptable}{enhanced,width=\linewidth,colback=white,colframe=UTPTeal!32, boxrule=0.7pt,arc=2mm, left=2.2mm,right=2.2mm,top=1.8mm,bottom=1.8mm, before skip=0.75em,after skip=0.85em, halign=center}
\newcommand{\utpbreakableframes}{%
  \renewtcolorbox{utpfigure}[1][]{enhanced,breakable,colback=white,colframe=UTPBlue!35, boxrule=0.5pt,arc=2mm,left=2.5mm,right=2.5mm,top=2.2mm,bottom=2.2mm, halign=center, fonttitle=\bfseries\color{white},colbacktitle=UTPBlue,title={##1}}%
  \renewtcolorbox{utptable}{enhanced,breakable,colback=white,colframe=UTPBlue!28, boxrule=0.45pt,arc=2mm,left=2.2mm,right=2.2mm,top=1.8mm,bottom=1.8mm}%
}

\captionsetup{
  hypcap=false,
  font={small},
  labelfont={bf,color=UTPNavy},
  textfont={color=UTPInk},
  labelsep=period,
  justification=centering,
  skip=4pt
}

\newenvironment{utpfig}[3][\relax]
  {\def\utpfigshort{#1}\def\utpfigcap{#2}\def\utpfiglab{#3}%
   \par\addvspace{0.55em}%
   \begin{tcolorbox}[enhanced,breakable,colback=white,colframe=UTPBlue!32,
     boxrule=0.5pt,arc=2mm,left=2.5mm,right=2.5mm,top=2.6mm,bottom=2.2mm]
   \centering}
  {\par\nopagebreak\vspace{0.45em}%
   % The caption is boxed: \captionof starts with a write and a zero
   % rule, and the glue after them would otherwise be a place where the
   % page could break, leaving the caption alone at the top of the next.
   % Inside the box the caption is told it sits below the drawing, which
   % it can no longer see, so the gap stays above it.
   \nointerlineskip
   \vbox{\setlength{\parskip}{0pt}\captionsetup{position=below}%
   \ifx\utpfigshort\relax
     \captionof{figure}{\utpfigcap}%
   \else
     \captionof{figure}[\utpfigshort]{\utpfigcap}%
   \fi
   \label{\utpfiglab}}%
   \end{tcolorbox}%
   \par\addvspace{0.55em}}

\newcommand{\sourceimage}[3][0.9\linewidth]{%
  \begin{utpfigure}[Source visual --- #2]%
    \includegraphics[width=#1,height=0.56\textheight,keepaspectratio]{#3}%
  \end{utpfigure}}
\newcommand{\compactsourceimage}[3][0.78\linewidth]{%
  \begin{utpfigure}[Source visual --- #2]%
    \includegraphics[width=#1,height=0.42\textheight,keepaspectratio]{#3}%
  \end{utpfigure}}

\newcommand{\fignote}[1]{\par\smallskip{\scriptsize\color{UTPMuted}#1\par}}

% --------------------------------------------------------------------
%  12.  LAB SHEETS
% --------------------------------------------------------------------
%  Every lab sheet opens with the banner, a full-width page break that
%  also puts the lab in the contents:
%
%      \labsection{3}{Scheduling Policies}{sec:lab03}
%
%  The "of N" comes from \utplabtotal, so a book states its lab count
%  once in its configuration rather than in every lab file.
%
%  The body of the sheet is a labproject box (section 8), and inside it
%  the same toolkit serves every course:
%
%      \labflow{chapter link}{goal}{what to explain}
%          the three-line brief under the lab title.
%      \labsourcecard{file}{what it does}{what to inspect}
%          the heading over a listing of one supplied source file.
%      \begin{labquestion}{3}{Short title} ... \end{labquestion}
%      \begin{labanswer} ... \expectedoutput ... \end{labanswer}
%          a question as the student meets it, then the worked code, the
%          output it should print, and what that output means.
%      \begin{checkanswers} ... \end{checkanswers}
%          answers to a quickcheck, set apart so they can be skipped.
%      labmode, verification, troubleshoot, labdeliverable (section 8)
%          how much is scaffolded, the evidence it worked, what to do if
%          it did not, and what to hand in.

\newcommand{\labsection}[3]{%
  \clearpage
  \phantomsection
  \addcontentsline{toc}{lab}{Lab #1 of \utplabtotal\space--- #2}%
  \markboth{Lab #1 of \utplabtotal\space--- #2}{}%
  \begin{tcolorbox}[enhanced,width=\linewidth,colback=UTPNavy,colframe=UTPNavy,
    boxrule=0pt,arc=2mm,left=5mm,right=5mm,top=5mm,bottom=5mm,
    before skip=0pt,after skip=1.2em]
    {\large\bfseries\color{UTPTeal!35!white} PRACTICAL MILESTONE\par}
    \vspace{0.25em}
    {\Huge\bfseries\color{white} Lab #1 of \utplabtotal\par}
    \vspace{0.20em}
    {\Large\bfseries\color{white} #2\par}
  \end{tcolorbox}
  \label{#3}%
}

\newcommand{\labflow}[3]{%
  \par\smallskip\noindent
  \textbf{Chapter connection.} #1\par
  \textbf{Goal.} #2\par
  \textbf{Explain.} #3\par\medskip}
\newcommand{\labsourcecard}[3]{%
  \par\medskip\noindent\textbf{#1 --- #2}\par
  \textbf{Inspect.} #3\par\smallskip}

\newtcolorbox{labquestion}[2]{enhanced,breakable,colback=UTPPaleOrange!55,colframe=UTPOrange,
  boxrule=0.8pt,arc=2mm,left=2.6mm,right=2.6mm,top=1.8mm,bottom=1.8mm,
  fonttitle=\bfseries\color{white},colbacktitle=UTPOrange,title={Question #1 --- #2}}
\newtcolorbox{labanswer}{enhanced,breakable,colback=white,colframe=UTPTeal,
  boxrule=0.8pt,arc=2mm,left=2.6mm,right=2.6mm,top=1.8mm,bottom=1.8mm,
  fonttitle=\bfseries\color{white},colbacktitle=UTPTeal,title=Worked answer}
\newcommand{\expectedoutput}{\par\smallskip\noindent\textbf{\color{UTPNavy}Reference output.}\nopagebreak\par\nopagebreak}
\utpnewcallout[boxrule=0.65pt,left=2.5mm,right=2.5mm,top=1.6mm,bottom=1.6mm,colbacktitle=UTPMuted]{checkanswers}{UTPPaleGray!72}{UTPMuted!62}{Quick-check answers}

% --------------------------------------------------------------------
%  13.  STAGE DIVIDER
% --------------------------------------------------------------------
%  A full-page divider used by the longer books to open a part:
%  \substage{3A}{Supervised Workflow}{one-line summary}.
%  Harmless to leave unused.

\newcommand{\substage}[3]{%
  \cleardoublepage
  \thispagestyle{empty}
  \phantomsection
  \addcontentsline{toc}{chapter}{Stage #1 --- #2}
  \begin{center}
    \vspace*{0.20\textheight}
    {\Large\bfseries\color{UTPTeal} Stage #1\par}
    \vspace{0.45em}
    {\Huge\bfseries\color{UTPNavy} #2\par}
    \vspace{0.65em}
    {\color{UTPTeal}\rule{0.46\linewidth}{1.5pt}\par}
    \vspace{1.1em}
    \begin{minipage}{0.76\linewidth}
      \centering\large\color{UTPMuted} #3
    \end{minipage}
  \end{center}
  \cleardoublepage
}

% --------------------------------------------------------------------
%  14.  FRONT MATTER
% --------------------------------------------------------------------
%  The opening every book shares: the cover, the contents in roman
%  numerals, the course map, then chapter one on an arabic page 1.
%
%      \begin{document}
%      \utpfrontmatter
%
%  What it prints comes from the configuration in section 1 --- the
%  cover text from \utpcoursetitle and \utpcourseoutline, the overview
%  chapter from \utpcoursemap --- and the contents page is set as
%  \utpcontentspage (section 5) describes. A book with no overview
%  chapter sets \def\utpcoursemap{} and the step is skipped.
%
%  Page anchors are off across the cover so hyperref does not warn
%  about the unnumbered pages, and back on before the contents links
%  are used.
%
%  The number of front-matter pages is recorded in the .aux file as
%  \utpfrontpages, which is how .shared/build.py turns the page numbers
%  in the contents into physical pages when it splits the book.

\newcommand{\utpfrontmatter}{%
  \hypersetup{pageanchor=false}%
  % ---------- Cover ----------
%  The title page, printed by \utpfrontmatter. Everything that changes
%  between books is a configuration macro from section 1 of utpnotes:
%  \utpcoursetitle, \utpcourseoutline, \utpoutlineheading, \utpauthor
%  and \utpauthorrole. The colours are the shared palette.

\begin{titlepage}
\thispagestyle{empty}
\begin{tikzpicture}[remember picture,overlay]
  \fill[UTPGold]
    (current page.north west) rectangle
    ([yshift=-0.12cm]current page.north east);
  \fill[UTPGold]
    (current page.south west) rectangle
    ([yshift=0.12cm]current page.south east);
\end{tikzpicture}

\centering
\includegraphics[
  height=2.85cm,
  clip
]{logo-UTP.png}\par

\vspace{0.50cm}
{\fontsize{29}{33}\selectfont\bfseries\color{UTPNavy}\utpcoursetitle\par}
\vspace{0.24cm}
{\Large\color{UTPMuted} \utpdocumentsubtitle\par}
\vspace{0.42cm}
{\color{UTPGold}\rule{0.26\linewidth}{1.3pt}\par}
\vspace{0.72cm}
\begin{tikzpicture}[
  core/.style={circle,draw=UTPGold,fill=UTPNavy,minimum size=1.95cm,
    align=center,font=\small\bfseries\color{white},line width=1pt},
  stage/.style={rounded corners=1.5mm,draw=UTPBlue!65,fill=UTPPaleBlue!42,
    text width=3.15cm,minimum height=1.42cm,align=center,
    font=\small\bfseries\color{UTPNavy},inner xsep=3mm,
    inner ysep=3.5mm,line width=0.65pt},
  cycle/.style={-{Stealth[length=2.1mm]},semithick,draw=UTPMuted,
    rounded corners=4mm,shorten <=7pt,shorten >=7pt}
]
\node[core] (guide) at (0,-2.55) {TA\\WORKFLOW};
\node[stage] (prepare) at (0,0) {SESSION\\PLANNING};
\node[stage] (explain) at (4.75,-2.55) {INSTRUCTIONAL\\SUPPORT};
\node[stage] (facilitate) at (0,-5.10) {STUDENT\\ENGAGEMENT};
\node[stage] (assess) at (-4.75,-2.55) {LEARNING\\ASSESSMENT};

\draw[cycle] (prepare.east) -- (explain.north);
\draw[cycle] (explain.south) -- (facilitate.east);
\draw[cycle] (facilitate.west) -- (assess.south);
\draw[cycle] (assess.north) -- (prepare.west);
\end{tikzpicture}

\vspace{0.72cm}
\begin{tcolorbox}[
  enhanced,width=0.82\linewidth,colback=white,colframe=UTPBlue!18,
  arc=1.5mm,boxrule=0.45pt,left=6mm,right=6mm,top=3.2mm,bottom=3.2mm,
  borderline west={1.6pt}{0pt}{UTPGold}]
\centering\normalsize\color{UTPInk}
{\large\bfseries\color{UTPBlue}\utpoutlineheading\par}
\vspace{0.3em}
\utpcourseoutline
\end{tcolorbox}

\vfill
\begin{minipage}{0.82\textwidth}
\begin{tabular}{@{}>{\RaggedRight\arraybackslash}m{0.42\linewidth}@{\hspace{0.10\linewidth}}>{\RaggedLeft\arraybackslash}m{0.48\linewidth}@{}}
{\normalsize\bfseries\color{UTPBlue}PREPARED BY\par\vspace{0.25em}
 \Large\bfseries\color{UTPNavy}\utpauthor\par\vspace{0.12em}
 \normalfont\small\color{UTPMuted}\utpauthorrole}
&
\includegraphics[height=2.10cm,keepaspectratio]{logo-DC.png}
\end{tabular}
\end{minipage}
\vspace{0.5cm}
\end{titlepage}%
  \cleardoublepage
  \pagenumbering{roman}%
  \setcounter{tocdepth}{0}%
  \utptableofcontents
  \newpage
  \hypersetup{pageanchor=true}%
  \ifx\utpcoursemap\empty\else\input{\utpcoursemap}\fi
  \cleardoublepage
  \ifdefined\ReadonlyShipoutCounter
    \immediate\write\csname @auxout\endcsname
      {\string\gdef\string\utpfrontpages{\the\ReadonlyShipoutCounter}}%
  \fi
  \pagenumbering{arabic}%
}

% --------------------------------------------------------------------
%  15.  LISTING STYLES
% --------------------------------------------------------------------
%  Two base looks, and one named style per language built on them, so
%  the colours and frame are the same in every book:
%
%      utpcode          the base for source code
%      utpterminal      the base for commands typed and output printed
%      utppython        Python
%      utpr             R
%      utpc             C
%      utpcpp           C++
%      utpbash          shell scripts and commands
%      utpcnumbered     C with line numbers, for a complete source file
%      utpbashnumbered  shell with line numbers, likewise
%
%      \begin{lstlisting}[style=utppython] ... \end{lstlisting}
%      \lstinputlisting[style=utpcnumbered]{all_experiments/Lab04/fcfs.c}
%
%  A book that sets most of its listings in one style makes it the
%  default once, after the input:  \lstset{style=utpcpp}

\lstdefinestyle{utpcode}{
  basicstyle=\ttfamily\scriptsize,
  keywordstyle=\bfseries\color{UTPBlue},
  commentstyle=\itshape\color{UTPMuted},
  stringstyle=\color{UTPGreen!70!black},
  backgroundcolor=\color{UTPPaleGray!72},
  frame=single,
  rulecolor=\color{UTPBlue!28},
  framesep=5pt,
  breaklines=true,
  breakatwhitespace=false,
  showstringspaces=false,
  keepspaces=true,
  columns=fullflexible,
  upquote=true,
  tabsize=4,
  aboveskip=0.7em,
  belowskip=0.8em
}
\lstdefinestyle{utpterminal}{
  basicstyle=\ttfamily\scriptsize,
  backgroundcolor=\color{UTPPaleGray!72},
  frame=single,
  rulecolor=\color{UTPMuted!35},
  framesep=5pt,
  breaklines=true,
  showstringspaces=false,
  keepspaces=true,
  columns=fullflexible,
  upquote=true,
  aboveskip=0.6em,
  belowskip=0.7em
}

\lstdefinestyle{utppython}{style=utpcode,language=Python}
\lstdefinestyle{utpr}{style=utpcode,language=R,tabsize=2}
\lstdefinestyle{utpc}{style=utpcode,language=C}
\lstdefinestyle{utpcpp}{style=utpcode,language=C++}
\lstdefinestyle{utpbash}{style=utpcode,language=bash,
  rulecolor=\color{UTPMuted!35},aboveskip=0.6em,belowskip=0.7em}
\lstdefinestyle{utpcnumbered}{style=utpc,numbers=left,numberstyle=\tiny\color{UTPMuted},
  numbersep=7pt,xleftmargin=1.7em,framexleftmargin=1.4em}
\lstdefinestyle{utpbashnumbered}{style=utpbash,numbers=left,numberstyle=\tiny\color{UTPMuted},
  numbersep=7pt,xleftmargin=1.7em,framexleftmargin=1.4em}

% Pseudocode. A book that sets algorithms asks for them once, after the
% input, and gets algorithm2e loaded and coloured to match:
%
%     \utpusealgorithms
%
% \algostyle then sets one algorithm environment a size smaller than the
% text around it.
\newcommand{\utpalgorithmstyle}{%
  \SetAlCapNameFnt{\bfseries\color{UTPNavy}}%
  \SetAlCapFnt{\bfseries\color{UTPNavy}}%
  \SetAlgoCaptionSeparator{.}%
  \SetKwInput{KwIn}{Input}%
  \SetKwInput{KwOut}{Output}%
  \SetNlSty{color}{UTPMuted}{}%
  \SetCommentSty{\itshape\color{UTPMuted}}%
  \SetKwSty{bfseries\color{UTPBlue}}%
  \SetArgSty{}%
  \providecommand{\algostyle}{\small\SetAlgoNlRelativeSize{-1}\DontPrintSemicolon\SetAlgoVlined}%
}
\newcommand{\utpusealgorithms}{%
  \usepackage[linesnumbered,ruled,vlined]{algorithm2e}%
  \utpalgorithmstyle
}
